\section{The shape theorem and the global theorem}\label{sec:pf:top-shape}

\providecommand{\topshapeCl}{\mathrm{Closed}}
\providecommand{\topshapeLarge}{\mathrm{Large}}
\providecommand{\topshapeDes}{\mathrm{Des}}
\providecommand{\topshapeD}{\mathcal{D}_3}
\providecommand{\topshapeF}{\mathcal{F}}
\providecommand{\topshapeWin}{\mathrm{Win}}
\providecommand{\topshapeEta}{\eta\Sigma}
\providecommand{\topshapeVV}{\mathrm{vv}}
\providecommand{\topshapewt}{\mathrm{wt}}

This section proves the shape theorem (\cref{thm:shape}): Green's parametrisation of the cuspidal characters
of $\GL_d(\F_g)$, together with the data and hypotheses of an element of the hypothesis package, implies the
proposition $\sfname{TopG}$ of \cref{def:topg}. It also proves the global theorem (\cref{thm:global-chain}),
which turns the family-level ``links'' into chain-tied records, and the results of \cref{sec:families:links}
which establish the individual links: the label weights (\cref{prop:families:admLq}), the irreducibility link
(\cref{thm:linkIRR}), the division link (\cref{prop:families:div}), the approximate orthogonality link
(\cref{cor:families:linkS}), the mass and slope statements (\cref{prop:families:famgl}) and the existence of
designed families (\cref{prop:families:design}).

\paragraph{The idea.}
The statement $\sfname{TopG}$ asks, for every exponent $\lambda>1$ and every requirement $\mathcal G$ with
$\sfname{HostExW}(\lambda,\mathcal G)$, for families in the class $\mathcal G.\mathrm{cls}(\lambda)$ which
are chain-tied at all large window scales. The proof has three layers.
\begin{enumerate}[label=(\roman*)]
\item \emph{Closedness and design.} Almost every hypothesis and every intermediate statement is a
\emph{closed} property of families: it holds for all members of a class whose parameter $\eps$ is small and
whose constants are large with respect to a staged threshold. The hypothesis $h_{\mathrm{Ex}}$ of the package
produces, for every small $\eps$ and every threshold, a \emph{designed} family whose constants exceed the
threshold. Closedness and design are dual, and the assembly theorem (\cref{pf:top-shape:assembly}) combines them.
\item \emph{Links.} Each link is an eventual inequality or structural condition on the members of a family.
Each link is proved closed on the class $\topshapeD(\mathcal A)$ of designed members of the family class
$\mathcal A$ of the package, from fields of the package (\cref{sec:pf:top-shape:irr,sec:pf:top-shape:links,sec:pf:top-shape:div,sec:pf:top-shape:gl}).
Green's theorem is used only once, for the irreducibility link.
\item \emph{The global theorem.} From the links and the eleven mass, slope and integrality properties of
\cref{prop:families:famgl}, the global theorem builds, at every large window scale, a record in the sense of
\cref{def:record} from the members of the family, and verifies the eleven systems of size conditions
\ref{S:setting}--\ref{S:fparams} and the identification of \cref{def:chaintied}
(\cref{sec:pf:top-shape:global}).
\end{enumerate}

\paragraph{What is proved and what is cited.}
The formal statement of the shape theorem is an implication from Green's theorem and about sixty hypotheses
on an arbitrary family class; its arithmetic content lies in the hypotheses, which are supplied by the
non-emptiness of the package (\cref{thm:hyps-nonempty}). We prove the theorem from the hypotheses completely at
the level of the assembly: the closedness calculus, the assembly theorem, the derivation of $\sfname{DIVinN}$
from the Sha hypothesis, the derivation of $\sfname{LinkMu}$, the label weights from Green's theorem, the
counting argument behind $\sfname{LinkIRR}$, a Mertens-type estimate for the tops, the wiring of the division
link, the global link, the hook-length bound for Fock norms, the construction of the record in the global
theorem, and the existence of designed families from Minkowski's bound. We also prove, at the level of their main
steps and with the estimates written out, the frame-number bound (with the slope lemma on the division side and the
gauge-form bound), the Arakelov slope inequalities, the sum-over-primes service theorem and the orbit-entry bounds
behind the admissibility of twisted periods (including the eight local orbit bounds, the linear-programming estimate for
the radii, the Laplace-operator count at the ramified places, the Hecke lattice count and the chain estimates for the
formal-group statements), the transport of Theorem~S and the thin-scale prime sums from $\sfname{LitThin}$; their
deepest remaining ingredients (the orbit letter bounds and the generic kernel bounds of the local letter analysis,
Dwork's weight estimates, some algebraic identities for the theta series of formal-group systems, the archimedean class
map, the exact slopes of column lattices, and the analytic estimates for Dirichlet $L$-functions) are cited by their
formal names. The local lattice link $\sfname{LinkG}$
is proved in \cref{sec:pf:linkG} and cited as \lref{Formalism.L2.G.linkG_concrete5}. The family statements
entering the global theorem are stated precisely where they are used and cited by their formal names.

\paragraph{Organisation.}
\Cref{sec:pf:top-shape:framework} recalls families, closed properties, the decorated classes and designed
families. \Cref{sec:pf:top-shape:statement} defines the links and the shape, and states the main theorem.
\Cref{sec:pf:top-shape:assemblysec} proves the assembly theorem and deduces the shape theorem from the link
theorems. \Cref{sec:pf:top-shape:irr} treats Green's theorem and $\sfname{LinkIRR}$,
\cref{sec:pf:top-shape:links} the links $\sfname{LinkG}$, $\sfname{LinkEta}$, $\sfname{SmallLevelRankEq}$ and
$\sfname{LinkS}$, \cref{sec:pf:top-shape:div} the division link, \cref{sec:pf:top-shape:gl} the mass and slope
statements and the global link, \cref{sec:pf:top-shape:global} the global theorem and
\cref{sec:pf:top-shape:design} the designed families. \Cref{sec:pf:top-shape:remarks} collects remarks on the
formal statements.

% ===========================================================================
\subsection{Families, closed properties and designed families}\label{sec:pf:top-shape:framework}

We use the families of \cref{sec:top} and \cref{def:families:closed}. A family $\topshapeF$ with parameter
$\eps$ has a threshold $x_1>1$, index sets $I_x$, a set $\topshapeWin(\topshapeF)$ of window scales (unbounded
above), members $(D_{x,i},L_{x,i})$ consisting of an arithmetic datum and a size record whose scale is
$(x,\eps,J)$ for $x>x_1$, extra data $\mathrm{ex}(x,i)$, and lists of family constants. For a member we write
$n,m,k,f,r,t$ for the ranks ($f=2n$, $r=k+m-1$), $\Lambda=x^{\eps}$, $b=e^{\Lambda}$, $a=xb$,
$\delta_1=x^{-\eps/8}$, $X'=\exp(x^{\eps/3})$, $X''=X''(n)$, $B_0=fmb\Lambda$, and $N_c$ for the scale
quantities of $L_{x,i}$, and $\eps_z=\eps/J$. A property $Q(x,i)$ of members holds \emph{eventually in
$\topshapeF$} if there is $x_0$ such that $Q(x,i)$ holds for all $x\ge x_0$ and all $i\in I_x$. A
\emph{family class} is a predicate on families (of any parameter); we write $\topshapeF\in\mathcal C$.

\begin{definition}[Staged thresholds and largeness]\label{pf:top-shape:large}
A \emph{staged threshold} is a quadruple $K=(K_{\rm pre},K_t,K_{\rm mid},K_{\rm last})$ of a real number and
three functions from finite lists of reals to $\R$. A family $\topshapeF$ carries three lists of constants: an
outer list $c_{\rm out}$ of sixteen constants, the list $c_{\rm pre}=[C_P]$, a middle list
\[
c_{\rm mid}(C_H)=[C_H,\ C_\Gamma,\ C_{\rm sz},\ C_{\rm sz},\ C_{\rm sz},\ C_{\rm sz},\ C_{\rm sz},\ c_{\rm err},\
G_K.\mathrm{La},\ C_F]
\]
depending on a real number $C_H$ (the \emph{level}), and a last list $c_{\rm last}$ of thirty-two constants.
The family is \emph{$K$-large at level $C_H$}, written $\topshapeLarge(\topshapeF,K,C_H)$, if every entry of
$c_{\rm pre}$ is $\ge K_{\rm pre}$, every entry of $c_{\rm mid}(C_H)$ is $\ge K_{\rm mid}(c_{\rm out})$, and
every entry of $c_{\rm last}$ is $\ge K_{\rm last}(c_{\rm out}\,c_{\rm mid}(C_H))$ (concatenation of lists).
The supremum $K\sqcup K'$ of two thresholds is the componentwise maximum (pointwise for the functions); then
$\topshapeLarge(\topshapeF,K\sqcup K',C_H)$ holds if and only if both $\topshapeLarge(\topshapeF,K,C_H)$ and
$\topshapeLarge(\topshapeF,K',C_H)$ hold.
\formal{Formalism.MainNotions.FamP.Large}\formallabel{Formalism.MainNotions.FamP.large_sup}
\end{definition}

\begin{definition}[Closed properties]\label{pf:top-shape:closed}
Let $\mathcal C$ be a family class and $P(\topshapeF,C_H)$ a property of families and real numbers.
$P$ is \emph{closed on $\mathcal C$}, written $\topshapeCl(\mathcal C,P)$, if there is $\eps_1>0$ such that
for every $\eps\in(0,\eps_1]$ there is a staged threshold $K$ with $P(\topshapeF,C_H)$ for every family
$\topshapeF\in\mathcal C$ with parameter $\eps$ and every $C_H$ with $\topshapeLarge(\topshapeF,K,C_H)$.
\formal{Formalism.MainNotions.Closed}
\end{definition}

\begin{lemma}[Calculus of closed properties]\label{pf:top-shape:calculus}
Let $\mathcal C,\mathcal C'$ be family classes and $P,Q$ properties.
\begin{enumerate}[label=(\alph*)]
\item If $P(\topshapeF,C_H)$ holds for all $\topshapeF\in\mathcal C$ and all $C_H$, then
$\topshapeCl(\mathcal C,P)$.
\item $\topshapeCl(\mathcal C,P)$ and $\topshapeCl(\mathcal C,Q)$ imply $\topshapeCl(\mathcal C,P\wedge Q)$.
\item If $\topshapeCl(\mathcal C,P)$, and $P(\topshapeF,C_H)$ implies $Q(\topshapeF,C_H)$ for all families
with positive parameter, then $\topshapeCl(\mathcal C,Q)$.
\item If $\mathcal C'\subseteq\mathcal C$, then $\topshapeCl(\mathcal C,P)$ implies $\topshapeCl(\mathcal C',P)$.
\item $\topshapeCl\bigl(\mathcal C,\ (\topshapeF,C_H)\mapsto(\topshapeF\in\mathcal C\Rightarrow
P(\topshapeF,C_H))\bigr)$ implies $\topshapeCl(\mathcal C,P)$.
\item For every class $\mathcal C$, $\topshapeCl(\mathcal C,\ C_H>0)$, and
$\topshapeCl(\mathcal C,\ c(c_{\rm out})\le C_{\rm sz})$ for every function $c$ of lists of reals.
\end{enumerate}
\formallabel{Formalism.MainNotions.Pk.Closed.of_class}
\formallabel{Formalism.MainNotions.Pk.Closed.sub}
\formallabel{Formalism.L2.Top3.closed_of_mem}
\formallabel{Formalism.MainNotions.Pk.closed_CH}
\formallabel{Formalism.MainNotions.DIV.closed_Csz_ge_fn}
\end{lemma}

\begin{proof}
(a) Take $\eps_1=1$ and the threshold all of whose components are $0$. (b) and (c) are
\lref{Formalism.MainNotions.Closed.and} and \lref{Formalism.MainNotions.Closed.imp}: for (b) take the smaller
of the two $\eps_1$ and, for given $\eps$, the supremum of the two thresholds, which works by the last
assertion of \cref{pf:top-shape:large}; (c) uses the same $\eps_1$ and thresholds. (d) and (e) use the same
$\eps_1$ and thresholds. (f) Take $\eps_1=1$ and $K_{\rm mid}$ the constant function $1$ (resp.\ the function
$c$), the other components being $0$. Since $C_H$ and $C_{\rm sz}$ are entries of $c_{\rm mid}(C_H)$, largeness
gives $C_H\ge1>0$ (resp.\ $C_{\rm sz}\ge c(c_{\rm out})$).
\end{proof}

\begin{definition}[Size conditions and choices]\label{pf:top-shape:sizes}
For a family $\topshapeF$ and members $(D,L)=(D_{x,i},L_{x,i})$:
\begin{itemize}
\item $\sfname{TuSize}(\topshapeF)$: eventually the size condition $\sfname{TuGe}$ of $L$ holds and
$\log T_u\le C_{Tu}\,m\,\Lambda$ for every word $u$ of $D$, where $T_u$ is the level prime of $u$ and $C_{Tu}$ a
family constant;
\item $\sfname{DGamSize}(\topshapeF)$: there is $C$ such that eventually $\log D_\Gamma(u)\le C\,fN(m)\log(fN(m))$
for every word $u$, where $D_\Gamma(u)$ is the integer attached to the member of $u$ and $N(j)=k+j(t+1)$;
\item $\sfname{DGamDvdF}(\topshapeF)$: eventually $D_\Gamma(u)$ divides the Minkowski number
$\mathrm{M}(2fN(m))$ for every $u$; $\sfname{DGamSmooth}(\topshapeF)$: eventually every prime divisor $p$ of
$D_\Gamma(u)$ satisfies $p\le2fN(m)+1$;
\item $\sfname{ChosenE}(\topshapeF,C_0,C_H)$: there are $C_5,c_6,C_b>0$ such that the seven \emph{choice
conditions} $\sfname{ThickSetEventually}$, $\sfname{ChosenTT}$, $\sfname{ChosenBad}(C_b)$,
$\sfname{ChosenQD}(C_5)$, $\sfname{ChosenThin}(c_6)$, $\sfname{ChosenProviso}(C_0)$ and
$\sfname{ChosenCode}(C_0,C_H)$ hold; each is an eventual condition on the members. For instance
$\sfname{ChosenTT}$ says that eventually $\delta_1n/2\le|TT|\le2\delta_1n$, where $TT$ is the set of tops of
the host.
\end{itemize}
The remaining conditions named below ($\sfname{TwGood}$, $\sfname{TameDistinctF}$, $\sfname{BoxesF}$,
$\sfname{QDRowsF}$, and the inequality $C_{YL}\le9$ between family constants) are conditions on families of the
development that enter only as bookkeeping.
\formal{Formalism.MainNotions.TuSize}\formal{Formalism.MainNotions.DGamSize}\formal{Formalism.MainNotions.ChosenE}
\end{definition}

\begin{definition}[Decorated classes and designed families]\label{pf:top-shape:decor}
For a family class $\mathcal C$ and $C_0\in\R$ put
\begin{gather*}
\mathrm{WithChN}(\mathcal C)=\mathcal C\wedge\sfname{DGamDvdF}\wedge\sfname{TameDistinctF}\wedge
\sfname{BoxesF}\wedge\sfname{QDRowsF}\wedge(C_{YL}\le9),\\
\mathrm{WithSm}(\mathcal C)=\mathcal C\wedge\sfname{DGamSmooth},\qquad
\mathrm{WithTw}(\mathcal C)=\mathcal C\wedge\sfname{TwGood},\\
\topshapeDes(\mathcal C,C_0)=\mathrm{WithTw}(\mathrm{WithSm}(\mathcal C))\wedge
\bigl(\exists C_H:\sfname{ChosenE}(C_0,C_H)\bigr)\wedge\sfname{TuSize}\wedge\sfname{DGamSize},
\end{gather*}
and $\topshapeD(\mathcal A)=\topshapeDes(\mathrm{WithChN}(\mathcal A),3)$ (denoted $D3(3,\mathcal A)$ in
\cref{sec:families}). $\sfname{DesignExists}(\mathcal C,C_0)$ means: there is $\eps_0>0$ such that for every
$\eps\in(0,\eps_0]$ and every staged threshold $K$ there are a family $\topshapeF\in\mathcal C$ with parameter
$\eps$ and $C_H\ge0$ with $\topshapeLarge(\topshapeF,K,C_H)$, $\sfname{ChosenE}(\topshapeF,C_0,C_H)$,
$\sfname{TuSize}(\topshapeF)$ and $\sfname{DGamSize}(\topshapeF)$. $\sfname{DesignChoices}(\mathcal C,C_0,c_T)$
is the analogous statement in which the family has, instead of the last three properties, the raw choices
$\sfname{ThickSetEventually}$, $\sfname{ChosenTT}$, $\sfname{ChosenBad}(C_b)$, $\sfname{ChosenQD}(C_5)$ (with
$C_5,C_b>0$), $\sfname{ChosenCode}(C_0,C_H)$, $\sfname{BoxesF}$, $\sfname{QDRowsF}$, $\sfname{DGamDvdF}$,
$\sfname{TameDistinctF}$, $\sfname{CodeUpF}(C_0,C_H)$, and $C_F=C_H$, $C_{YL}\le9$ and
$c_T(\eps,C_{YL})\le C_{Tu}$.
\formal{Formalism.MainNotions.Des}\formallabel{Formalism.MainNotions.DesignExists}\formal{Formalism.MainNotions.DesignChoices}
\end{definition}

Closedness says nothing about classes without large members, and $\sfname{DesignExists}$ says nothing about
properties of the families it produces beyond the listed ones; the assembly theorem combines the two.

% ===========================================================================
\subsection{The links, the shape and the main theorem}\label{sec:pf:top-shape:statement}

The links of \cref{def:links} are eventual conditions on the members $(D,L)=(D_{x,i},L_{x,i})$ of a family. We
state them in the form used below; $j$ ranges over the set $\mathrm{lifts}(D)=\{1,\dots,m\}$ of lift indices,
$u,u'$ over the words of $D$, and $w$ over the places of the coefficient field $\mathit{LL}$ in the support of
$L$.

\begin{definition}[Links]\label{pf:top-shape:links}
Put
\[
\topshapeEta(L,j)=\sum_{w:\ \mu_w\ne\mu_0}\mathrm{lw}(w)\,e_w\,\eta_j(w),\qquad
\topshapeVV(D,\ell)=\log\ell\sum_{e=1}^{t-1}\mathrm{sh}_\ell(D,e)\quad(\ell\text{ prime}),
\]
and $\topshapeVV(D,\ell)=0$ for $\ell$ not prime; here the sum runs over the twisted places $w$ of the support,
$e_w$ is the ramification index, $\mathrm{lw}(w)$ a weight, $\eta_j(w)\ge0$ the local indices of $D$, and
$\mathrm{sh}_\ell(D,e)\in\Q_{\ge0}$ the Sha lengths recorded in $D$.
\begin{itemize}
\item $\sfname{LinkN}$: there is $C$ such that eventually, if all Galois-cohomological lengths
$\mathrm{shaLen}(\mathrm{ld}(p),e)$ ($p$ prime, $0<e<t$) of the local data of $D$ are finite
(\emph{$\sfname{ShaFinOdd}(D)$}) and $\ell\mapsto\topshapeVV(D,\ell)$ is summable, then for all $j$
\[
\topshapeEta(L,j)\le d_{Z_0}(j)\Bigl(C\bigl(fmb\Lambda^2+(\log x)^2N_c\bigr)+\sum_\ell\topshapeVV(D,\ell)\Bigr).
\]
\item $\sfname{DIVinN}$: there is $C$ such that eventually $\topshapeEta(L,j)\le D_0\cdot C\,x^{\eps}(B_0+N_c)$
for all $j$.
\item $\sfname{LinkEta}$: eventually, for $1\le j\le m$ and every place $w$ with $\mu_w\ne\mu_0$, the local
modules $\mathit{Mt}_i$, $i\in\mathrm{locs}(j,w)$, of $D$ have total rank $\dim\mathit{Mt}(\mu_w,j)$.
\item $\sfname{SmallLevelRankEq}$: eventually $d_{Z_0}(j)=D_0$ for all $j$.
\item $\sfname{LinkIRR}$: eventually $0\le\delta_p(D)\le\exp(-x^{\eps_z/7})$, where $\delta_p$ is defined by
$D_{\rm rows}=(1-\delta_p)D_{\rm all}$.
\item $\sfname{LinkDIV}(C_{\rm abs},C_H)$: there is $C$ such that eventually the following nine conditions hold:
(1) the functionals $B_u$ map the lattices $N_u$ into the lattice $\mathit{CC}$ (``one map into
$\mathit{CC}$''); (2) $s_{\rm Frame}(N,h)+\mu_{\max}(\mathit{RR})\le x^{C_H}b$; (3) $\sfname{ZeroAbove}$: for
every word $u$ and every place $w$ whose residue characteristic is larger than $X''$ and different from $T_u$,
$g_0$, $g_1$ and $d$, and does not divide $D_\Gamma(u)$, the gain $\mathrm{gain}(j,w)$ and the term $E(u,j,w)$ of
$L$ vanish; (4) the Bessel identity $\sfname{BesselId}(D)$: the functionals $B_u$ are rational (they are the
embeddings of the functionals $B_k(u)$ of the division data), and there is an isometry $V$ of the space
$\mathit{EE}$ such that for every word $u$ there are a scalar $a_u\ne0$ and an isometry $U_u$ of the space of $u$
with $B_u(\varphi,F)=a_u\,\mathrm{bes}_u(U_u\varphi,VF)$ for a fixed pairing $\mathrm{bes}_u$; (5) $\mathit{RR}$
is finitely generated and spans its ambient space; (6) $\mathrm{ht}(B_k(u))\le h(u)$ whenever $B_k(u)\ne0$; (7) for
every $u$,
\[
\mu_{\rm dual}(\mathit{RR})+h(u)-F_o\le\mathrm{MAIN}(u)+C|TT|aL_a+C\,m(B_0+N_c)x^{C_{\rm abs}\eps};
\]
(8) $\bigl|\sum_w\mathrm{sz}(w,E(u,j,w))-\bigl(\sum_{t\in TT}a_T(t)\bigr)\log(g_0g_1)\bigr|\le CB_0$ for all
$u,j$; (9) the quantity $h(u)+\bigl(\sum_ta_T(t)\bigr)\log T_u$ varies with $u$ by at most $C(|TT|aL_a+mB_0)$.
\item $\sfname{LinkDIVt}(C_{\rm abs},C_H)$: $\sfname{LinkEta}\wedge\sfname{DIVinN}\Rightarrow
\sfname{LinkDIV}(C_{\rm abs},C_H)$.
\item $\sfname{LinkG}$: there are $C$ and a function $\delta(x)\to0$ such that eventually the local lattice
conditions $\sfname{LinkG\_at}$ of \cref{thm:linkG} hold with $b^{\delta(x)}$ in place of $b$.
\item $\sfname{LinkS}(C_0)$: eventually there are $\eta_0,\eta_1\in\R$ and $e_0\ge0$ with
$\eta_0,\eta_1,\sqrt{e_0}\le\exp(-x^{\eps/8})$ such that, with $\Gamma=\exp(C_\Gamma m^5f\Lambda)$: if
$C_0^{-1}\le d_u\Gamma^2/D_*\le C_0$ for all words $u$, then $\mathrm{hs}(u,u')\le(1+\eta_0)d_ud_{u'}/D_*$ for
$u\ne u'$, $\|K_{uu}-\mathrm{id}\|_{\rm HS}^2\le e_0d_u$ for all $u$, and $D_{\rm all}\le(1+\eta_1)D_*$.
\item $\sfname{LinkMu}$: there is $C$ such that eventually
$\bigl|\mu(N_u)-\mu(N^H_u)-\sum_w\mathrm{sz}(w,\sigma(u,w))\bigr|\le C\,X''$ for all $u$.
\item $\sfname{NuFullF}$: eventually every lattice $N_u$ is finitely generated and spans its ambient space;
$\sfname{MuRank}$: eventually $\rank N_u>0$ for all $u$.
\end{itemize}
Finally $\sfname{GlobalTheoremInputs}(C_{\rm abs},C_H,C_0)$ is the conjunction
$\sfname{LinkG}\wedge\sfname{LinkS}(C_0)\wedge\sfname{LinkIRR}\wedge\sfname{LinkDIV}(C_{\rm abs},C_H)\wedge
\sfname{LinkMu}$, and $\sfname{LinkGLOBAL}(C_{\rm abs},C_H,C_0,e_D,e_L)$ is the statement: if
$\sfname{ChosenE}(C_0,C_H)$ and $\sfname{GlobalTheoremInputs}(C_{\rm abs},C_H,C_0)$, then there are
$K,c_{HH},e_K,C_L,x_0$ such that $\topshapeF$ is chain-tied with exponents $(e_D,e_L)$ and constants
$(K,c_{HH},e_K,C_L)$ (\cref{def:chaintied}) at every window scale $x\ge x_0$.
\formal{Formalism.MainNotions.LinkN}\formal{Formalism.MainNotions.DIVinN}\formal{Formalism.MainNotions.LinkEta}
\formal{Formalism.MainNotions.SmallLevelRankEq}\formal{Formalism.MainNotions.LinkIRR}
\formal{Formalism.MainNotions.LinkDIV}\formal{Formalism.MainNotions.LinkS}\formal{Formalism.MainNotions.LinkMu}
\formallabel{Formalism.MainNotions.GlobalTheoremInputs}\formallabel{Formalism.MainNotions.LinkGLOBAL}
\end{definition}

The Sha hypothesis of \cref{def:topg} is: $\sfname{HypV}(\topshapeF)$ holds if eventually
$\sfname{VVle}(E_{x,i},t,\sqrt b)$ (\cref{def:vwin}), where $E_{x,i}$ is the CM field of the host of $D_{x,i}$.
A family class $\mathcal E$ is \emph{connected} if there are $e_D,e_L$ and $\eps_1>0$ such that for every
$\eps\in(0,\eps_1)$ some family $\topshapeF\in\mathcal E$ with parameter $\eps$ satisfies: there are
$K,c_{HH},e_K,C_L,x_0$ such that $\topshapeF$ is chain-tied with exponents $(e_D,e_L)$ and these constants at
every window scale $x\ge x_0$. $\sfname{TopG}$ (\cref{def:topg}) says that $\mathcal G.\mathrm{cls}(\lambda)$
is connected whenever $\lambda>1$, $\sfname{HostExW}(\lambda,\mathcal G)$ holds and every family of
$\mathcal G.\mathrm{cls}(\lambda)$ satisfies $\sfname{HypV}$. Here a family belongs to
$\mathcal G.\mathrm{cls}(\lambda)$ if for every $y$ there is $Y\ge y$ with $[Y,Y^\lambda]\subseteq
\topshapeWin(\topshapeF)$, and $\mathcal G$ holds for the data of all members with $x>x_1$.

\begin{definition}[The shape]\label{pf:top-shape:shape}
Let $V,N_R$ be propositions, $N$ a predicate on family classes and $G$ a proposition. The \emph{shape}
$\sfname{TopShapeAlt}(V,N_R,N,G)$ is the statement: for every family class $\mathcal A$, every function
$c_T\colon\R^2\to\R$ and all of the following data and hypotheses, $G$ holds. We write $\topshapeD=\topshapeD(\mathcal A)$;
``selected members'' are the members $(x,i)$ of families $\topshapeF\in\mathcal A$ with $x>x_1$ and
$x\ge x_{\rm sel}(\topshapeF)$.
\begin{enumerate}[label=(\alph*)]
\item \emph{Design}: $h_{\rm Ex}$: $\sfname{MinkDiv}\Rightarrow\sfname{ExS}(\mathcal A,c_T)$, where
$\sfname{MinkDiv}$ says that for every $n$ the order of every finite subgroup of $\GL_n(\Q)$ divides
$\mathrm M(n)$, and $\sfname{ExS}(\mathcal A,c_T)$ says that $\sfname{DesignChoices}(\mathrm{WithTw}(\mathcal
A\wedge\mathcal G.\mathrm{cls}(\lambda)),3,c_T)$ for all $\lambda>1$ and all requirements $\mathcal G$ with
$\sfname{HostExW}(\lambda,\mathcal G)$.
\item \emph{Duty}: every $\topshapeF\in\mathcal A$ satisfies $\sfname{Duty}$, the conjunction of seven
conditions, among them $\sfname{TameDistinctF}$ and the wall condition $\sfname{ChosenWall}$.
\item \emph{Selections}: a twist selection $\sigma$, which assigns to $\topshapeF\in\mathcal A$ a threshold
$x_{\rm sel}(\topshapeF)$ and to each selected member a tied record, and dependent selections $\sigma_S$
(spectral ties), $\sigma_P$ (a base-level row-side family), $\sigma_{DH}$ (definite heights), $\sigma_C$
(column choices) and $\sigma_B$ (a bundle containing an indefinite realisation).
\item \emph{Conditions at selected members}: $h_{RB}$ (bounded rows of the indefinite realisation), $h_{Tw}$
(ties of the twisting data), $h_{PA}$ (all realisations primitive), $h_{\Lambda P}$, $h_{DivG}$, $h_{DcPtG}$ (row
conditions of $\sigma_P$), $h_{InstR}$ (independence of the volume of the signature, and class items),
$\sigma_J$ (a join of the indefinite realisations), and $h_{Mt}$: eventually the lattices $\widetilde M_\mu$ are
stable under $\OO_{\mathit{LL}}$.
\item \emph{Closed conditions}: $h_{Tu}$ ($c_T(\eps,C_{YL})\le C_{Tu}\Rightarrow\sfname{TuSize}$) and $h_{n4}$ (a
class-number formula for the levels) closed on $\mathrm{WithTw}(\mathrm{WithSm}(\mathrm{WithChN}(\mathcal
A)))$; and closed on $\topshapeD$: $h_{TF}$ (a trace frame), $s_O$ (Theorem~S for the selection), $h_{30}$ (a
point count), $j_{On}$ (the base vector image is full), $r_{\rm Rows}$ (rows along the selection), $a_{Pt0}$,
$a_{PtTw}$, $a_{Per0}$, $a_{PerTw}$, $a_{P10}$, $a_{AtG}$, $a_{C2}$ (integrality of points and periods, column
bounds at the level, used tuples and second-complement inputs), $\mathrm{irr}_C$ (\cref{pf:top-shape:irrthm}),
$z_L$ (surjectivity of level projections), $\mathrm{tie}$ (the real-place second-layer conditions), and
$n_{\rm Rest}\colon N(\topshapeD)$.
\item \emph{Global propositions}: $\sfname{LitThin}$, $\sfname{WuDecayOfGap}$, $\sfname{RowsFinDim}$,
$\sfname{TypeP}$, $\sfname{RowNGP}$, $\sfname{RowMainEqP}$, the loop rows $s_{Lp}$, $h_U$ (for every datum,
$\sfname{VolTie}\Rightarrow\sfname{VolumeRatio}_2$) and $h_G$ (for every datum, the genus mass in rank $2$).
\item \emph{Numerical data}: counting constants $c_{30}$ with $h_{c30}$, a real number $c_{Ic}>0$, and the
service hypothesis $h_{\rm Serv}$.
\item \emph{Parameters}: $h_{VC}\colon V$ and $n_R\colon N_R$.
\end{enumerate}
\formal{GlobalTheory.Fc.TopShapeAlt}
\end{definition}

\begin{theorem}[\cref{thm:shape}]\label{pf:top-shape:main}
Assume Green's theorem in the form $\sfname{Green}$ of \cref{def:green}. Then
\[
\sfname{TopShapeAlt}\bigl(\mathrm{True},\ \mathrm{True},\ \mathcal C\mapsto\topshapeCl(\mathcal C,\sfname{LinkN}),\
\sfname{TopG}\bigr).
\]
In words: for every family class $\mathcal A$ with data and hypotheses (a)--(g) of \cref{pf:top-shape:shape}
such that $\sfname{LinkN}$ is closed on $\topshapeD(\mathcal A)$, and for every $\lambda>1$ and every requirement
$\mathcal G$ with $\sfname{HostExW}(\lambda,\mathcal G)$ such that all families of $\mathcal G.\mathrm{cls}(\lambda)$
satisfy $\sfname{HypV}$, the class $\mathcal G.\mathrm{cls}(\lambda)$ is connected, with $e_D=2$ and $e_L=0$.
\formallabel{UpperTheory.Top.shape_closed_link_n_of_green}
\end{theorem}

\begin{remark}\label{pf:top-shape:rem:logic}
\begin{enumerate}[label=(\arabic*)]
\item The fields of an element of the package $\sfname{HypsAlt}$ of \cref{def:hyps}, with
$N=\topshapeCl(\cdot,\sfname{LinkN})$ and the definite-side condition, are exactly the data (a)--(g), the five
dependent selections being packaged as one selection, and the rows field being converted into $r_{\rm Rows}$; the
deduction of \cref{thm:shape} in the form stated in Part~I, from \cref{pf:top-shape:main} and an element of
the package, is \lref{GlobalTheory.Fc.forall_goal_of_shape_of_hyps}.
\item The parameters $V$ and $N_R$ are the true proposition; the hypotheses $h_{VC}$ and $n_R$ play no role.
\item The closedness of $\sfname{LinkN}$ on $\topshapeD(\mathcal A)$ is the parameter $N$ of the shape, i.e.\
a hypothesis; it is established in \cref{thm:linkN}.
\item Green's theorem is used only in the proof of $\sfname{LinkIRR}$ (\cref{pf:top-shape:irrthm}).
\end{enumerate}
\end{remark}

% ===========================================================================
\subsection{The assembly theorem and the proof of the shape theorem}\label{sec:pf:top-shape:assemblysec}

\subsubsection*{From the Sha hypothesis to $\sfname{DIVinN}$}

Each datum $D$ stores an algebraic closure $\Omega$ of its CM field $E$ and, for each prime $p$, local data
$\mathrm{ld}(p)$ of the kind quantified over in $\sfname{VVle}$, together with the tie: for $p$ prime and
$0<e<t$, if the length $\mathrm{shaLen}(\mathrm{ld}(p),e)$ is a finite number $n$, then $\mathrm{sh}_p(D,e)=n$.

\begin{lemma}\label{pf:top-shape:vvrecord}
Let $D$ be a datum and $B\in\R$ with $\sfname{VVle}(E,t,B)$. Then $\sfname{ShaFinOdd}(D)$ holds, the function
$\ell\mapsto\topshapeVV(D,\ell)$ is summable, and $\sum_\ell\topshapeVV(D,\ell)\le B$.
\formallabel{UpperTheory.Top.vvRecord}
\end{lemma}

\begin{proof}
Apply $\sfname{VVle}(E,t,B)$ to the algebraic closure and the local data stored in $D$. Its first clause is
$\sfname{ShaFinOdd}(D)$. For a prime $p$ the lengths $\mathrm{shaLen}(\mathrm{ld}(p),e)$, $0<e<t$, are therefore
finite, and the tie gives
\[
\sum_{e=1}^{t-1}\mathrm{sh}_p(D,e)=\sum_{e=1}^{t-1}\mathrm{shaLen}(\mathrm{ld}(p),e).
\]
Hence $\topshapeVV(D,\ell)\ge0$ for all $\ell$, and $\topshapeVV(D,\ell)=\log\ell\sum_e\mathrm{shaLen}(\mathrm{ld}(\ell),e)$
for $\ell$ prime. For a finite set $s$ of natural numbers, $\sum_{\ell\in s}\topshapeVV(D,\ell)$ equals the
sum over the primes of $s$, which is at most $B$ by the second clause of $\sfname{VVle}$. A non-negative function
on $\N$ whose finite partial sums are bounded by $B$ is summable with sum at most $B$.
\end{proof}

\begin{lemma}\label{pf:top-shape:divinat}
Let $(D,L)$ be a member, $j\in\mathrm{lifts}(D)$ and $C\in\R$. Assume the inequality of $\sfname{LinkN}$ at $j$
with constant $C$ (under summability), that $\ell\mapsto\topshapeVV(D,\ell)$ is summable with sum at most $b$,
that $d_{Z_0}(j)=D_0$, and that $(\log x)^2\le x^{\eps}$. Then
$\topshapeEta(L,j)\le D_0\,(|C|+1)\,x^{\eps}(B_0+N_c)$.
\formallabel{UpperTheory.Top.divinN_at}
\end{lemma}

\begin{proof}
Write $y=x^{\eps}=\Lambda\ge1$; then $B_0=fmby$, so $fmby^2=yB_0$. We have $f\ge1$, $m\ge1$ (as $j\in\{1,\dots,m\}$),
$b>0$ and $N_c\ge0$. The hypotheses give
\[
\topshapeEta(L,j)\le D_0\Bigl(C\bigl(yB_0+(\log x)^2N_c\bigr)+\sum_\ell\topshapeVV(D,\ell)\Bigr).
\]
Since $0\le yB_0+(\log x)^2N_c\le y(B_0+N_c)$, the first term is at most $|C|\,y(B_0+N_c)$. Since $fmy\ge1$ we
have $b\le B_0\le y(B_0+N_c)$, so the second term is at most $y(B_0+N_c)$. As $D_0\ge0$, the claim follows.
\end{proof}

\begin{lemma}\label{pf:top-shape:divinN}
Let $\eps>0$ and let $\topshapeF$ be a family with parameter $\eps$ satisfying $\sfname{HypV}$,
$\sfname{LinkN}$ and $\sfname{SmallLevelRankEq}$. Then $\topshapeF$ satisfies $\sfname{DIVinN}$, with constant
$|C|+1$ where $C$ is the constant of $\sfname{LinkN}$.
\formallabel{UpperTheory.Top.divinN_of_hypV}
\end{lemma}

\begin{proof}
Let $x_0$ be the maximum of the thresholds of $\sfname{HypV}$, $\sfname{LinkN}$ and $\sfname{SmallLevelRankEq}$,
of a threshold beyond which $(\log x)^2\le x^{\eps}$ (as $\eps>0$), and of $x_1+1$. Let $x\ge x_0$, $i\in I_x$
and $j\in\mathrm{lifts}(D_{x,i})$. By $\sfname{HypV}$ and \cref{pf:top-shape:vvrecord} with $B=\sqrt b$, the datum
satisfies $\sfname{ShaFinOdd}$ and $\topshapeVV$ is summable with sum at most $\sqrt b\le b$ (as $b\ge1$). So
the inequality of $\sfname{LinkN}$ applies, and \cref{pf:top-shape:divinat} (with $d_{Z_0}(j)=D_0$ from
$\sfname{SmallLevelRankEq}$, and the scale of $L_{x,i}$ equal to $(x,\eps)$ because $x>x_1$) gives the
inequality of $\sfname{DIVinN}$ with constant $|C|+1$.
\end{proof}

\subsubsection*{From $\sfname{LinkDIV}$ to $\sfname{LinkMu}$}

\begin{lemma}\label{pf:top-shape:linkmu}
Let $\topshapeF$ be a family with parameter $\eps>0$ satisfying $\sfname{LinkDIV}(C_{\rm abs},C_H)$ for some
$C_{\rm abs},C_H$, $\sfname{NuFullF}$, $\sfname{MuRank}$, $\sfname{TuSize}$ and $\sfname{DGamSize}$. Then
$\topshapeF$ satisfies $\sfname{LinkMu}$.
\formallabel{Formalism.MainNotions.J.linkMu_of_linkDIV}
\formallabel{Formalism.MainNotions.J.linkMu}
\formallabel{Formalism.MainNotions.muTu_of_tuSize}
\formallabel{Formalism.MainNotions.J.muDGam_of_fact}
\end{lemma}

\begin{proof}
We collect four eventual statements.
\begin{enumerate}[label=(\roman*)]
\item $\sfname{ZeroAbove}$ holds eventually: it is clause (3) of $\sfname{LinkDIV}$.
\item Eventually every $N_u$ is finitely generated (from $\sfname{NuFullF}$) and of positive rank (from
$\sfname{MuRank}$).
\item Eventually $\log T_u\le2X''$ for all words $u$. Indeed, as $\eps>0$, eventually $C_{Tu}\le x^{\eps/100}$,
and then $\sfname{TuSize}$ gives $\log T_u\le x^{\eps/100}m\Lambda=m\,x^{\eps/100+\eps}\le mx\le2nx\le2X''$,
using $\eps/100+\eps\le1$, $m\le2n$ and $nx\le X''(n)$.
\item There is $C$ such that eventually $\log D_\Gamma(u)\le10|C|X''$ for all $u$: by $\sfname{DGamSize}$,
$\log D_\Gamma(u)\le C\,fN(m)\log(fN(m))$, and the size facts of the family give
$0\le fN(m)\log(fN(m))\le10X''$.
\end{enumerate}
At a member where (i)--(iv) hold, \lref{Formalism.MainNotions.J.nprimeMuR_at} gives, for every $u$,
\[
\Bigl|\mu(N_u)-\mu(N^H_u)-\sum_w\mathrm{sz}(w,\sigma(u,w))\Bigr|\le2\bigl(\log4+3+2+10|C|\bigr)X'' .
\]
(That lemma states: if $\sfname{ZeroAbove}$ holds, the lattices $N_u$ are finitely generated of positive rank,
$\log T_u\le C_TX''$ and $\log D_\Gamma(u)\le C_GX''$ for all $u$, then, under the size facts $\sfname{SizeQD}$,
$\sfname{SizeG}$ of $L$ and the defining tie of $N^H_u$, the displayed difference is at most
$2(\log4+3+C_T+C_G)X''$.) Taking the maximum of the four thresholds and $x_1+1$ gives $\sfname{LinkMu}$.
\end{proof}

\subsubsection*{The assembly theorem}

\begin{theorem}[Assembly]\label{pf:top-shape:assembly}
Let $\mathcal C,\mathcal C',\mathcal E$ be family classes and $C_{\rm abs},C_0,e_D,e_L\in\R$. Assume:
\begin{enumerate}[label=(\arabic*)]
\item $\sfname{DesignExists}(\mathrm{WithTw}(\mathrm{WithSm}(\mathcal C')),C_0)$;
\item $\mathcal C'\subseteq\mathcal C$, $\mathcal C'\subseteq\mathcal E$, and every family of $\mathcal E$
satisfies $\sfname{HypV}$;
\item the following nine properties are closed on $\topshapeDes(\mathcal C,C_0)$: $\sfname{LinkN}$,
$\sfname{LinkEta}$, $\sfname{LinkIRR}$, $\sfname{SmallLevelRankEq}$,
$\sfname{TuSize}\Rightarrow\sfname{DGamSize}\Rightarrow\sfname{LinkDIVt}(C_{\rm abs},C_H)$,
$\sfname{ChosenE}(C_0,C_H)\Rightarrow\sfname{NuFullF}\wedge\sfname{MuRank}$, $\sfname{LinkG}$,
$\sfname{TuSize}\Rightarrow\sfname{LinkS}(C_0)$, and
$\sfname{TuSize}\Rightarrow\sfname{DGamSize}\Rightarrow\sfname{LinkGLOBAL}(C_{\rm abs},C_H,C_0,e_D,e_L)$.
\end{enumerate}
Then $\mathcal E$ is connected, with exponents $(e_D,e_L)$.
\formallabel{UpperTheory.Top.topDesWin}
\end{theorem}

\begin{proof}
By \cref{pf:top-shape:calculus}(b) the nine properties of (3) combine into one closed property; let $\eps_1>0$
be its constant, and let $\eps_3>0$ be the constant of (1). Put $\eps^*=\min(\eps_1,\eps_3)$ and let
$\eps\in(0,\eps^*)$. The combined closedness gives a staged threshold $K$ for $\eps$. By (1), applied to $\eps$
and $K$, there are a family $\topshapeF\in\mathrm{WithTw}(\mathrm{WithSm}(\mathcal C'))$ with parameter $\eps$ and
$C_H\ge0$ with $\topshapeLarge(\topshapeF,K,C_H)$, $\sfname{ChosenE}(\topshapeF,C_0,C_H)$, $\sfname{TuSize}$ and
$\sfname{DGamSize}$. As $\mathcal C'\subseteq\mathcal C$, the family lies in $\topshapeDes(\mathcal C,C_0)$; hence
the nine properties hold for $(\topshapeF,C_H)$.

As $\mathcal C'\subseteq\mathcal E$, $\topshapeF$ satisfies $\sfname{HypV}$; with $\sfname{LinkN}$ and
$\sfname{SmallLevelRankEq}$, \cref{pf:top-shape:divinN} gives $\sfname{DIVinN}$. With $\sfname{LinkEta}$,
$\sfname{TuSize}$ and $\sfname{DGamSize}$, the closed property $\sfname{LinkDIVt}$ gives
$\sfname{LinkDIV}(C_{\rm abs},C_H)$. From $\sfname{ChosenE}(C_0,C_H)$ we get $\sfname{NuFullF}$ and
$\sfname{MuRank}$, and \cref{pf:top-shape:linkmu} gives $\sfname{LinkMu}$. With $\sfname{LinkG}$ and
$\sfname{LinkS}(C_0)$ (from $\sfname{TuSize}$) and $\sfname{LinkIRR}$, this is
$\sfname{GlobalTheoremInputs}(C_{\rm abs},C_H,C_0)$. Now $\sfname{LinkGLOBAL}$, applied with
$\sfname{TuSize}$, $\sfname{DGamSize}$ and $\sfname{ChosenE}(C_0,C_H)$, yields $K',c_{HH},e_K,C_L,x_0$ such that
$\topshapeF$ is chain-tied with exponents $(e_D,e_L)$ at every window scale $x\ge x_0$. Since $\topshapeF\in\mathcal
E$, this proves connectedness with $\eps_1=\eps^*$.
\end{proof}

\subsubsection*{Proof of the shape theorem}

The link theorems are proved in the following subsections. Given them, the shape theorem is a direct
application of the assembly theorem.

\begin{proof}[Proof of \cref{pf:top-shape:main}]
Let $\mathcal A$, $c_T$ and the data (a)--(h) be given, with $n_{\rm Rest}\colon\topshapeCl(\topshapeD,
\sfname{LinkN})$ where $\topshapeD=\topshapeD(\mathcal A)$. Let $\lambda>1$ and let $\mathcal G$ be a requirement
with $h_H\colon\sfname{HostExW}(\lambda,\mathcal G)$ such that all families of $\mathcal G.\mathrm{cls}(\lambda)$
satisfy $\sfname{HypV}$. Apply \cref{pf:top-shape:assembly} with
\begin{gather*}
\mathcal C=\mathrm{WithChN}(\mathcal A),\qquad\mathcal C'=\mathrm{WithChN}(\mathcal A\wedge\mathcal G.\mathrm{cls}(\lambda)),
\qquad\mathcal E=\mathcal G.\mathrm{cls}(\lambda),\\
C_{\rm abs}=2,\qquad C_0=3,\qquad e_D=2,\qquad e_L=0 .
\end{gather*}
Then $\topshapeDes(\mathcal C,3)=\topshapeD$, and the inclusions of (2) are immediate from the definitions. The
inputs of (1) and (3) are:
\begin{itemize}
\item \emph{Design}: \cref{pf:top-shape:designthm}, applied to the class $\mathcal A\wedge\mathcal
G.\mathrm{cls}(\lambda)$ with the hypothesis $h_{\rm Ex}(\cdot)(\lambda)(\mathcal G)(h_H)$, the selection $\sigma$
and $\sfname{Duty}$ restricted to this class, the restrictions of $h_{Tu}$ and $h_{n4}$
(\cref{pf:top-shape:calculus}(d)) and $\sfname{LitThin}$.
\item $\sfname{LinkN}$: the hypothesis $n_{\rm Rest}$.
\item $\sfname{LinkEta}$: \cref{pf:top-shape:eta}.
\item $\sfname{LinkIRR}$: \cref{pf:top-shape:irrthm}, from Green's theorem, $c_{Ic}$ and $\mathrm{irr}_C$.
\item $\sfname{SmallLevelRankEq}$: \cref{pf:top-shape:small}.
\item $\sfname{LinkDIVt}$: \cref{pf:top-shape:divmain,pf:top-shape:bridge}.
\item $\sfname{ChosenE}\Rightarrow\sfname{NuFullF}\wedge\sfname{MuRank}$ and $\sfname{LinkGLOBAL}(2,C_H,3,2,0)$:
\cref{pf:top-shape:linkgl}.
\item $\sfname{LinkG}$: \cref{pf:top-shape:linkG}.
\item $\sfname{LinkS}(3)$: \cref{pf:top-shape:linkS}.
\end{itemize}
The assembly theorem shows that $\mathcal G.\mathrm{cls}(\lambda)$ is connected, which is the conclusion of
$\sfname{TopG}$.
\end{proof}

% ===========================================================================
\subsection{Green's theorem and the irreducibility link}\label{sec:pf:top-shape:irr}

\subsubsection*{Label weights}

Let $D$ be a datum and $k\in\{0,1\}$. Let $g=g_k$ be the $k$-th type prime of the host; it satisfies
$g\ge5$, and the ranks satisfy $r=\kappa+2s+1$ with $\kappa=k_D\ge1$ the rank $k$ of $D$ and $s=m/2-1$. A
\emph{label} is a monic polynomial $f\in\F_g[X]$ of degree $r$ with $f(0)\ne0$; a \emph{prime label of degree
$d$} is a monic irreducible polynomial of degree $d$ with non-zero constant term. The type data of $D$ attach to
every polynomial $f$ an element $e_{\rm Ser}(f)$ of the group algebra $\C[\GL_r(\F_g)]$, and for
$a\in\C[\GL_r(\F_g)]$ we put
\[
\topshapewt_k(a,f)=\frac{\Re\,(a\,e_{\rm Ser}(f))(1)}{\Re\,a(1)},
\]
$h(1)$ denoting the coefficient of the identity. For a non-trivial additive character $\psi_0$ of $\F_g$ let
$e_U=e_U(g,r,\psi_0)$ and $e_J=e_J(g,r,\kappa,s,\psi_0)$ be the Gelfand--Graev and Jacquet-type idempotents of
\cref{prop:families:admLq}.

\begin{definition}[Admissible label weights]\label{pf:top-shape:admLq}
$\sfname{admLq}(D)$ means: for $k\in\{0,1\}$ and every non-trivial additive character $\psi_0$ of $\F_{g_k}$,
\begin{enumerate}[label=(\arabic*)]
\item $\topshapewt_k(e_U,f)\ge0$ and $\topshapewt_k(e_J,f)\ge0$ for all $f\in\F_g[X]$;
\item $\topshapewt_k(e_U,f)=\prod_{P}(g^{\deg P}-1)^{-1}$ for every squarefree label $f$, the product running
over the monic irreducible factors $P$ of $f$;
\item $(1-g^{-1})\,\topshapewt_k(e_U,f)\le\topshapewt_k(e_J,f)$ for every label $f$;
\item $\topshapewt_k(e_J,f)\le(1+g^{-1})^r\,\topshapewt_k(e_U,f)$ for every squarefree label $f$.
\end{enumerate}
\formal{Formalism.MainNotions.ArithmeticCylinders.admLq}
\end{definition}

A \emph{label datum} for $g$ consists of an injective multiplicative map $\mathrm{teich}$ from
$\overline\F_g^\times$ to $\C^\times$ and norm-compatible generators of the groups $\F_{g^d}^\times$; Green's
theorem (\cref{def:green}) is formulated relative to such a datum, and attaches to every prime label $P$ of
degree $d$ a cuspidal irreducible character $\tau(P)$ of $\GL_d(\F_g)$. The type data of $D$ also label the
simple cuspidal representations $V$ of $\GL_d(\F_g)$ by prime labels $O(V)$ of degree $d$, injectively up to
isomorphism.

\begin{lemma}\label{pf:top-shape:perm}
Assume $\sfname{Green}$. Let $\mathcal L$ be a label datum for $g$ and $O$ a labelling of the cuspidal
representations as above. There are maps $\sigma,\pi$ of $\F_g[X]$ to itself such that for every prime label $P$
of degree $d$, $\sigma(P)$ is a prime label of degree $d$ and $\pi(\sigma(P))=P$, and such that for every $d\ge1$
and every simple cuspidal representation $V$ of $\GL_d(\F_g)$ the character of $V$ is $\tau(\sigma(O(V)))$.
\formallabel{GlobalTheory.Top.Lq.exists_perm}
\formallabel{GlobalTheory.Top.Lq.exists_perm_deg}
\end{lemma}

\begin{proof}
It suffices to construct, for each $d\ge1$, maps $\sigma_d,\pi_d$ with these properties for prime labels of
degree $d$ and representations of $\GL_d(\F_g)$, and to put $\sigma(P)=\sigma_{\deg P}(P)$,
$\pi(P)=\pi_{\deg P}(P)$. Fix $d\ge1$ and let $\tau$ be given by $\sfname{Green}$ for $\mathcal L$ and $d$. For each
prime label $P$ of degree $d$, $\tau(P)$ is a cuspidal irreducible character; choose a simple representation
$W(P)$ with character $\tau(P)$. It is cuspidal, because for $\GL_d(\F_g)$ cuspidality of a simple
representation is equivalent to the cuspidality of its character
(\lref{GlobalTheory.T.isCuspidalCharacter_iff_isCuspidal_GLnFiniteField}). Put $p(P)=O(W(P))$, a prime label of
degree $d$. The map $p$ is injective: if $O(W(P))=O(W(Q))$, then $W(P)\cong W(Q)$, so $\tau(P)=\tau(Q)$ and $P=Q$ by
the injectivity of $\tau$. As there are finitely many prime labels of degree $d$, $p$ is a bijection; let
$\sigma_d=p^{-1}$ and $\pi_d=p$ on prime labels of degree $d$ (and the identity elsewhere). Let $V$ be a simple
cuspidal representation. By the surjectivity part of $\sfname{Green}$, its character is $\tau(P)$ for some prime
label $P$ of degree $d$. Representations of a finite group with equal characters are isomorphic, so $V\cong
W(P)$, hence $O(V)=p(P)$ and $\sigma_d(O(V))=P$, i.e.\ the character of $V$ is $\tau(\sigma_d(O(V)))$.
\end{proof}

Extend $\sigma$ multiplicatively: for monic $f$ with $f(0)\ne0$ let $\sigma(f)$ be the product of the $\sigma(P)$
over the monic irreducible factors $P$ of $f$, counted with multiplicity. Then $\sigma$ maps labels to labels and
squarefree labels to squarefree labels, the factors of $\sigma(f)$ are the $\sigma(P)$, of the same degrees as the
$P$, $\sigma$ is multiplicative on products of prime labels, and $\pi(\sigma(f))=f$. Labels of irreducible
representations of $\GL_r(\F_g)$ are defined through Harish-Chandra induction. An irreducible character $\chi$ has
\emph{Green label} $f$ if there are cuspidal characters $\chi_1,\dots,\chi_s$ of groups $\GL_{d_c}(\F_g)$,
$\sum_cd_c=r$, with $\chi_c=\tau(P_c)$ for prime labels $P_c$, such that $f=\prod_cP_c$ and $\chi$ occurs in the
Harish-Chandra induction of $\chi_1\otimes\dots\otimes\chi_s$; by
\lref{GlobalTheory.T.harishChandraSeriesLabel_exists_unique_GLn}, for $g\ge5$ under $\sfname{Green}$ every
irreducible character has exactly one Green label, labels are monic of degree $r$ with non-zero constant term, and
the irreducible characters with a squarefree label $f$ form a single series. A simple representation has
\emph{type-data label} $f$ if it occurs in the Harish-Chandra induction of simple cuspidal representations
$\lambda_1,\dots,\lambda_s$ (in blocks of sizes $d_c$) with $f=\prod_cO(\lambda_c)$. Let $e_{\mathcal L}(f)$ be the sum
of the central idempotents of the irreducible characters with Green label $f$, and
$\topshapewt_{\mathcal L}(a,f)=\Re(a\,e_{\mathcal L}(f))(1)/\Re a(1)$.

\begin{lemma}\label{pf:top-shape:eser}
In the situation of \cref{pf:top-shape:perm}, with $g\ge5$ and $\mathcal L$, $O$ the label datum and labelling given
by normalized type data of $D$, one has $e_{\rm Ser}(f)=e_{\mathcal L}(\sigma(f))$ for every label $f$, and
$e_{\rm Ser}(f)=0$ if $f$ is not a label.
\formallabel{GlobalTheory.Top.Lq.eSer_eq}
\formallabel{GlobalTheory.Top.Lq.hasLabel_of}
\formallabel{GlobalTheory.Top.Lq.exists_labelO}
\end{lemma}

\begin{proof}
For labels, $e_{\rm Ser}(f)$ is by definition the sum of the central idempotents of the simple representations
whose type-data label is $f$. It suffices to show that a simple representation $V$ has type-data label $f$ if and
only if its character has Green label $\sigma(f)$. If $V$ has type-data label $f$, realised by cuspidal
$\lambda_1,\dots,\lambda_s$, then by \cref{pf:top-shape:perm} the characters of the $\lambda_c$ are
$\tau(\sigma(O(\lambda_c)))$, and $\sigma(f)=\prod_c\sigma(O(\lambda_c))$; the character of the Harish-Chandra induced
representation is the product character of the $\chi_c$, and occurrence of $V$ is equivalent to the non-vanishing of
the pairing of characters (\lref{GlobalTheory.T.trace_harishChandra_induced_eq_product_characters} and \lref{GlobalTheory.T.isConstituent_harishChandraRep_iff_pairing_ne_zero});
so the character of $V$ has Green label $\sigma(f)$. Conversely, if the character of $V$ has Green label
$\sigma(f)$, the cuspidal characters of its Green label are characters of simple cuspidal representations
$\lambda_c$ (\lref{GlobalTheory.T.exists_blocks_of_cuspidal_character_list}), and $V$ occurs in their
Harish-Chandra induction; hence $V$ has the type-data label $f'=\prod_cO(\lambda_c)$, monic with $f'(0)\ne0$. By
the first part its character has Green label $\sigma(f')$; by uniqueness of Green labels $\sigma(f')=\sigma(f)$,
and applying $\pi$ gives $f'=f$. The vanishing for non-labels is a tie of the type data
(\lref{Formalism.L2.eSer_zero_of_tie}).
\end{proof}

\begin{proposition}[\cref{prop:families:admLq}]\label{pf:top-shape:admLqthm}
Assume $\sfname{Green}$. Let $D$ be a datum, $j_E$ an embedding of its CM field into $\Qbar$, and $T$ a normalized
types compatibility record for $(D,j_E)$. Then $\sfname{admLq}(D)$ holds.
\formallabel{GlobalTheory.Top.Lq.admissibleLabelWeights_of_greenCuspidalParametrisation}
\end{proposition}

\begin{proof}
Fix $k$, $g=g_k$ and $\psi_0\ne1$. The record $T$ provides a label datum $\mathcal L$ (the Teichm\"uller map of
the type data and compatible generators) and a labelling $O$ of the cuspidal representations; let $\sigma$ be as
in \cref{pf:top-shape:perm}. By \cref{pf:top-shape:eser}, $\topshapewt_k(a,f)=\topshapewt_{\mathcal L}(a,\sigma(f))$
for labels $f$ and $\topshapewt_k(a,f)=0$ otherwise. We use four results on $\GL_N(\F_g)$, proved in the
development and shared with other parts of the proof:
\begin{itemize}
\item $e_U^2=e_U$, $e_U(1)=g^{-r(r-1)/2}$, and, since $r=\kappa+2s+1$, $e_J^2=e_J$ and $e_J(1)=g^{\kappa}g^{-r(r-1)/2}$
(\lref{GlobalTheory.T.whittakerJacquetIdempotents_properties_primeField});
\item for an idempotent $a$ with $\Re a(1)\ge0$, $\topshapewt_{\mathcal L}(a,f)\ge0$ for all $f$
(\lref{GlobalTheory.T.harishChandraSeriesWeight_nonneg_of_idempotent_glnPrime});
\item for $g\ge5$, under $\sfname{Green}$ and $\psi_0\ne1$: for every squarefree monic $f$ of degree $N$ with
$f(0)\ne0$, $\topshapewt_{\mathcal L}(e_U(g,N,\psi_0),f)=\prod_P(g^{\deg P}-1)^{-1}$
(\lref{GlobalTheory.T.harishChandraSeriesWeight_squarefree_label_gelfandGraev});
\item for $N=\kappa+2s+1$, $\kappa\ge1$ and $\psi_0\ne1$:
$(1-g^{-1})\topshapewt_{\mathcal L}(e_U,f)\le\topshapewt_{\mathcal L}(e_J,f)$ for all $f$
(\lref{GlobalTheory.T.harishChandraSeriesWeight_lower_comparison}), and, for $g\ge5$ under $\sfname{Green}$,
$\topshapewt_{\mathcal L}(e_J,f)\le(1+g^{-1})^N\topshapewt_{\mathcal L}(e_U,f)$ for squarefree labels $f$
(\lref{GlobalTheory.T.harishChandraSeriesWeight_le_squarefree}).
\end{itemize}
(1): for labels, apply the second item to $e_U$ and $e_J$, which are idempotents with $e_U(1),e_J(1)>0$ by the
first; for non-labels both weights vanish. (2): if $f$ is a squarefree label, so is $\sigma(f)$, and its factors
$\sigma(P)$ have the degrees of the factors $P$ of $f$; the third item gives the formula. (3) and (4) follow from
the fourth item applied to $\sigma(f)$.
\end{proof}

For squarefree $f$ the classical computation behind (2) is the following: the Harish-Chandra series of $f$
contains exactly one generic representation, which occurs with multiplicity one in the Gelfand--Graev
representation \cite{GelfandGraev1962,Zelevinsky1981} and has degree
$\prod_{i\le r}(g^i-1)/\prod_P(g^{\deg P}-1)$ by Green's theorem \cite{Green1955}.

\subsubsection*{Counting with all-rows records}

For $k\in\{0,1\}$ the type data of $D$ provide a finite set $\mathcal P_k$ of $f$ places and characters
$\chi_k\colon\mathcal P_k\times\{1,\dots,m\}\to\F_{g_k}^\times$. A \emph{label tuple} at $k$ is a map
$t\colon\mathcal P_k\to\F_{g_k}[X]$. It is \emph{good} at $k$ if it is good for $\chi_k$ and for $\chi_k^{-1}$, where
good for $\chi$ means: every $t(v)$ is a squarefree label; for every prime $\ell$ outside an explicit finite set,
the prime-to-$\ell$ parts of the roots of each $t(v)$ in $\overline\F_g$ are pairwise distinct; for every prime
$\ell\ne g$ the prime-to-$\ell$ parts of the roots of $t(v)$ differ from that of $\chi(v,j)$, for all $j$; for every
$1\le a\le r/2$ some $t(v)$ has no factor of degree $a$; and some $t(v)$ has an irreducible factor of prime degree
$q$ with $r<2q$ and $q+3\le r$. A \emph{cylinder} is a set $\mathrm{cyl}(k,S)$ of pairs $(t_0,t_1)$ of label tuples
with $t_k\in S$.

An \emph{all-rows record} $A$ for $D$ consists of rows data, label tuples $\mathrm{lab}(k,\tau)$ (all of whose
entries are labels) for the $D_{\rm all}$ indices $\tau$, an injection $\mathrm{inc}$ of the $D_{\rm rows}$ rows into
the indices, idempotents $e_{k,v}\in\C[\GL_r(\F_{g_k})]$ ($v\in\mathcal P_k$), group-algebra functions
$e_T(S)$ attached to sets $S$ of pairs of label tuples with $e_T(\text{all})=e$, dimension counts
$\dim_A(S)\in\N$, and the normalisation
\[
\Re\,e_T(\mathrm{cyl}(k,S))(1)=\Re\,e(1)\cdot\mathrm{sh}_k(S),\qquad
\mathrm{sh}_k(S)=\sum_{t\in S}\prod_{v\in\mathcal P_k}\topshapewt_k(e_{k,v},t(v)).
\]
The \emph{expected value} $D^*(h)$ of a group-algebra function $h$ is the main term of the rows data (formal
degree times volume times $\Re h(1)$); put $D_*=D^*(e)$, so that $D^*(e_T(\mathrm{cyl}(k,S)))=D_*\,\mathrm{sh}_k(S)$.

\begin{definition}[Conditions on all-rows records]\label{pf:top-shape:allrows}
\begin{itemize}
\item $\sfname{LabelCountGoodInjective}(A)$: $\dim_A(S)$ is the number of indices $\tau$ with
$(\mathrm{lab}(0,\tau),\mathrm{lab}(1,\tau))\in S$; $\mathrm{inc}$ is injective; the labels of $\mathrm{inc}(\rho)$ are the
labels of the row $\rho$; and the pair of label tuples of $\tau$ is good (at both $k$) if and only if $\tau$ is in
the image of $\mathrm{inc}$.
\item $\sfname{CountAllRows}(A,\eps_C)$: $0\le\eps_C<1$, $D_*>0$, and
$\bigl|\dim_A(\mathrm{cyl}(k,S))-D_*\mathrm{sh}_k(S)\bigr|\le\eps_C\,D_*\mathrm{sh}_k(S)$ for all $k$ and $S$.
\item $\sfname{WeightAllRows}(A)$: all weights $\topshapewt_k(e_{k,v},f)$ are $\ge0$, and
$\mathrm{sh}_k(S)\le W_k$ for every finite set $S$ of label tuples none of which is good at $k$, where
\begin{multline*}
W_k=\frac{3f}{g_k}+\frac{23f\log g_k}{\sqrt{g_k}}+2fm\Bigl(\frac1{4\min(q,d)}+\frac{1+\log r}{g_k-1}\Bigr)
+\frac r2\Bigl(\frac78\Bigr)^f+\exp\Bigl(-\frac{f}{5\log r}\Bigr),
\end{multline*}
$q,d$ being the primes of the host.
\item $\sfname{AllJoinA}(A)$: a rows tie and a sharp level condition; it implies that there are non-trivial
additive characters $\psi_0$ of $\F_{g_k}$ such that every $e_{k,v}$ is $e_U$ or $e_J$ for $\psi_0$
(\lref{Formalism.L2.GlobalTheoryTheory.AllJoin.tieVecW_of_join}).
\item $\sfname{SizesRec}(D,x,\eps,C_s)$: for $k=0,1$, $\exp(8x^{\eps})\le g_k$, $x^2\le f\le x^{C_s}$,
$m\le x^{C_s}$, $100\le r\le x^{C_s}$ and $\exp(x^{\eps/6})\le\min(q,d)$.
\end{itemize}
\formal{Formalism.MainNotions.ArithmeticCylinders.LabelCountGoodInjective}
\formal{Formalism.L2.GlobalTheoryTheory.ConcN.CountAllRows}
\formal{Formalism.L2.GlobalTheoryTheory.ConcN.WeightAllRows}
\formal{Formalism.L2.GlobalTheoryTheory.ConcN.SizesRec}
\end{definition}

\begin{lemma}\label{pf:top-shape:transfer}
Let $\iota$ be a finite set, $(T_v)_{v\in\iota}$ finite sets, $g\ge1$, and $w_U,w_J$ functions on $\bigsqcup_vT_v$
with $w_U\ge0$, $\sum_{t\in T_v}w_U(v,t)=\sum_{t\in T_v}w_J(v,t)=1$ and $w_J\ge(1-g^{-1})w_U$. For every set
$\mathcal A$ of tuples $x=(x_v)\in\prod_vT_v$,
\[
\sum_{x\in\mathcal A}\prod_vw_J(v,x_v)\le\sum_{x\in\mathcal A}\prod_vw_U(v,x_v)+\frac{|\iota|}{g}.
\]
\formallabel{Formalism.IrreducibleCountingLemmas.GaloisIrreducibility.comparison_transfer_holds}
\end{lemma}

\begin{proof}
Put $\alpha=1-g^{-1}\in[0,1]$ and write $\mu_U,\mu_J$ for the two product measures; both have total mass $1$, and
$\mu_J\ge\alpha^{|\iota|}\mu_U$ pointwise. Hence
$\mu_J(\mathcal A)=1-\mu_J(\mathcal A^c)\le1-\alpha^{|\iota|}(1-\mu_U(\mathcal A))\le(1-\alpha^{|\iota|})+\mu_U(\mathcal
A)$, and $1-\alpha^{|\iota|}\le|\iota|/g$ by Bernoulli's inequality.
\end{proof}

\begin{lemma}\label{pf:top-shape:weight}
Let $A$ be an all-rows record for $D$ with $\sfname{AllJoinA}(A)$, and assume $\sfname{admLq}(D)$, $100\le r$,
$2^{12}\le g_k$ ($k=0,1$), and that $q,d$ are distinct odd primes with $16\le\max(q,d)$ dividing $g_k-1$, as does
$16$. Then $\sfname{WeightAllRows}(A)$ holds.
\formallabel{Formalism.L2.GlobalTheoryTheory.AllJoin.weight_of_join}
\formallabel{Formalism.L2.GlobalTheoryTheory.ConcNWeight.weightAllRows_of}
\formallabel{Formalism.L2.GlobalTheoryTheory.ConcNWeight.locShare_le_Wg}
\end{lemma}

\begin{proof}
By $\sfname{AllJoinA}$, each $e_{k,v}$ is $e_U$ or $e_J$ for a non-trivial $\psi_0$, so the weights are $\ge0$ by
$\sfname{admLq}$(1). Fix $k$ and a finite set $S$ of non-good tuples. Identify a label with its coefficient vector
$c\in\F_g^r$ (a label is $X^r+\sum_{i<r}c_iX^i$), and put $P(c)=\topshapewt_k(e_U,f_c)$ and
$w(v,c)=\topshapewt_k(e_{k,v},f_c)$, $f_c$ being the monic polynomial with coefficient vector $c$. Since the series
idempotents $e_{\rm Ser}(f)$ sum to $1$, $\sum_cP(c)=\sum_cw(v,c)=1$. By $\sfname{admLq}$(3) (and trivially when
$e_{k,v}=e_U$, or when $f_c$ is not a label, where both weights vanish), $w\ge(1-g^{-1})P$. In $\mathrm{sh}_k(S)$ the
tuples with a non-label entry contribute $0$; the other tuples of $S$ correspond to coefficient tuples which are
not good. Hence $\mathrm{sh}_k(S)$ is at most the $\prod_vw(v,\cdot)$-measure of the set of non-good coefficient
tuples, and by \cref{pf:top-shape:transfer} (with $|\iota|=|\mathcal P_k|=f$) at most its $P^{\otimes\mathcal
P_k}$-measure plus $f/g$. By $\sfname{admLq}$(1),(2), $P$ is a probability vector equal to $\prod_P(g^{\deg P}-1)^{-1}$
on the coefficient vectors of squarefree labels. Under the stated conditions on $g,q,d,r$,
\lref{Formalism.IrreducibleCountingLemmas.GaloisIrreducibility.property_W_of_estimates} bounds the
$P^{\otimes\mathcal P_k}$-measure of the non-good tuples (for $\chi_k$ with values in $\F_g^\times$, indexed by
$\mathcal P_k\times\{1,\dots,m\}$) by
\[
\frac{2f}{g}+\frac{23f\log g}{\sqrt g}+2fm\Bigl(\frac1{4\min(q,d)}+\frac{1+\log r}{g-1}\Bigr)+\frac r2\Bigl(\frac78\Bigr)^f
+\exp\Bigl(-\frac f{5\log r}\Bigr).
\]
Adding $f/g$ gives $W_k$.
\end{proof}

\begin{lemma}\label{pf:top-shape:delp}
Let $A$ be an all-rows record for $D$ with $\sfname{LabelCountGoodInjective}(A)$, $\sfname{CountAllRows}(A,\eps_C)$
and $\sfname{WeightAllRows}(A)$, and let $x,\eps$ be real numbers with
$\frac{1+\eps_C}{1-\eps_C}(W_0+W_1)\le\exp(-x^{\eps/7})$. Then $0\le\delta_p(D)\le\exp(-x^{\eps/7})$.
\formallabel{Formalism.L2.GlobalTheoryTheory.ConcN.delp_of_allRows}
\formallabel{Formalism.IrreducibleCountingLemmas.GaloisIrreducibility.share_from_count}
\end{lemma}

\begin{proof}
For $k=0,1$ let $S_k$ be the set of label tuples $\mathrm{lab}(k,\tau)$ which are not good at $k$, and $B_k$ the number
of indices $\tau$ whose tuple at $k$ is not good. By $\sfname{LabelCountGoodInjective}$, $\dim_A(\mathrm{cyl}(k,S_k))=B_k$,
the number of good indices is $D_{\rm rows}$ (the good indices are the image of the injection $\mathrm{inc}$), and an
index which is not good is not good at $0$ or at $1$; hence $D_{\rm all}-D_{\rm rows}\le B_0+B_1$. The count for the
cylinder over all tuples gives $|D_{\rm all}-D_*|\le\eps_CD_*$, so $D_{\rm all}\ge(1-\eps_C)D_*>0$. The counts for the
cylinders over $S_k$ give $B_k\le(1+\eps_C)\mathrm{sh}_k(S_k)D_*\le(1+\eps_C)W_kD_*$, using $0\le\mathrm{sh}_k(S_k)\le
W_k$ from $\sfname{WeightAllRows}$. Therefore
\[
\frac{D_{\rm all}-D_{\rm rows}}{D_{\rm all}}\le\frac{(1+\eps_C)(W_0+W_1)D_*}{(1-\eps_C)D_*}\le\exp(-x^{\eps/7}).
\]
Since $\delta_pD_{\rm all}=D_{\rm all}-D_{\rm rows}$ and $D_{\rm all}>0$, the bound for $\delta_p$ follows; and
$\delta_p\ge0$ because $D_{\rm rows}\le D_{\rm all}$, $\mathrm{inc}$ being injective.
\end{proof}

\begin{lemma}\label{pf:top-shape:numbers}
\begin{enumerate}[label=(\alph*)]
\item For $c>0$, for all large $x$: if $N\ge100$ and $f\ge x^2$, then $\exp(-cNf\log N)\le\frac12$.
\item For $0<\eps\le1$ and $C_s$, for all large $x$: if $\sfname{SizesRec}(D,x,\eps,C_s)$ then $W_k\le
x^{3C_s+3}\exp(-x^{\eps/6})$ for $k=0,1$.
\item For $C\in\R$ and $\eps>0$, for all large $x$: $x^C\exp(-x^{\eps/6})\le\frac1{20}\exp(-x^{\eps/7})$.
\item Consequently, for $C$ and $\eps>0$, for all large $x$: if $0\le\eps_C\le\frac12$ and
$W_0,W_1\le x^C\exp(-x^{\eps/6})$, then $\frac{1+\eps_C}{1-\eps_C}(W_0+W_1)\le\exp(-x^{\eps/7})$.
\end{enumerate}
\formallabel{Formalism.L2.GlobalTheoryTheory.ConcN.epsilon_bound_large_x}
\formallabel{Formalism.IrreducibleCountingLemmas.GaloisIrreducibility.dominant_numbers_sharp_holds}
\formallabel{Formalism.L2.GlobalTheoryTheory.ConcN.numbers_dominant_large}
\end{lemma}

\begin{proof}
(a) For $x\ge\max(1,1/c)$ we have $cf\ge cx^2\ge1$ and $N\log N\ge1$, so $cNf\log N\ge1$ and
$\exp(-cNf\log N)\le e^{-1}<\frac12$. (b) is \lref{Formalism.L2.GlobalTheoryTheory.ConcN.weight_numbers_bound}:
each of the five terms of $W_k$ is bounded by a multiple of $x^{3C_s+3}\exp(-x^{\eps/6})$, using $g\ge\exp(8x^\eps)$
for the first three terms, $\min(q,d)\ge\exp(x^{\eps/6})$, and $f\ge x^2$ with $r\le x^{C_s}$ for the last two.
(c) For large $x$, $x^{\eps/6}\ge3x^{\eps/7}$, $C\log x\le\frac13x^{\eps/6}$ and $x^{\eps/6}\ge3\log20$; then
$C\log x-x^{\eps/6}\le-\log20-x^{\eps/7}$. (d) If $\eps_C\le\frac12$ then $\frac{1+\eps_C}{1-\eps_C}\le3$; if
$W_0+W_1\ge0$ the left side is at most $6x^C\exp(-x^{\eps/6})\le\frac6{20}\exp(-x^{\eps/7})$ by (c), and otherwise it
is $\le0$.
\end{proof}

\begin{theorem}[\cref{thm:linkIRR}]\label{pf:top-shape:irrthm}
Assume $\sfname{Green}$. Let $\mathcal A$ be a family class with a twist selection $\sigma$, let $c_{Ic}>0$, and
assume that the following property $\mathrm{irr}_C$ is closed on $\topshapeD(\mathcal A)$: eventually every member
$D=D_{x,i}$ admits an all-rows record $A$ with $\sfname{LabelCountGoodInjective}(A)$, $\sfname{AllJoinA}(A)$ and
$\sfname{CountAllRows}(A,\exp(-c_{Ic}\,rf\log r))$. Then $\sfname{LinkIRR}$ is closed on $\topshapeD(\mathcal A)$.
\formallabel{GlobalTheory.Top.Irr.closed_linkIRR_of_join_count}
\formallabel{GlobalTheory.Top.Irr.linkIRR_of_join_ex}
\formallabel{Formalism.L2.GlobalTheoryTheory.ConcN.linkIRR_of_allRows}
\formallabel{Formalism.L2.GlobalTheoryTheory.ConcN.link_irreducible_all_rows}
\end{theorem}

\begin{proof}
\emph{Step 1: the records carry admissible weights and sizes.} Let $\topshapeF\in\topshapeD(\mathcal A)$ and let $x_0$
be the threshold of $\mathrm{irr}_C$. For $x\ge\max(x_0,x_{\rm sel}(\topshapeF),x_1+1,4)$ and $i\in I_x$, the selection
provides a tied record whose normalized types give, by \cref{pf:top-shape:admLqthm}, $\sfname{admLq}(D_{x,i})$; and
the size facts of the family give $\sfname{SizesRec}(D_{x,i},x,\eps_z,6)$
(\lref{GlobalTheory.Top.Irr.sizesRec_at}). By \cref{pf:top-shape:calculus}(c),(e) the property ``eventually
every member admits $A$ with $\sfname{LabelCountGoodInjective}$, $\sfname{AllJoinA}$, $\sfname{admLq}$,
$\sfname{CountAllRows}(A,\exp(-c_{Ic}rf\log r))$ and $\sfname{SizesRec}(D,x,\eps_z,6)$'' is closed on
$\topshapeD(\mathcal A)$. So is $\eps\le1$ (take $\eps_1=1$). By \cref{pf:top-shape:calculus}(b),(c) it suffices to
show that a family with parameter $0<\eps\le1$ having this property satisfies $\sfname{LinkIRR}$.

\emph{Step 2: weights.} For $x\ge(3/2)^{1/\eps_z}$ we have $x^{\eps_z}\ge3/2$, and $\sfname{SizesRec}$ then gives
$r\ge100$ and $g_k\ge\exp(8x^{\eps_z})\ge2^{12}$; the host conditions on $q,d$ of \cref{pf:top-shape:weight} are ties of
the type data. Hence \cref{pf:top-shape:weight} gives $\sfname{WeightAllRows}(A)$.

\emph{Step 3: numbers.} Let $\eps_C=\exp(-c_{Ic}rf\log r)$. By \cref{pf:top-shape:numbers}(a) (with $N=r\ge100$
and $f\ge x^2$ from $\sfname{SizesRec}$), (b) (with $\eps_z\le1$ and $C_s=6$) and (d) (with $C=3\cdot6+3$), for all
large $x$ we have $\frac{1+\eps_C}{1-\eps_C}(W_0+W_1)\le\exp(-x^{\eps_z/7})$. By \cref{pf:top-shape:delp},
$0\le\delta_p(D_{x,i})\le\exp(-x^{\eps_z/7})$ for all large $x$ and all $i$, which is $\sfname{LinkIRR}$.
\end{proof}

% ===========================================================================
\subsection{The links \texorpdfstring{$\sfname{LinkG}$, $\sfname{LinkEta}$, $\sfname{SmallLevelRankEq}$ and
$\sfname{LinkS}$}{LinkG, LinkEta, SmallLevelRankEq and LinkS}}\label{sec:pf:top-shape:links}

\subsubsection*{The local lattice link}

For a member $(D,L)$ with group instance $X=G_x(x,i)$, let $\mathcal P(D,L)$ be the set of primes $p<X'$ in the
residue class attached to the host (the class ``zero'') such that the additive order of the element $h_P(p)$ is
at least $x^{\eps/4}$, and for such $p$ let $\mathrm{bad}(p)\subseteq TT$ be the set of tops $t$ at which the
instance is not good at $p$. Put $B_{TT}=\sum_{p\in\mathcal P}\frac{\log p}p|\mathrm{bad}(p)|$. The condition
$\sfname{Adm2F}$ of \cref{thm:linkG} says that eventually
\[
\tfrac12n\delta_1\le|TT|\le2n\delta_1\qquad\text{and}\qquad B_{TT}\le|TT|\,x^{-2+\eps}.
\]
The choice condition $\sfname{ChosenBad}(C_b)$ says that eventually
\[
\sum_{p\in\mathcal P}\frac{\log p}p|\mathrm{bad}(p)|\le4C_bx^{-2}\delta_1n\sum_{p\in\mathcal P}\frac{\log p}p .
\]

\begin{lemma}[A Mertens-type estimate]\label{pf:top-shape:mertens}
Let $\topshapeF$ be a family with parameter $\eps$ satisfying $\sfname{ChosenTT}$ and $\sfname{ChosenBad}(C_b)$ for
some $C_b>0$. Then $\topshapeF$ satisfies $\sfname{Adm2F}$.
\formallabel{Formalism.MainNotions.Pk.admissible_tops_mertens}
\formallabel{Formalism.MainNotions.Pk.sum_log_div_le_two_mul_mul_add_two}
\end{lemma}

\begin{proof}
First, for $T\ge0$ and every set $s$ of integers $p$ with $1\le p<\lceil e^T\rceil$,
\begin{equation}\label{eq:pf:top-shape:mertens}
\sum_{p\in s}\frac{\log p}p\le2T(T+2).
\end{equation}
Indeed $\log p\le T$ and $1/p\le2/(p+1)$, and with $M=\lceil e^T\rceil\le e^T+1$ we have
$\sum_{0\le p<M}(p+1)^{-1}=H_M\le1+\log M\le T+2$.

Let $T=x^{\eps/3}$, so that $X'=e^T$, and let $x$ be beyond the thresholds of $\sfname{ChosenTT}$ and
$\sfname{ChosenBad}$, beyond $x_1+1$ and with $T\ge16C_b+2$ (i.e.\ $x\ge(16C_b+2)^{3/\eps}$). The first condition of
$\sfname{Adm2F}$ is $\sfname{ChosenTT}$. Every prime counted in $B_{TT}$ lies in $\mathcal P$, so by
$\sfname{ChosenBad}$ and \eqref{eq:pf:top-shape:mertens}
\[
B_{TT}\le4C_bx^{-2}\delta_1n\cdot2T(T+2).
\]
Since $T\ge16C_b+2$ we have $16C_b(T+2)\le T^2$, hence $8C_bT(T+2)\le\frac12T^3$; with $|TT|\ge\frac12\delta_1n$
and $x^{-2+\eps}=x^{-2}T^3$ this gives $B_{TT}\le\frac12\delta_1n\,x^{-2}T^3\le|TT|\,x^{-2+\eps}$.
\end{proof}

\begin{proposition}\label{pf:top-shape:linkG}
Let $\mathcal A$ be a family class with a twist selection $\sigma$ all of whose families satisfy $\sfname{Duty}$.
Then $\sfname{LinkG}$ is closed on $\topshapeD(\mathcal A)$.
\formallabel{GlobalTheory.Top.closed_linkG_of_duty}
\end{proposition}

\begin{proof}
Every member of $\topshapeD(\mathcal A)$ satisfies $\sfname{ChosenE}(3,C_H)$ for some $C_H$, hence
$\sfname{ChosenTT}$ and $\sfname{ChosenBad}(C_b)$ with $C_b>0$, hence $\sfname{Adm2F}$ by \cref{pf:top-shape:mertens};
so $\sfname{Adm2F}$ is closed on $\topshapeD(\mathcal A)$ by \cref{pf:top-shape:calculus}(a). For
$\topshapeF\in\mathcal A$, the tied records of the selection and the conditions of $\sfname{Duty}$ show that
$\topshapeF$ belongs to the second-layer family class (eventually the extra data of the members are recursive,
and the linked families are aligned; \lref{Formalism.L2.Top3.layerTwoFamilyClass_of_linked_aligned} and \lref{Formalism.L2.Top3.linkedFamiliesAligned_of_duty}),
and second-layer families satisfy the group compatibility of the extra data
(\lref{Formalism.L2.group_compatibility_of_second_layer_family}); so the group compatibility is closed on
$\topshapeD(\mathcal A)$. The local lattice link \lref{Formalism.L2.G.linkG_concrete5} (\cref{thm:linkG}: if the
group compatibility and $\sfname{Adm2F}$ are closed on a class, so is $\sfname{LinkG}$) concludes.
\end{proof}

\subsubsection*{The links $\sfname{LinkEta}$ and $\sfname{SmallLevelRankEq}$}

\begin{lemma}\label{pf:top-shape:etaclosed}
Let $\topshapeF$ be a family such that eventually every member $D$ admits a family $M$ of local modules with a
normalized local compatibility and a Hecke tie. Then $\topshapeF$ satisfies $\sfname{LinkEta}$.
\formallabel{Formalism.L2.linkEta_closed}
\end{lemma}

\begin{proof}
$\sfname{LinkEta}$ is the eventual statement that the local places cover the twists: for $1\le j\le m$ and every place
$w$ with $\mu_w\ne\mu_0$, the local modules $\mathit{Mt}_i$, $i\in\mathrm{locs}(j,w)$, have total rank
$\dim\mathit{Mt}(\mu_w,j)$. At a member with $M$ as in the hypothesis this holds by
\lref{Formalism.L2.linkEtaTw_closed}: at each twisted place $w$ and lift index $j$ it compares the ranks of the
local modules over the places $i\in\mathrm{locs}(j,w)$ with the dimension of $\mathit{Mt}(\mu_w,j)$ by means of Hecke
standard modules (\lref{ReciprocityTheory.ChartRestriction.L3H.heckeStd} and \lref{Formalism.L2.locs_cover_twists_of_hecke}).
The threshold is that of the hypothesis.
\end{proof}

\begin{proposition}\label{pf:top-shape:eta}
Let $\mathcal A$ be a family class with a twist selection. Then every family of $\topshapeD(\mathcal A)$ satisfies
$\sfname{LinkEta}$; in particular $\sfname{LinkEta}$ is closed on $\topshapeD(\mathcal A)$.
\formallabel{GlobalTheory.Top.closed_linkEta_of_duty}
\end{proposition}

\begin{proof}
Let $\topshapeF\in\topshapeD(\mathcal A)$. The selection provides, for $x>x_1$ beyond its threshold and all $i$, a
tied record $T$ (\lref{Formalism.L2.Top3.tied_families_of_choice_alt}). Its normalized local data consist of a
family $M$ of local modules with a normalized local compatibility and a Hecke tie. By \cref{pf:top-shape:etaclosed},
$\sfname{LinkEta}$ holds.
Closedness follows from \cref{pf:top-shape:calculus}(a). (The formal statement also assumes that the families of
$\mathcal A$ satisfy $\sfname{Duty}$; this hypothesis is not used.)
\end{proof}

\begin{proposition}\label{pf:top-shape:small}
Let $\mathcal A$ be a family class with selections $\sigma,\sigma_S,\sigma_P,\sigma_{DH},\sigma_C,\sigma_B$ and the
hypotheses $h_{PA}$, $h_{\Lambda P}$, $h_{DivG}$ of \cref{pf:top-shape:shape}, and assume that $z_L$ is closed on
$\topshapeD(\mathcal A)$: for every $\topshapeF\in\topshapeD(\mathcal A)$ there is $x_r$ such that at selected members
with $x\ge x_r$, for every $k<m$, the level projection $l_{\rm Prj}(k,\mu_0)$ of the row-side family $\sigma_P$ is
surjective. Then $\sfname{SmallLevelRankEq}$ is closed on $\topshapeD(\mathcal A)$.
\formallabel{GlobalTheory.Top.link_alt_alt_alt}
\formallabel{Formalism.L2.N.RankZeroOfLevelThree.rank_zero_of_row_conditions_alt}
\end{proposition}

\begin{proof}
At selected members the hypothesis $h_{DivG}$, together with the fact that the indefinite realisation of
$\sigma_B$ uses all tuples, gives the division rows of the row-side family; $h_{PA}$ gives the row condition at
$\mu_0$ (\lref{Formalism.L2.Top4.selP_condition_zero_of_primAll}); $h_{\Lambda P}$ gives the
$\Lambda$-rows; and the tied data of $\sigma_P$ contain an index $t$ with tuple $\mu_0$. For $x$ beyond
$\max(x_{\rm sel},x_r)$ and $j=k+1\in\{1,\dots,m\}$, the surjectivity of $l_{\rm Prj}(k,t)$ and these row
conditions imply, by \lref{Formalism.L2.N.ExtNRows3.thetaT_injective_and_finrank_Zj_eq_D0_of_surjective},
that the theta map at level $j$ is injective and that the space $Z_j$ has dimension $D_0$; this dimension is
$d_{Z_0}(j)$. Hence $\sfname{SmallLevelRankEq}$ holds for every family of $\topshapeD(\mathcal A)$ for which the
property in $z_L$ holds, and closedness follows from \cref{pf:top-shape:calculus}(c),(e).
\end{proof}

\subsubsection*{Approximate orthogonality}

\begin{proposition}[Transport of Theorem~S]\label{pf:top-shape:thmS}
Let $\mathcal C$ be a family class with a spectral selection $s$, let $C_0\ge1$, and assume $\sfname{WuDecayOfGap}$,
$\sfname{RowsFinDim}$, $\sfname{TypeP}$, $\sfname{RowNGP}$, $\sfname{RowMainEqP}$, the loop rows for $s$, the local ties
and instance ties (with $q\ge1$) at all selected members, counting constants $c_{30}$ with $0\le C$ and
$2(C_{YL}+1)\le A$ for the constants $C,A$ of $c_{30}$, and the closedness on $\mathcal C$ of the trace frame, of the
eventual control of residues, of Theorem~S for the selection, and of the point count with constants $c_{30}$. Then
$\sfname{TuSize}\Rightarrow\sfname{LinkS}(C_0)$ is closed on $\mathcal C$.
\formallabel{Formalism.L2.SG.ThmSOwn.theoremS_transport_link_zero}
\end{proposition}

\begin{proof}
\emph{Step 1 (standing setup).} $\sfname{WuDecayOfGap}$ and $\sfname{RowsFinDim}$ give the finiteness of the relevant
row spaces (\lref{Formalism.L2.SG.Pf.finXC_of}); with the local ties and the closed trace frame, the selection
provides the standing setup at its members, as a closed property
(\lref{Formalism.L2.SG.ThmSOwn.selOK_of}). Consequently the concentration statement $\sfname{ConcEv}(C_0)$ is
closed, and so is the existence, eventually at every member, of a rows record with the rows tie, which is built from
the selection record, the instance tie and $q\ge1$ (\lref{Formalism.L2.SG.ThmSOwn.rowsTiePT'_of}); both by
\cref{pf:top-shape:calculus}(a).

\emph{Step 2 (four kernel estimates).} The closed concentration and the selection give the four hypotheses on the
kernels of the rows: the off-diagonal estimate for the kernels $K_{pp}$
(\lref{Formalism.L2.SG.ThmSOwn.hypOwnE_kppOff}); a diagonal positivity estimate for $K_{pp}$ derived from
$\sfname{TypeP}$ (\lref{Formalism.L2.SG.ThmSOwn.hypOwnE_kppDiagP}); the off-diagonal Hilbert--Schmidt estimate given
by Theorem~S for the selection (\lref{Formalism.L2.SG.ThmSOwn.hypOwnS_off}); and the diagonal positivity of the
loop rows derived from $\sfname{RowNGP}$ and $\sfname{RowMainEqP}$ (\lref{Formalism.L2.SG.ThmSOwn.hypOwnL_diagP}).
In each case the absolute convergence and the exponential decay of the tied kernels are part of the input.

\emph{Step 3 (conclusion).} With the counting constants, the closed residue control, the rows ties, the four
estimates and the closed point count, the concrete form of the $\sfname{LinkS}$ theorem
(\lref{Formalism.L2.SG.ThmSOwn.theoremS_link_concrete}) gives the closedness of
$\sfname{TuSize}\Rightarrow\sfname{LinkS}(C_0)$: at each member, either the window condition $C_0^{-1}\le
d_u\Gamma^2/D_*\le C_0$ fails for some word, and then the conclusion of $\sfname{LinkS}$ holds with $\eta_0=\eta_1=e_0=0$,
or the standing hypotheses give the three estimates with $\eta_0,\eta_1,\sqrt{e_0}\le\exp(-x^{\eps/8})$
(\lref{Formalism.MainNotions.S.linkS_of_standing_hypotheses}).
\end{proof}

\begin{corollary}[\cref{cor:families:linkS}]\label{pf:top-shape:linkS}
Let $\mathcal A$ be a family class with a twist selection $\sigma$ and a spectral selection $\sigma_S$, and assume
$\sfname{WuDecayOfGap}$, $\sfname{RowsFinDim}$, $\sfname{TypeP}$, $\sfname{RowNGP}$, $\sfname{RowMainEqP}$, the
loop rows $s_{Lp}$ for the induced spectral selection, $h_{InstR}$, counting constants $c_{30}$ with $h_{c30}$, and the
closedness on $\topshapeD(\mathcal A)$ of the trace frame $h_{TF}$, of Theorem~S for the selection $s_O$ and of the
point count $h_{30}$. Then $\sfname{TuSize}\Rightarrow\sfname{LinkS}(3)$ is closed on $\topshapeD(\mathcal A)$.
\formallabel{GlobalTheory.Top.link_alt_alt_bis}
\end{corollary}

\begin{proof}
The spectral selection induces a selection on $\topshapeD(\mathcal A)$ (threshold $x_{\rm sel}$ and spectral ties). At
selected members it provides local ties, and instance ties are obtained from $h_{InstR}$, from the real theta tie of
the tied record and from the scale tie of the family (\lref{Formalism.L2.SG.ThmSOwn.instTie_of_record}); the
instance has $q\ge1$. The eventual control of residues is closed on every class
(\lref{Formalism.L2.SG.ResidEvFree.residEv_closed_free}). \Cref{pf:top-shape:thmS}, applied with $C_0=3$ to
these ties, the closed trace frame, the counting constants, the closed residue control, $\sfname{TypeP}$,
$\sfname{RowNGP}$, $\sfname{RowMainEqP}$, the loop rows, and the closedness of Theorem~S for the selection and of the
point count, gives the claim.
\end{proof}

% ===========================================================================
\subsection{The division link}\label{sec:pf:top-shape:div}

For a family $\topshapeF$ in a class $\mathcal D$ suppose given \emph{kernel data} $R_k(\topshapeF)$: a threshold and,
for members $(x,i)$ beyond it with $x>x_1$, division kernel data for the extra data $\mathrm{ex}(x,i)$. For the
designed class of the package they are obtained from the indefinite realisations of the selection bundle and the
hypothesis $h_{PA}$. A \emph{division package} for kernel data, with constants $C,C_{\rm abs}$, is a record
containing an indefinite package of lattices, compatibility data, lattices $\mathit{Nu}_u$ attached to the words,
numerical records $R_u$ with ties $R_u.h-\mathrm{meanLog}(\kappa)=h(u)$ and $2F_o(\text{core})=F_o$, and the
inequalities $R_u.\mathrm{rhs}\le\mathrm{MAIN}(u)+C|TT|aL_a+C\,m(B_0+N_c)x^{C_{\rm abs}\eps}$.

\begin{definition}[Division kernel constant bounds]\label{pf:top-shape:dkcb}
Let $\topshapeF$ be a family with kernel data $R_k$ and $C_{\rm abs},C_H\in\R$. The \emph{division kernel constant
bounds} $\sfname{DKCB}(\topshapeF,R_k,C_{\rm abs},C_H)$ consist of:
\begin{enumerate}[label=(\roman*)]
\item (package) if $\sfname{DIVinN}$ holds, there is $C$ such that eventually the kernel data admit a division
package with constants $C,C_{\rm abs}$;
\item ($\mathrm{adm}_2$) eventually, for every word $u$, every place $w$, every $\varphi\in\mathit{NuLat}(u)$ and every
$F\in\widetilde M_{\mu_w}\cap V_0(\mu_w)$, the period $\mathrm{per}_u(\mu_w)(j_N\varphi,F)$ is the image of an
element $y\in\mathit{LL}$ with $y=0$ or $v_w(y)\ge\lfloor u_r(D,w,\sigma(u,w))\rfloor$;
\item (cut) eventually, for every used tuple $\mu$, every element of $V_0(\mu)$ that is an
$\OO_{\mathit{LL}}$-combination of the generators of $\widetilde M_\mu$ is an $\OO_{\mathit{LL}}$-combination of
elements of $\widetilde M_\mu\cap V_0(\mu)$;
\item ($\mathrm{Compl}_3$) there is $C$ such that eventually the comparison of the explicit heights
$h_{\rm expl}(u)+\bigl(\sum_ta_T(t)\bigr)\log T_u$ across words holds with constant $C$;
\item (frame number) if $\sfname{DIVinN}$ holds, then eventually $s_{\rm expl}+\mu_{\max}(\mathit{RR})\le x^{C_H}b$;
\item (second complement) there is $C$ such that eventually, for all $u$ and $j$,
\[
\Bigl|\sum_w\mathrm{sz}(w,E(u,j,w))-\Bigl(\sum_ta_T(t)\Bigr)\log(g_0g_1)\Bigr|\le CB_0\quad\text{and}\quad
\sum_{w}{}'\,|\mathrm{sz}(w,E(u,j,w))|\le CB_0,
\]
the primed sum running over the places whose residue characteristic is not a type prime;
\item (complement count) eventually the second half of the Bessel identity holds: there is an isometry $V$ of
$\mathit{EE}$ such that for every $u$, $B_u(\varphi,F)=a_u\,\mathrm{bes}_u(U_u\varphi,VF)$ for some $a_u\ne0$ and an
isometry $U_u$;
\item eventually the functionals $B_u$ are rational, $\mathit{RR}$ is finitely generated and spans, and $L_a\ge0$.
\end{enumerate}
The \emph{remaining division bounds} $\sfname{RDB}(\topshapeF,C_{\rm abs},C_H)$ are: (v) under the second
complement condition (vi), and the three conditions of (viii).
\formallabel{Formalism.L2.DIV.DivisionKernelConstantBounds}
\formallabel{Formalism.L2.Top4.RemainingDivisionBounds}
\end{definition}

\begin{proposition}[\cref{prop:families:div}, first part]\label{pf:top-shape:divmain}
Let $\mathcal D$ be a family class with kernel data $R_k(\topshapeF)$ for its families, and $C_{\rm abs}\in\R$. If
the property $\sfname{TuSize}\Rightarrow\sfname{DGamSize}\Rightarrow\sfname{DKCB}(\topshapeF,R_k(\topshapeF),
C_{\rm abs},C_H)$ is closed on $\mathcal D$, then so is $\sfname{TuSize}\Rightarrow\sfname{DGamSize}\Rightarrow
\sfname{LinkDIVt}(C_{\rm abs},C_H)$.
\formallabel{Formalism.L2.DIV.closed_linkDIVt_of_closed_bridge_kernel}
\formallabel{Formalism.L2.DIV.linkDIVt_of_bridge_kernel}
\formallabel{Formalism.L2.DIV.linkDIVt_of_rowsT}
\formallabel{Formalism.L2.DIV.linkDIVt_of_rowsF}
\formallabel{Formalism.L2.DIV.linkDIV_at_words}
\end{proposition}

\begin{proof}
We use the same $\eps_1$ and thresholds; it suffices to show that $\sfname{DKCB}(\topshapeF,R_k,C_{\rm abs},C_H)$
implies $\sfname{LinkDIVt}(C_{\rm abs},C_H)$ for a single family. The kernel data are converted into row-column
data, and the bounds into the corresponding bounds for these data, without change of content. Assume
$\sfname{DIVinN}$ (the hypothesis $\sfname{LinkEta}$ of $\sfname{LinkDIVt}$ is not used). Let $C_1$, $C_3$, $C_2$ be
the constants of (i), (iv), (vi), let $C=\max(C_1,C_3,C_2,0)$, and let $x$ be beyond all thresholds, beyond the
threshold of the kernel data and $x_1+1$. At a member $(D,L)$ the family ties give $h(u)=h_{\rm expl}(u)$ for all
words. We verify the nine clauses of $\sfname{LinkDIV}$ with constant $C$.
\begin{itemize}
\item Clauses (1), (6), (7): the division package of (i), together with (ii), (iii) and the defining tie of
$N^H_u$, gives for every word $u$ the \emph{word facts}: $B_u$ maps $N_u$ into $\mathit{CC}$, the height bound
$\mathrm{ht}(B_k(u))\le h(u)$, and the inequality (7) with constant $C_1$
(\lref{Formalism.L2.DIV.DivisionPackageColumns.wordFacts}); since $L_a\ge0$ by (viii) and the error terms are
non-negative, the inequality persists with $C\ge C_1$.
\item Clause (2) is (v), as $h=h_{\rm expl}$.
\item Clause (3) holds at every member with $x>x_1$ (\lref{Formalism.MainNotions.WeightDegreeSizes.zero_gain_above_family}).
\item Clause (4) is the rationality of (viii) together with (vii); clause (5) is (viii).
\item Clause (8) is the first inequality of (vi), the product of the type primes being $g_0g_1$; it persists with
$C\ge C_2$.
\item Clause (9) is (iv) with $C\ge C_3$, as $h=h_{\rm expl}$.\qedhere
\end{itemize}
\end{proof}

\begin{lemma}[Division packages from the selection]\label{pf:top-shape:wiresel}
Let $\mathcal A$ be a family class with a twist selection $\sigma$, a spectral condition, a base-level tied choice, a
definite height condition, a compatible column choice and a selection bundle for them, and assume $h_{Tw}$ (the twisted
records of the tied choice are tied to the size data), $h_{PA}$ (all indefinite realisations of the bundle are
primitive), $h_{DivG}$ and $h_{DcPtG}$ (the general forms of the division rows and of the point conditions for the tied
choices). Then for every family $\topshapeF$ of $\mathcal A$ and every $C_H$: under $0\le C_{\rm sz}$, $\sfname{DGamSize}$,
the rows and single-map conditions of the bundle and the rationality of the $B_u$, the condition $\sfname{DIVinN}$
implies that there is $C$ such that eventually the kernel data admit a division package with constants $(C,2)$.
\formallabel{GlobalTheory.Top.wire_sel_alt}
\end{lemma}

\begin{proof}
The coefficient fields of the members are totally complex. In the indefinite realisation of the selection bundle all
tuples are used (\lref{Formalism.IndefR.LevelRealisationRankEquality.tup_all_used_of_realisation}); hence
$h_{DivG}$ gives the division rows of the tied choice at every selected member
(\lref{Formalism.L2.DIV.DivRows3G.toRows}), and $h_{DcPtG}$ gives the point rows
(\lref{Formalism.L2.DIV.dcPt_of_G}). The realisation also identifies the base rank $D_0$ of the member with the
rank $d_Z$ of the level of the realisation (\lref{Formalism.IndefR.LevelRealisationRankEquality.indefinite_realisation_base_rank_eq}).
With $h_{Tw}$, $h_{PA}$ and the single-map condition, these are the inputs of the wiring theorem for division packages
(\lref{Formalism.L2.Top4Wire2.wire4F_alt}), applied with the constant $C_{\rm abs}=2$, which gives the claim.
\end{proof}

\begin{proposition}[\cref{prop:families:div}, second part]\label{pf:top-shape:bridge}
Let $\mathcal A$ carry the data (b)--(e) of \cref{pf:top-shape:shape} and let $R_k$ be the kernel data given by the
selection bundle and $h_{PA}$. Then $\sfname{TuSize}\Rightarrow\sfname{DGamSize}\Rightarrow\sfname{DKCB}(\topshapeF,
R_k(\topshapeF),2,C_H)$ is closed on $\topshapeD(\mathcal A)$. Consequently
$\sfname{TuSize}\Rightarrow\sfname{DGamSize}\Rightarrow\sfname{LinkDIVt}(2,C_H)$ is closed on $\topshapeD(\mathcal A)$.
\end{proposition}

\begin{proof}
By \cref{pf:top-shape:calculus}(a),(b) the following eight properties are jointly closed on $\topshapeD(\mathcal A)$:
\begin{enumerate}[label=(\alph*)]
\item $\sfname{TuSize}\Rightarrow\sfname{DGamSize}\Rightarrow\sfname{RDB}(\topshapeF,2,C_H)$, by
\cref{pf:top-shape:brest};
\item the cut condition (iii), which follows at every family from the $\OO_{\mathit{LL}}$-stability $h_{Mt}$ of the
lattices $\widetilde M_\mu$ (\lref{Formalism.L2.Top4.cut_of_stable});
\item the rows condition $r_{\rm Rows}$;
\item the single-map condition along the selection, a field of the bounded rows $h_{RB}$;
\item $0\le C_{\rm sz}$ (\cref{pf:top-shape:calculus}(f));
\item eventually $\mathrm{adm}_2$, by \cref{pf:top-shape:adm2};
\item the second complement condition, which follows from the closed inputs $a_{C2}$
(\lref{Formalism.L2.Top3.complement_two_of_items}: the second-complement inputs of a family imply the
second complement condition);
\item the complement count condition, by \lref{GlobalTheory.Top.selection_closed_of_wu_decay_alt}: under
$\sfname{WuDecayOfGap}$, $\sfname{RowsFinDim}$, $\sigma_J$, $j_{On}$ and $r_{\rm Rows}$ it is closed on
$\topshapeD(\mathcal A)$ (members of $\topshapeD(\mathcal A)$ eventually satisfy the gap hypothesis; the row spaces are
finite-dimensional; the multiplicity condition is part of $r_{\rm Rows}$; and the base vector image is full).
\end{enumerate}
Let $\topshapeF$ satisfy (a)--(h), $\sfname{TuSize}$ and $\sfname{DGamSize}$. The package condition (i) holds by
\cref{pf:top-shape:wiresel}. Condition (iv) holds for every family, with constant $20+6/\eps$
(\lref{Formalism.MainNotions.DIV.HookProducts.RBCompl3.rbCompl3}). Conditions (ii), (iii), (vi), (vii) are
(f), (b), (g), (h), and (v), (viii) come from (a) and (g). Thus $\sfname{DKCB}(\topshapeF,R_k,2,C_H)$ holds
(\lref{Formalism.L2.Top4.RemainingDivisionBounds.toRows_alt}). The last assertion follows from
\cref{pf:top-shape:divmain}.
\end{proof}

\begin{proposition}[Remaining division bounds]\label{pf:top-shape:brest}
Let $\mathcal A$ carry the selections, $h_{RB}$, $h_{PA}$, $h_{DivG}$, $h_{DcPtG}$, and the closed conditions
$r_{\rm Rows}$ and $\mathrm{tie}$ on $\topshapeD(\mathcal A)$. Then for every $C_{\rm abs}$ the property
$\sfname{TuSize}\Rightarrow\sfname{DGamSize}\Rightarrow\sfname{RDB}(\topshapeF,C_{\rm abs},C_H)$ is closed on
$\topshapeD(\mathcal A)$.
\formallabel{GlobalTheory.Top.link_brest_alt}
\formallabel{Formalism.L2.DIV.RestAll.bRest_all_alt}
\end{proposition}

\begin{proof}
The conditions $L_a\ge0$ hold eventually at every family (\lref{Formalism.L2.DIV.Wire.la_of}), and the
rationality of the $B_u$ and the fullness of $\mathit{RR}$ are closed on $\topshapeD(\mathcal A)$ under the stated
hypotheses (\lref{Formalism.L2.DIV.RatFullFam.buRat_rrFull_closed_alt}). For the frame-number bound, combine the closed properties $\mathrm{tie}$, $\sfname{LinkG}$ under
$\sfname{TuSize}\wedge\sfname{DGamSize}$ (\lref{Formalism.L2.DIV.LinkGCl.closed_linkG_of_tuSize_dGamSize}) and the
closed sharp-function threshold (\lref{Formalism.L2.DIV.FunSharpTop.closed_funSharp_thr_alt}), which provides a
function $C_{F0}$ of the outer constants such that the family is large for an explicit threshold built from $C_{F0}$
and, under $\sfname{DIVinN}$, eventually $\mu_{\max}(\mathit{RR})+\mathrm{mean}_e\log\kappa_R(e)\le x^{C_{F0}}b$. For a
family with these properties, $\sfname{TuSize}$, $\sfname{DGamSize}$, the second complement condition and
$\sfname{DIVinN}$: the ties and $\sfname{TuSize}$ give the level statement
(\lref{Formalism.L2.DIV.LvlWire.lvlF_of_ties}), hence, with the choice condition $\sfname{ChosenProviso}(3)$ of
the designed class, $\sfname{MuRank}$ and $\sfname{NuFullF}$
(\lref{Formalism.L2.DIV.LvlWire.muRank_of_lvl} and \lref{Formalism.L2.DIV.LvlWire.nuFullF_of_lvl}), hence
$\sfname{LinkMu}$ (\lref{Formalism.SlopeId.linkMu_of_full}). The column valuations given by the selection and
$\sfname{TuSize}$ yield the eventual bound $\mu_{\max}(M_u)\le x^{E_b}b$ for the level-zero lattices $M_u$ of the
words, with $E_b$ an explicit function of the outer constants (\cref{pf:top-shape:boundB}). With these inputs the frame-number bound
\cref{pf:top-shape:hh} gives the frame-number bound with constants $(C_{\rm sz},C_H)$; it combines the slope lemma on
the division side (\cref{pf:top-shape:slopelemma}, which uses $\sfname{LinkG}$, $\sfname{LinkMu}$ and the place
bounds of \cref{pf:top-shape:places}, resting on a Mertens-type estimate for weight sums,
\cref{pf:top-shape:tausum}) with the gauge-form bound of \cref{pf:top-shape:gauge}, derived from the second
complement condition.
\end{proof}

\subsubsection*{The frame-number bound}

We now prove the frame-number bound and its ingredients. Throughout, $\topshapeF$ is a family, $C_t$ is the
first of its outer constants and
\[
M_t=\max(C_t,5)+3 .
\]
``Eventually'' means: for all $x$ beyond a threshold with $x>x_1$ and all members $(D,L)=(D_{x,i},L_{x,i})$; the
scale of $L$ is then $x$. We write $b=e^\Lambda$ for the scale parameters of $L$, $a$, $b'$, $X'$, $X''$ for its other
scale parameters, $f,k,m,t$ for the ranks of $D$, $n$ for its degree, $N_m$ for its top rank, and put
\[
P=1+f+k+m+t+N_m+\Lambda,\qquad B_0=fmb\Lambda
\]
(the polynomial size and the loss bound). For a word $u$ let $N_u$ be its lattice in the division data, of rank
$d_u$, $M_u$ its level-zero lattice, and $\mathit{RR}$ the remainder lattice. Let $W$ be the support of $L$,
$\ell_w\ge0$ ($w\in W$) the weights of its places, and $\tau_w$, $\mathrm{gain}(j,w)$, $F(j,w)$, $E(u,j,w)$
($1\le j\le m$) its exponents; the gauge sums of $L$ are
\begin{gather*}
\mathrm{Ga}=2\sum_{w\in W}\ell_w\Bigl(\sum_j\mathrm{gain}(j,w)-\tau_w\Bigr),\\
\mathrm{Fo}=2\sum_j\sum_{w\in W}\ell_wF(j,w),\qquad
\mathrm{Ee}(u)=2\sum_j\sum_{w\in W}\ell_wE(u,j,w)
\end{gather*}
(\lref{Formalism.MainNotions.Lnk.Ga_eq}); in terms of the size function $\mathrm{sz}$ of $L$,
$\mathrm{Ee}(u)=\sum_{j=1}^m\sum_{w\in W}\mathrm{sz}(w,E(u,j,w))$. Let $h_u$ be the explicit height of $u$ (it
depends on $C_{\rm sz}$), $\kappa_R(e)$ the remainder constants and $\mathrm{mean}_e$ the mean over their index set.
The \emph{frame mean} and the \emph{frame-number bound} are
\[
s(N,h)=\frac{\sum_ud_u\bigl(\mu(N_u)-h_u\bigr)}{\sum_ud_u},\qquad
\sfname{FNB}(C_{\rm sz},C_H):\quad s(N,h)+\mu_{\max}(\mathit{RR})\le x^{C_H}b .
\]

\begin{lemma}[Polynomial sizes]\label{pf:top-shape:bpol}
Let $(D,L)$ satisfy the size conditions of the family at scale $x$.
\begin{enumerate}[label=(\alph*)]
\item If $C\in\R$, $e\in\N$ and $x\ge\max(9,C)$, then $C\,P^e\,b\le x^{eM_t+1}b$.
\item For every $C\in\R$, $C\,m\,B_0\le|C|\,P^4\,b$.
\end{enumerate}
\formallabel{Formalism.MainNotions.DIV.Proofs.bPol_le}
\formallabel{Formalism.MainNotions.DIV.Proofs.const_rank_base_le_poly_fourth}
\end{lemma}

\begin{proof}
(a) For $x\ge9$ the size conditions give $P\le x^{M_t}$
(\lref{Formalism.MainNotions.DIV.Proofs.szPol_le_rpow}: each of $f,k,m,t,N_m,\Lambda$ is bounded by an explicit
power of $x$). As $P\ge1$ and $b\ge1$, $C\,P^e\le x\,P^e\le x\cdot x^{eM_t}$, and multiplying by $b$ gives the claim.
(b) The numbers $m$, $f$, $\Lambda$ are summands of $P$, hence at most $P$, and $b\ge1$; so
$0\le m\,B_0=m\cdot f\cdot m\cdot b\cdot\Lambda\le P^4b$ and $C\,mB_0\le|C|\,mB_0\le|C|P^4b$.
\end{proof}

\begin{lemma}[Gauge-form bound]\label{pf:top-shape:gauge}
\begin{enumerate}[label=(\alph*)]
\item Let $C\in\R$ satisfy the second complement condition at $(D,L)$: for every word $u$ and $1\le j\le m$,
\begin{gather*}
\Bigl|\sum_{w\in W}\mathrm{sz}(w,E(u,j,w))-\Bigl(\sum_{t\in TT}a_T(t)\Bigr)\log\prod_{g\in G}g\Bigr|\le CB_0,\\
\sum_{w\in W,\ p_w\notin G}\bigl|\mathrm{sz}(w,E(u,j,w))\bigr|\le CB_0,
\end{gather*}
where $G$ is the set of type primes of $L$ and $p_w$ the residue characteristic of $w$. Then $-\mathrm{Ee}(u)\le
C\,m\,B_0$ for every $u$.
\item If the family satisfies the second complement condition (with one constant $C>0$ for all $x>1$ and all
members), then eventually $\mathrm{Ga}+\mathrm{Fo}-\mathrm{Ee}(u)\le x^{E_c}b$ for all $u$, where $E_c=9M_t+2$.
\end{enumerate}
\formallabel{Formalism.MainNotions.DIV.HookProducts.Row4C.negEe_le}
\formallabel{Formalism.MainNotions.DIV.HookProducts.Row4C.gauge_form_bound_of_uniform_size}
\end{lemma}

\begin{proof}
(a) The coefficients satisfy $a_T(t)\ge a>0$, and $\prod_{g\in G}g$ is a natural number,
so $\bigl(\sum_ta_T(t)\bigr)\log\prod_gg\ge0$ (with $\log0=0$). The first inequality gives
$\sum_w\mathrm{sz}(w,E(u,j,w))\ge-CB_0$ for each $j$, and summing over $1\le j\le m$ gives $\mathrm{Ee}(u)\ge-mCB_0$.

(b) Eventually $\mathrm{Ga}\le x^{4M_t+1}b$ and $\mathrm{Fo}\le x^{9M_t+1}b$
(\lref{Formalism.MainNotions.DIV.HookProducts.Row4.ga_pow} and \lref{Formalism.MainNotions.DIV.HookProducts.HookProductWeightBounds.fo_pow}).
Let $x$ also satisfy $x\ge\max(9,2+|C|)$ and put $Q=x^{9M_t+1}$. Then $\mathrm{Ga}\le Qb$ and $\mathrm{Fo}\le Qb$, since
$x\ge1$. By (a) and \cref{pf:top-shape:bpol}(b), $-\mathrm{Ee}(u)\le CmB_0\le|C|P^4b$, and $P^4\le x^{4M_t}\le Q$ by the
bound $P\le x^{M_t}$; so $-\mathrm{Ee}(u)\le|C|Qb$. Adding,
$\mathrm{Ga}+\mathrm{Fo}-\mathrm{Ee}(u)\le(2+|C|)Qb\le xQb=x^{9M_t+2}b$.
\end{proof}

The slope lemma rests on the following decomposition of the weighted slope excess; it is an identity and two
estimates for finite sums.

\begin{lemma}[Weighted slope excess]\label{pf:top-shape:wse}
Let $W_T\subseteq W$ and $J$ be finite sets, and let $\ell_w\ge0$, $\mathrm{gain}(j,w)$, $F(j,w)$, $E(j,w)$, $\tau_w$,
$\mathrm{io}_w$, $\mathrm{len}_w$, $\sigma'_w$ be real numbers ($w\in W$, $j\in J$). Put
$\sigma_w=\sum_{j\in J}(\mathrm{gain}(j,w)+F(j,w)-E(j,w))$, define $\mathrm{Ga},\mathrm{Fo},\mathrm{Ee}$ by the formulas above
(with $E(j,w)$ in place of $E(u,j,w)$), and put
\begin{gather*}
\mathrm{Sg}(s)=\sum_{w\in W}2\ell_w(s_w-\tau_w),\qquad\mathrm{Ta}=2\sum_{w\in W}\ell_w\tau_w,\\
\mathrm{IO}=\sum_{w\in W}2\ell_w\mathrm{io}_w,\qquad\mathrm{Be}=\sum_{w\in W}2\ell_w\frac{\mathrm{len}_w-d\tau_w}{d}
\end{gather*}
for a number $d>0$.
\begin{enumerate}[label=(\alph*)]
\item $\mathrm{Sg}(\sigma)=\mathrm{Ga}+\mathrm{Fo}-\mathrm{Ee}$.
\item (Upper bound.) If $\sigma'_w\le\sigma_w$ for $w\in W$, $\mu^{\max}_N\le\mu^{\max}_\mu+\mathrm{Sg}(\sigma')$ and
$\mu^{\max}_\mu\le\mu^{\max}_M$, then $\mu^{\max}_N\le\mu^{\max}_M+\mathrm{Ga}+\mathrm{Fo}-\mathrm{Ee}$.
\item (Lower bound.) Let $d_M>0$, $0\le\delta$, $f\ge0$ and $T>2$ with $\delta\le2f/T\le\frac12$ and $d=(1-\delta)d_M$.
Assume: $\tau_w=0$ for $w\in W_T$; $\tau_w\ge0$ and $\mathrm{io}_w\ge0$ for $w\in W$; $\mathrm{len}_w\le
d_M(\tau_w+\mathrm{io}_w)$ for $w\in W\setminus W_T$; $\sum_{w\in W_T}2\ell_w\mathrm{len}_w/d\le P_T\le L_T$;
$\sum_{w\in W}2\ell_w(\sigma_w-\sigma'_w)\le R_d\le R_X$; $\mathrm{Ta}\le b^2$;
$\frac{4f}Tb^2+f\frac{\log(T-1)}{T-2}\le1$; $\mathrm{IO}=2Q$; $\mu_\mu=\mu_M-f\frac{\log(T-1)}{T-2}$; and
$\mu_N=\mu_\mu+\mathrm{Sg}(\sigma')-\mathrm{Be}$. Then
\[
\mu_M+\mathrm{Ga}+\mathrm{Fo}-\mathrm{Ee}-(4Q+L_T+R_X+1)\le\mu_N .
\]
\end{enumerate}
\formallabel{Formalism.LFunctionChainParameters.UnitWeightPartition.weightedSlopeExcess_decomposition}
\formallabel{Formalism.LFunctionChainParameters.UnitWeightPartition.claimHigh_valid}
\formallabel{Formalism.LFunctionChainParameters.UnitWeightPartition.lower_bound_simplified}
\end{lemma}

\begin{proof}
(a) Expanding, $\mathrm{Sg}(\sigma)=\sum_w2\ell_w\bigl(\sum_j\mathrm{gain}(j,w)-\tau_w\bigr)+\sum_w2\ell_w\sum_jF(j,w)
-\sum_w2\ell_w\sum_jE(j,w)$, and exchanging the order of summation in the last two sums gives
$\mathrm{Ga}+\mathrm{Fo}-\mathrm{Ee}$.

(b) As $\ell_w\ge0$, $\mathrm{Sg}(\sigma')\le\mathrm{Sg}(\sigma)$, so
$\mu^{\max}_N\le\mu^{\max}_M+\mathrm{Sg}(\sigma)$, and (a) concludes.

(c) Split $\mathrm{Be}$ into the sums over $W_T$ and $W\setminus W_T$. On $W_T$, $\tau_w=0$ and the summand is
$2\ell_w\mathrm{len}_w/d$, so this part is at most $P_T$. For $w\in W\setminus W_T$, using $d_M/d=1/(1-\delta)$,
\[
2\ell_w\frac{\mathrm{len}_w-d\tau_w}{d}\le2\ell_w\frac{d_M(\tau_w+\mathrm{io}_w)-d\tau_w}{d}
=\frac1{1-\delta}\,2\ell_w\mathrm{io}_w+\frac\delta{1-\delta}\,2\ell_w\tau_w ,
\]
since the difference of the two sides is $2\ell_w(d_M(\tau_w+\mathrm{io}_w)-\mathrm{len}_w)/d\ge0$. Summing, and using
$\mathrm{io}_w\ge0$ on $W_T$ and $\tau_w=0$ on $W_T$ (so that the $\tau$-sum over $W\setminus W_T$ is $\mathrm{Ta}$),
\[
\mathrm{Be}\le P_T+\frac{2Q}{1-\delta}+\frac\delta{1-\delta}\mathrm{Ta}.
\]
Moreover $\mathrm{Sg}(\sigma)-\mathrm{Sg}(\sigma')=\sum_w2\ell_w(\sigma_w-\sigma'_w)\le R_d$. Hence, by (a),
\[
\mu_N\ge\mu_\mu+\mathrm{Ga}+\mathrm{Fo}-\mathrm{Ee}-R_d-P_T-\frac{2Q}{1-\delta}-\frac\delta{1-\delta}\mathrm{Ta}.
\]
Now $\delta\le\frac12$ gives $2Q/(1-\delta)\le4Q$ (as $Q=\frac12\mathrm{IO}\ge0$) and
$\delta/(1-\delta)\le2\delta\le4f/T$, so $\frac\delta{1-\delta}\mathrm{Ta}\le\frac{4f}Tb^2$ (as $\mathrm{Ta}\ge0$). With
$\mu_\mu=\mu_M-f\log(T-1)/(T-2)$, $R_d\le R_X$, $P_T\le L_T$ and the hypothesis
$\frac{4f}Tb^2+f\frac{\log(T-1)}{T-2}\le1$ we obtain the claim.
\end{proof}

\begin{lemma}[A Mertens-type bound]\label{pf:top-shape:mertensplaces}
For every real $y\ge1$, $\sum_{p<\lceil y\rceil}\frac{\log p}p\le2\log y$, the sum running over primes.
\formallabel{Formalism.L2.GlobalTheoryTheory.Places.mertens}
\end{lemma}

\begin{proof}
Write $\lceil y\rceil=n+1$, so that $n<y$ and the sum runs over the primes $p\le n$. If $y\ge4$, then $n\ge3$ and
$\sum_{p\le n}\frac{\log p}p\le\log n+\log4$ (\lref{Formalism.LogFactorialBounds.sum_primesLE_log_div_le}),
which is at most $2\log y$ since $n<y$ and $4\le y$. If $1\le y<4$, then $n\le3$: for $n\le1$ the sum is empty;
for $n=2$ we have $y>2$ and the sum is $\frac{\log2}2\le2\log y$; for $n=3$ we have $y>3$ and the sum is
$\frac{\log2}2+\frac{\log3}3\le2\log3\le2\log y$.
\end{proof}

\begin{lemma}[Weight sums]\label{pf:top-shape:tausum}
For every family there are $C$, a threshold and a function $\delta(x)\to0$ ($x\to\infty$) such that eventually
\[
\sum_{w\in W}\mathrm{sz}(w,\tau_w)\le C\,m\,|TT|\,a\log X'+b^{\delta(x)} .
\]
\formallabel{Formalism.MainNotions.G.TopChoiceBounds.tau_sum_family_of_mertens}
\end{lemma}

\begin{proof}
The exponent sizes of the members of a family satisfy the bounds of
\lref{Formalism.MainNotions.G.TopChoiceBounds.family_exponent_sizes_hold}, and every member satisfies the
Mertens-type bound of \cref{pf:top-shape:mertensplaces}. From these two inputs,
\lref{Formalism.MainNotions.G.TopChoiceBounds.stmtFTa_proved} (the weight-sum statement for families whose
members satisfy the exponent-size bounds and the Mertens-type bound) gives $C$, the threshold and $\delta$ as
stated.
\end{proof}

\begin{proposition}[Place bounds]\label{pf:top-shape:places}
Let the family satisfy the level statement $\sfname{LevelT}$ (there are $C_{26}$ and a threshold such that eventually
the members satisfy the level admissibility $\mathrm{adm}_{26}(L,C_{26})$), the commensurability $\sfname{NuCommensF}$,
$\sfname{TuSize}$, $\sfname{MuRank}$, $\sfname{LinkG}$ and $\sfname{LinkMu}$. Let $I_L$ be the inertia sum of $L$, so
that $\mathrm{IO}=2I_L$ for the local inertia orders $\mathrm{io}_w$ of $L$. Then there is $C\ge4$ such that eventually
$I_L\ge0$ and, for every word $u$, there are real numbers $\sigma'_w$ ($w\in W$) and
$\delta,\mu_\mu,\mu^{\max}_\mu,P_T,R_d,L_T,R_X$ satisfying the hypotheses of \cref{pf:top-shape:wse}(b) and (c) with
$J=\{1,\dots,m\}$, $W_T$ the set of places of $W$ of residue characteristic $T_u$, $\mathrm{len}_w$ the local lengths
of $N_u$ relative to the ramified lattice, $d_M=d^S_u$ the rank of the full local module of $u$, $d=d_u=\rank N_u$,
$f$ the rank $f$ of $D$, $T=T_u$, $\mu_M=\mu(M_u)$, $\mu_N=\mu(N_u)$, $\mu^{\max}_M=\mu_{\max}(M_u)$,
$\mu^{\max}_N=\mu_{\max}(N_u)$, $Q=I_L$, and
\[
L_T\le C\,|TT|\,m\Bigl(\frac ak+b'\Bigr)\log T_u,\qquad R_X\le C\,X''(n).
\]
\formallabel{Formalism.L2.GlobalTheoryTheory.Places.placesF_min'}
\end{proposition}

\begin{proof}
Take $C=\max(C_{26},C_{\mu},4)$, where $C_\mu$ is the constant of $\sfname{LinkMu}$. Choose the threshold beyond the
thresholds of all hypotheses and of the following eventual facts, and let $x$ be beyond it.
\begin{itemize}
\item \emph{Sizes.} $x^9\le e^{x^\eps}=b$ (\lref{Formalism.MainNotions.Lnk.ev_pow_le_exp}); with the size conditions
this gives the auxiliary size bounds of \lref{Formalism.L2.GlobalTheoryTheory.Places.sizes_at}, and
$m\ge x^{2/5}\ge6$. By $\sfname{TuSize}$, $T_u\ge b^3\ge b\ge X'$ for every $u$, and then
$T_u>2$, $2f/T_u\le\frac12$ and $\frac{4f}{T_u}b^2+f\frac{\log(T_u-1)}{T_u-2}\le1$
(\lref{Formalism.L2.GlobalTheoryTheory.Places.size_alt}).
\item \emph{Weight sums.} By \cref{pf:top-shape:tausum}, with $\delta(x)\le1$ and $x\ge4|C_T|+1$ ($C_T$ the constant
there), $\sum_{w}\mathrm{sz}(w,\tau_w)\le b^2$ (\lref{Formalism.L2.GlobalTheoryTheory.Places.ta_le}).
\item \emph{Exponents.} Each $\tau_w$ is a natural number; $\mathrm{io}_w\ge0$ and $I_L\ge0$
(\lref{Formalism.MainNotions.Lnk.ioP_nonneg} and \lref{Formalism.MainNotions.Lnk.IOsum_nonneg}); $\tau_w=0$ if
the residue characteristic of $w$ is $T_u$, because $T_u\ge X'$
(\lref{Formalism.L2.GlobalTheoryTheory.Places.tau_eq_zero_of_le}).
\item \emph{Local lattices.} By $\sfname{LinkG}$, at every place $w$ whose residue characteristic is not $T_u$, the
lattice of $u$ is locally contained in the ramified lattice up to the exponent $\lfloor\mathrm{ur}(D,w,\tau_w)\rfloor$,
and $\mathrm{len}_w\le d^S_u(\tau_w+\mathrm{io}_w)$. Outside the support $W$, either the residue characteristic is $T_u$
and the local length vanishes (\lref{Formalism.L2.GlobalTheoryTheory.Places.lenLoc_atT_zero}), or the residue
characteristic exceeds $X''(n)$ and the length vanishes by the lattice conditions of the family.
\end{itemize}
With these facts, the level admissibility $\mathrm{adm}_{26}(L,C_{26})$, the definitions of $N'$ and $N_u$, the
commensurability, $\rank N_u>0$ and the bound of $\sfname{LinkMu}$, the construction of
\lref{Formalism.L2.GlobalTheoryTheory.Places.places_at_of} produces the required numbers: $\sigma'_w$ is the local
slope contribution of $w$, $\delta$ is the relative rank defect $1-d_u/d^S_u\le2f/T_u$, the bound
$\sum_{W_T}2\ell_w\mathrm{len}_w/d_u\le L_T$ comes from the lengths at the places above $T_u$, and $R_X\le CX''(n)$
from $\sfname{LinkMu}$.
\end{proof}

\begin{corollary}[Slope lemma, division side]\label{pf:top-shape:slopelemma}
Under the hypotheses of \cref{pf:top-shape:places} (with the level statement and commensurability derived from the
level statement $\sfname{LvlF}$ and $\sfname{TuSize}$), there are $C\ge0$ and a threshold such that eventually, for
every word $u$,
\begin{gather*}
\mu_{\max}(N_u)\le\mu_{\max}(M_u)+\mathrm{Ga}+\mathrm{Fo}-\mathrm{Ee}(u)+C\,m\,B_0,\\
\mu(M_u)+\mathrm{Ga}+\mathrm{Fo}-\mathrm{Ee}(u)-\mathrm{Lo}_u(C)\le\mu(N_u),
\end{gather*}
where $\mathrm{Lo}_u(C)=C\,I_L+C\,|TT|\,m\bigl(\frac ak+b'\bigr)\log T_u+C\,X''(n)+1$.
\formallabel{Formalism.L2.DIV.LvlWire.slope_lemma_division_side}
\formallabel{Formalism.MainNotions.Lnk.slopesN_at_of}
\formallabel{Formalism.MainNotions.Lnk.slopesNF_of}
\end{corollary}

\begin{proof}
$\sfname{LvlF}$ and $\sfname{TuSize}$ give the level statement $\sfname{LevelT}$ and the commensurability
(\lref{Formalism.MainNotions.J.lemma_five_one_of_theorem_two} and \lref{Formalism.MainNotions.nuCommensF_min}). Let
$C\ge4$ be as in \cref{pf:top-shape:places}. By \cref{pf:top-shape:wse}(b), $\mu_{\max}(N_u)\le\mu_{\max}(M_u)+
\mathrm{Ga}+\mathrm{Fo}-\mathrm{Ee}(u)$, and $CmB_0\ge0$. By \cref{pf:top-shape:wse}(c),
$\mu(N_u)\ge\mu(M_u)+\mathrm{Ga}+\mathrm{Fo}-\mathrm{Ee}(u)-(4I_L+L_T+R_X+1)$, and
$4I_L+L_T+R_X+1\le\mathrm{Lo}_u(C)$ since $I_L\ge0$ and $C\ge4$.
\end{proof}

\begin{theorem}[Frame-number bound]\label{pf:top-shape:hh}
Let $\topshapeF$ be a family, $C_H\in\R$, and $E_b$, $C_{F0}$ functions of the outer constants such that
\[
C_H\ge\max\bigl(E_b,E_c,C_{F0},E_d\bigr)+1,\qquad E_c=9M_t+2,\quad E_d=9M_t+1,
\]
at the outer constants of $\topshapeF$ (this is the largeness of $\topshapeF$ for the corresponding threshold).
Assume $\sfname{LvlF}$, $\sfname{TuSize}$, $\sfname{MuRank}$, $\sfname{LinkG}$, $\sfname{LinkMu}$ and the second
complement condition, and that eventually
\[
\text{(i)}\ \ \mu_{\max}(M_u)\le x^{E_b}b\ \text{for all }u,\qquad
\text{(ii)}\ \ \mu_{\max}(\mathit{RR})+\mathrm{mean}_e\log\kappa_R(e)\le x^{C_{F0}}b .
\]
Then eventually $\sfname{FNB}(C_{\rm sz},C_H)$ holds.
\formallabel{Formalism.L2.DIV.LvlWire.hh_of_lvl}
\formallabel{Formalism.MainNotions.DIV.Proofs.height_bound_of_component_bounds}
\formallabel{Formalism.MainNotions.DIV.Proofs.hhRest_of_parts}
\formallabel{Formalism.MainNotions.DIV.Proofs.height_bound_of_slope_fullness}
\formallabel{Formalism.MainNotions.DIV.Proofs.height_bound_of_word_slopes}
\end{theorem}

\begin{proof}
\emph{Step 1 (inputs).} \Cref{pf:top-shape:slopelemma} gives $C$ with the two slope inequalities; $\sfname{LvlF}$
gives that eventually every $N_u$ is finitely generated and spans its space
(\lref{Formalism.L2.DIV.LvlWire.nuFullF_of_lvl}); \cref{pf:top-shape:gauge}(b) gives eventually
$\mathrm{Ga}+\mathrm{Fo}-\mathrm{Ee}(u)\le x^{E_c}b$; and $\sfname{TuSize}$ gives eventually
$-(h_u+\mathrm{mean}_e\log\kappa_R(e))\le x^{E_d}b$
(\lref{Formalism.MainNotions.DIV.Proofs.hExpl6_lower_pow}).

\emph{Step 2 (the rest bound).} We show: for every $C'$, eventually for all $u$,
\[
\mu_{\max}(M_u)+\mathrm{Ga}+\mathrm{Fo}-\mathrm{Ee}(u)+C'mB_0-h_u+\mu_{\max}(\mathit{RR})\le x^{C_H}b .
\]
Let $E=\max(E_b,E_c,C_{F0},E_d)$ and take $x$ beyond the thresholds of Step~1, (i), (ii), and $x\ge\max(9,|C'|)$.
Since $x\ge1$, $x^{a}b\le x^Eb$ for $a\le E$. The five quantities
\begin{gather*}
\mu_{\max}(M_u),\qquad\mathrm{Ga}+\mathrm{Fo}-\mathrm{Ee}(u),\qquad\mu_{\max}(\mathit{RR})+\mathrm{mean}_e\log\kappa_R(e),\\
-\bigl(h_u+\mathrm{mean}_e\log\kappa_R(e)\bigr),\qquad C'mB_0
\end{gather*}
are each at most $x^Eb$: the first four by (i), Step~1, (ii) and Step~1, and the last because
$C'mB_0\le|C'|P^4b\le x^{4M_t+1}b$ by \cref{pf:top-shape:bpol} and $4M_t+1\le E_d\le E$. Their sum is the left side
above, so it is at most $5x^Eb\le x\cdot x^Eb=x^{E+1}b\le x^{C_H}b$, as $x\ge9$ and $E+1\le C_H$.

\emph{Step 3 (word slopes).} Eventually $\rank N_u>0$ and $N_u$ is full, so $N_u\ne0$ and
$\mu(N_u)\le\mu_{\max}(N_u)$ (\lref{Formalism.MainNotions.DIV.Proofs.ScaledLatticeNorms.mu_le_muMax}). With the
first slope inequality and Step~2 (for $C'=C$),
\[
\mu(N_u)-h_u+\mu_{\max}(\mathit{RR})\le\mu_{\max}(M_u)+\mathrm{Ga}+\mathrm{Fo}-\mathrm{Ee}(u)+CmB_0-h_u+\mu_{\max}(\mathit{RR})
\le x^{C_H}b .
\]

\emph{Step 4 (averaging).} Multiply by $d_u\ge0$ and sum over $u$:
$\sum_ud_u(\mu(N_u)-h_u)\le\bigl(\sum_ud_u\bigr)\bigl(x^{C_H}b-\mu_{\max}(\mathit{RR})\bigr)$. As $\sum_ud_u\ge d_{u_0}>0$,
dividing gives $s(N,h)\le x^{C_H}b-\mu_{\max}(\mathit{RR})$, which is $\sfname{FNB}(C_{\rm sz},C_H)$.
\end{proof}

\begin{proposition}[Maximal slopes at level zero]\label{pf:top-shape:boundB}
Let $\topshapeF$ carry twist-selection functions and satisfy $\sfname{TuSize}$. Then eventually
$\mu_{\max}(M_u)\le x^{4M_t+1}b$ for all words $u$.
\formallabel{Formalism.L2.DIV.SlopeMaxB.eventually_muMax_le_pow_ebMu_mul_of_tuSize}
\formallabel{Formalism.L2.DIV.SlopeMaxB.boundB_of_values}
\formallabel{Formalism.L2.DIV.SlopeMaxB.boundB_of_slopes}
\end{proposition}

\begin{proof}
The selection gives the column ties of the family
(\lref{Formalism.L2.DIV.SlopeMaxB.colTies_of_sel_alt}), hence, by \cref{pf:top-shape:tcv}, eventually every word
has column valuations. With $\sfname{TuSize}$, the size condition of the family on the divisor ranks, and the two
Arakelov slope inequalities (\cref{pf:top-shape:lowslope,pf:top-shape:upslope}),
\lref{Formalism.MainNotions.J.slopes_family_of_column_values} gives $C$ such that eventually
\[
\mu_{\max}(M_u)\le C\Bigl(|TT|\,a\log(kv)+f(b'+S_c)\bigl(k\log(kv)+\log T_u\bigr)+fk^2\Lambda\Bigr)
\]
for all $u$, where $v$ is the half root discriminant of the host and $S_c$ a scale parameter. Each of the three terms
is bounded by a fixed multiple of $P^4b$ (using $\log T_u\le C_{Tu}m\Lambda$ from $\sfname{TuSize}$), which gives
$\mu_{\max}(M_u)\le|C|(12+2|C_{Tu}|)P^4b$ (\lref{Formalism.L2.DIV.SlopeMaxB.muMax_le_pol}). For
$x\ge\max(9,|C|(12+2|C_{Tu}|))$, \cref{pf:top-shape:bpol}(a) with $e=4$ gives $\mu_{\max}(M_u)\le x^{4M_t+1}b$.
\end{proof}

\begin{proposition}[Admissible twisted periods]\label{pf:top-shape:adm2}
Let $\mathcal A$ carry a twist selection $\sigma$, the closed conditions $a_{Pt0}$, $a_{PtTw}$, $a_{Per0}$, $a_{PerTw}$,
$a_{P10}$, $a_{AtG}$ on $\topshapeD(\mathcal A)$, and the service hypothesis $h_{\rm Serv}$ for the induced selection.
Then ``eventually $\mathrm{adm}_2(D_{x,i},L_{x,i})$'' is closed on $\topshapeD(\mathcal A)$.
\formallabel{GlobalTheory.Top.closed_family_admissible_twist_link_alt}
\end{proposition}

\begin{proof}
$\sfname{DGamSize}$ holds for every family of $\topshapeD(\mathcal A)$, so it is closed. By
\cref{pf:top-shape:servT}, a served selection on a class on which $\sfname{DGamSize}$ is closed yields the closed
service property $\sfname{ServF}$ (for some $C\ge0$, eventually the members satisfy the local
``Weil--Jacquet--Shalika law'' $\sfname{Serv}(C)$ of \cref{pf:top-shape:serv}); its proof is the sum-over-primes
service theorem \cref{pf:top-shape:serv}, which rests on the formal-group statements of \cref{pf:top-shape:packet}
and the orbit-entry bounds of \cref{pf:top-shape:orbit}. The level primes of designed
members are eventually coprime to the stabilisers
(\lref{Formalism.MainNotions.closed_eventually_levelPrimeCoprimeStabiliser}). These closed properties and the
six hypotheses form the record of inputs of \cref{pf:top-shape:adminc}, which derives $\mathrm{adm}_2$ (and a
companion condition $\mathrm{adm}_{12}$) at every member satisfying them.
\end{proof}

\subsubsection*{Service and orbit-entry bounds}

We now prove the service property used above, at the level of its main steps. Its local ingredients are $p$-adic
estimates for the theta kernels of a member, organised by orbits of places. We recall the objects involved.

\emph{Orbit entries.} For a datum $D$ an \emph{orbit-entry record} $E$ consists of Young diagrams
$\nu_{u,i,z}$ (for words $u$, levels $i<m$ and places $z$ of the auxiliary place set of $D$) whose row lengths are
the rows of the weights of $u$ at level $i+1$, theta kernels $\Theta_{u,\mu,w,i}(L,L')$ (power series over $\C_p$,
$p$ the residue characteristic of $w$, in variables indexed by places and matrix entries, for classes $L,L'$ at levels
$i,i+1$ of the twisted chain of $u,\mu$), theta constants $\varepsilon_w$, $\sigma_{w,z}$ and power series $\beta_{w,z}$ over
$\C_p$ in the variables $\Z/n_p$. For a real $B$,
$\sfname{EntGe}(u,\mu,w,i,B)$ means that every kernel $\Theta_{u,\mu,w,i}(L,L')$ has kernel valuation at least $B$
with respect to the diagrams $\nu_{u,i,\cdot}$ ($\sfname{KvalGe}$). For a prime $p$ the places form orbits under
$z\mapsto$ the $s$-th translate ($s\in\Z/n_p$), and $\sfname{OrbGe}(r,\varepsilon,\sigma,\beta,B)$ means that every kernel
with an orbit theta expansion of type $(\varepsilon,\sigma,\beta)$ on the orbit of $r$ has kernel valuation at least
$B$ with respect to the diagrams along the orbit.

\emph{Constants.} $\sfname{CLarge}(C_6,C_d,K,c_e)$ means $C_d\ge0$, $C_6\le c_6$, $3C_d\le1+\delta_1$, $\delta_1\ge0$,
$\delta_2\ge0$ and $3C_d\le c_{\rm dig}$, for the constants $\delta_1,\delta_2$ of $K$ and $c_6,c_{\rm dig}$ of the error
constants $c_e$. The \emph{orbit letter bounds} $\sfname{OLB}(C_6,C_d)$ are the two local theorems on letters: for
every admissible local datum, Theorem~A with constant $C_d$ and its Corollary~2 with constants $(C_6,C_d)$, and for every
admissible one-place local datum, Theorem~One with constant $C_d$. The \emph{shape} of $E$ is the list of explicit
inequalities between the sizes $|\nu_{u,i,z}|$, their first rows, the weights $w(i+1,z)>0$ of the places and the scale
parameters: $0\le|\nu_{u,i,z}|-(k-1)w(i+1,z)\le4w(i+1,z)$ at non-top places, the analogous inequalities with the first
row removed and $(k-3)w$ in place of $(k-1)w$ at top places together with $2w(i+1,z)\le$ the first row, $|\nu_{u,i,z}|\le
a^2$, $2\log_5(a^2)+2\le k-1$, $X''(n)\le a$, and bounds for the capacities and the first two rows.

\begin{theorem}[Orbit-entry bounds]\label{pf:top-shape:orbit}
Let $D$ be a datum with size data $L$, instance $X$, constants $K$, divisor data $D_x$, an orbit-entry record $E$ and
error constants $c_e$, and let $C_6,C\in\R$. Assume $\sfname{OLB}(C_6,C)$, that the theta kernels at the places over
the ramified class of $D$ have the shape $\exp(-S\sum_tu_t\sum_{a,b}X_{(0,t),a,b}X_{(1,t),a,b})\cdot\sigma'$ with $|S|\le
p^{1/2}$, $|u_t|\le1$ and integral divided-power coefficients of $\sigma'$, that $E$ has the shape above, and
$\sfname{CLarge}(C_6,C,K,c_e)$. Then the following eight statements hold.
\begin{enumerate}[label=(O\arabic*)]
\item\label{O:add} (Additivity over orbits.) For every prime $p$ with admissible orbits, every word $u$, level $i<m$,
set $R$ of orbit representatives, admissible $(\varepsilon,\sigma,\beta)$ and numbers $\beta_r$ ($r\in R$): if every
orbit of $r\in R$ satisfies $\sfname{OrbGe}$ with bound $\beta_r$, then every kernel with a theta expansion over all
places of type $(R,\varepsilon,\sigma,\beta)$ has kernel valuation at least $\sum_{r\in R}\beta_r$.
\item (Additivity at conjugate primes.) The analogue of \ref{O:add} at primes $p\ge5$ with $n_p=1$, for kernels with a
real-place theta expansion, with the one-place bounds.
\item (Restriction.) The restriction of a kernel with a theta expansion over all places of type $(R,\varepsilon,\sigma,
\beta)$ to the orbit of $r\in R$ has an orbit theta expansion of type $(\varepsilon,\sigma_r,\beta_r)$.
\item (Restriction at conjugate primes.) The restriction to a place of a kernel with a real-place theta expansion at
a prime $p\ge5$ with $n_p=1$ is an admissible one-place local datum.
\item\label{O:val} (Orbit valuation bound.) For every prime $p\le X''(n)$ with admissible orbits which is of residue
class zero with $h_p\ne0$ or of signal type, every $u$, $i<m$, every orbit representative $r$ of kind B or C and every
admissible $(\varepsilon,\sigma,\beta)$, $\sfname{OrbGe}$ holds on the orbit of $r$ with the bound
$g_0(r)+F_{\rm Orb}(i+1,p,r)-\mathrm{err}_F(i+1,p,r)$, where $g_0(r)$ is the positive part of the base gain and
\begin{align*}
\mathrm{err}_F(j,p,r)={}&\gamma_p\,\mathrm{tops}(r)\,(b+\eta+1+3S_c)+c_6\frac{\sum_sw(j,r_s)}p
+n_pN_{j-1}N_j\bigl(\log_p(4n_p\,\mathrm{Nm}\,a)+1\bigr)\\
&+\delta_1n_p\,\mathrm{Nm}^2\log_p(2a)+\delta_2\,\mathrm{tops}(r)\log_p(2a)^2,
\end{align*}
with $\gamma_p=p/(p^2-1)$, $\mathrm{tops}(r)$ the number of top places in the orbit of $r$, $r_s$ its translates and
$\eta,S_c,\mathrm{Nm}$ size parameters.
\item (Discriminant places.) At every level-free place $w$ over the ramified class of $D$,
$\sfname{EntGe}(u,\mu,w,i,-\frac12\sum_t\mathrm{ct}_D(u,i+1,t))$ for all $u,\mu$ and $i<m$.
\item (Valuation sums.) For every prime $5\le p\le X''(n)$ of residue class zero with $h_p=0$, every $u$ and $i<m$,
\[
F_{\rm inert}(i+1,p)-\mathrm{Eat}_C(i+1,p)\le\rho_p\sum_z|\nu_{u,i,z}|-C\,fN_iN_{i+1}\bigl(\log_p(a^2)+1\bigr).
\]
\item (Scaled generators.) For every prime $5\le p\le X''(n)$ of type $C_s$, every $u$ and $i<m$: every kernel of theta
product form has kernel valuation at least $\mathrm{low}_{C_s}(i+1,p)$, and, if the error constants are large,
$F_{\rm inert}(i+1,p)-\mathrm{Eat}_I(i+1,p)\le\mathrm{low}_{C_s}(i+1,p)$.
\end{enumerate}
\formallabel{Formalism.MainNotions.DIV.Lemmas.OrbitEntryBounds.orbit_entry_bounds_hold}
\formallabel{Formalism.MainNotions.DIV.Lemmas.OrbitEntryBounds.orbit_entry_bounds_of_family}
\end{theorem}

\begin{proof}
The eight statements are proved below for an arbitrary orbit-entry record: (O1) is \cref{pf:top-shape:o1}, (O2) is
\cref{pf:top-shape:o2}, (O3) and (O4) are \cref{pf:top-shape:o34}, (O5) is \cref{pf:top-shape:o5} (which uses
$\sfname{OLB}(C_6,C)$, the shape of $E$ and $\sfname{CLarge}$), (O6) is \cref{pf:top-shape:o6} (which uses the shape of
the kernels at the ramified places), (O7) is \cref{pf:top-shape:o7}, and (O8) is \cref{pf:top-shape:o8}.
\end{proof}

\subsubsection*{The local orbit bounds}

We prove the eight local statements. Fix a prime $p\ge5$; $\C_p$ carries the absolute value with $|p|=p^{-1}$. Let
$T$ be a finite set of places, $N_0<N_1$, and consider the variables $Z_{z,a,b}$ and $X_{z,i,a}$ ($z\in T$, $a<N_0$,
$b<N_1$, $i\in\N$). For an increasing map $I\colon\{0,\dots,k-1\}\to\{0,\dots,N_1-1\}$ let
\[
M_z(I)=\det\Bigl(\sum_{a<N_0}X_{z,i,a}Z_{z,a,I(c)}\Bigr)_{0\le i,c<k}\in\Z[X,Z].
\]
For Young diagrams $\nu=(\nu_z)_{z\in T}$ a \emph{column choice} $I$ assigns to every $z$ and every column $c$ of
$\nu_z$ an increasing map $I_{z,c}$ of length the length of the column, and $G_\nu(I)=\prod_z\prod_cM_z(I_{z,c})$. For a
power series $\Theta$ in the $Z$-variables and a polynomial $P$ put
\[
\Theta(\partial)P=\sum_\mu\mathrm{coeff}_\mu(P)\,\mu_Z!\,\mathrm{coeff}_{\mu_Z}(\Theta)\,X^{\mu_X},
\]
where $\mu=(\mu_X,\mu_Z)$ runs over the exponents of $P$ and $\mu_Z!=\prod_v\mu_Z(v)!$; this is $\Theta(\partial/\partial Z)P$
evaluated at $Z=0$. We say that $\Theta$ \emph{has kernel valuation at least $B$ on $\nu$} if every coefficient of every
$\Theta(\partial)G_\nu(I)$ has absolute value at most $p^{-B}$ (this is the predicate $\sfname{KvalGe}$ used above).
Along an orbit $z\mapsto z_s$ ($s\in\Z/n_p$) of places the \emph{jet logarithm} of the variable $v=(z,a,b)$ is
\[
\lambda(v)=\sum_{e\ge0}p^{-e}\,Z_{(z_{-2e},a,b)}^{\,p^{2e}} .
\]

\begin{lemma}[Kernel valuations]\label{pf:top-shape:kvalbasic}
\begin{enumerate}[label=(\alph*)]
\item If $|\kappa!\,\mathrm{coeff}_\kappa(\Theta)|\le p^{-B}$ for every exponent $\kappa$ with $\sum_{a,b}\kappa(z,a,b)=|\nu_z|$
for all $z$, then $\Theta$ has kernel valuation at least $B$ on $\nu$.
\item If $T$ is the disjoint union of blocks $T_r$, $\Theta=\prod_r\Theta_r$ with $\Theta_r$ a series in the $Z$-variables of
$T_r$, and $\Theta_r$ has kernel valuation at least $\beta_r$ on $\nu|_{T_r}$, then $\Theta$ has kernel valuation at least
$\sum_r\beta_r$ on $\nu$.
\item If $\Theta$ has kernel valuation at least $B$ and $B'\le B$, it has kernel valuation at least $B'$.
\end{enumerate}
\formallabel{UpperTheory.kvalGe_of_norm_mfact_mul_coeff_le}
\end{lemma}

\begin{proof}
(a) A minor $M_z(I)$ of size $k$ is homogeneous of degree $k$ in the variables $Z_{z,\cdot,\cdot}$ and does not involve
the $Z$-variables of other places, so $G_\nu(I)$ is homogeneous of degree $\sum_c(\text{length of column }c)=|\nu_z|$ in
the $Z$-variables of $z$. Hence every exponent $\mu$ of $G_\nu(I)$ has $\mu_Z$ of the stated type, every coefficient of
$\Theta(\partial)G_\nu(I)$ is a $\Z$-linear combination of numbers $\mu_Z!\,\mathrm{coeff}_{\mu_Z}(\Theta)$ with such $\mu_Z$,
and the ultrametric inequality gives the claim. (b) The column choice and $G_\nu(I)=\prod_rG_{\nu|T_r}(I|_{T_r})$ split
along the blocks, with disjoint sets of $X$- and $Z$-variables; a monomial in the $Z$-variables splits uniquely into its
block parts and $\mu_Z!$ is the product of the block factorials, so $\Theta(\partial)G_\nu(I)=\prod_r\Theta_r(\partial)
G_{\nu|T_r}(I|_{T_r})$ (\lref{PI.OrbJet.OJ.applyAt0_prod}). A coefficient of a product of polynomials in disjoint
variables is the product of coefficients, so its absolute value is at most $\prod_rp^{-\beta_r}$. (c) is clear.
\end{proof}

\begin{lemma}[Triangular substitution]\label{pf:top-shape:tri}
Let $|\varepsilon|=1$. A power series $\Phi$ has no monomials of total degree $<n$ if and only if $\Phi(\varepsilon\lambda)$
has none; and every power series $\Psi$ is $\Phi(\varepsilon\lambda)$ for a unique $\Phi$.
\end{lemma}

\begin{proof}
$\lambda(v)=Z_v+(\text{terms of degree }p^{2e}\ge25)$, so the substitution $Z_v\mapsto\varepsilon\lambda(v)$ maps a monomial
of degree $d$ to $\varepsilon^d$ times itself plus terms of degree $>d$. Hence the lowest-degree part of
$\Phi(\varepsilon\lambda)$ is $\varepsilon^{\deg}$ times that of $\Phi$, which gives the first claim, and $\Phi$ is found from
$\Psi$ degree by degree (\lref{PI.OrbShared.Tri.vanish_of_subst} and \lref{PI.OrbShared.Tri.subst_sol}); this is the formal
inverse function theorem.
\end{proof}

\begin{proposition}[(O1) Additivity over orbits]\label{pf:top-shape:o1}
Statement (O1) of \cref{pf:top-shape:orbit} holds for every orbit-entry record.
\formallabel{PI.OrbAddUp}
\formallabel{PI.OrbJet.AU.orbAddUp_main}
\formallabel{PI.OrbJet.AU.subst_Pm}
\end{proposition}

\begin{proof}
Let $R$ be a set of orbit representatives, $(\varepsilon,\sigma,\beta)$ admissible, $\beta_r$ the orbit bounds, and let
$\Theta(\varepsilon\lambda)=U\cdot E_{\rm all}(\lambda)$ with $U$ integral, where $E_{\rm all}=\exp\bigl(\sum_{r\in R}\sigma_r\sum_{a,b}
\beta_r(Z_{(r_\cdot,a,b)})\bigr)$. The series $\beta_r$ have no constant term (condition B3: every exponent has non-zero
even and odd parts), and $\varepsilon\ne0$.

\emph{Orbit factors.} For $r\in R$ and an exponent $k$ on the orbit of $r$ let $F_r(k)$ be the series with
$F_r(k)(\varepsilon\lambda)=Z^k\cdot E_r(\lambda)$, $E_r=\exp(\sigma_r\sum_{a,b}\beta_r)$ (\cref{pf:top-shape:tri}). Since $Z^k$ has
integral coefficients, $F_r(k)$ has an orbit theta expansion of type $(\varepsilon,\sigma_r,\beta_r)$, so by hypothesis it has
kernel valuation at least $\beta_r$ on the diagrams of the orbit. For an exponent $m$ on all places let
$P_m=\prod_{r\in R}F_r(m|_r)$ (each factor transported to the variables of the orbit of $r$). Since the orbits partition
the places and the substitution is multiplicative, $P_m(\varepsilon\lambda)=Z^m\cdot E_{\rm all}(\lambda)$, and
by \cref{pf:top-shape:kvalbasic}(b) $P_m$ has kernel valuation at least $\sum_r\beta_r$.

\emph{Truncation.} Fix a column choice $I$ and let $n$ exceed the $Z$-degree of $G=G_\nu(I)$. Put
$\Theta'=\sum_{|m|<n}\mathrm{coeff}_m(U)P_m$. Then $(\Theta-\Theta')(\varepsilon\lambda)=(U-U_{<n})E_{\rm all}(\lambda)$ has no
monomials of degree $<n$, so neither has $\Theta-\Theta'$ (\cref{pf:top-shape:tri}), and $\Theta(\partial)G=\Theta'(\partial)G$
since only coefficients of degree $<n$ enter.

\emph{Estimate.} Every coefficient of $\Theta'(\partial)G=\sum_m\mathrm{coeff}_m(U)\,P_m(\partial)G$ is a finite sum of
products of $\mathrm{coeff}_m(U)$, of absolute value $\le1$, and coefficients of $P_m(\partial)G$, of absolute value
$\le p^{-\sum_r\beta_r}$; the ultrametric inequality gives the bound $p^{-\sum_r\beta_r}$.
\end{proof}

\begin{proposition}[(O2) Additivity at conjugate primes]\label{pf:top-shape:o2}
Statement (O2) of \cref{pf:top-shape:orbit} holds for every orbit-entry record.
\formallabel{PI.OrbAddOne}
\formallabel{PI.OrbJet.AO.orbitsAddAtConjugate}
\formallabel{PI.OrbJet.AO.bound_place}
\end{proposition}

\begin{proof}
Here $n_p=1$, so every place is its own orbit. Write the real-place expansion as $\Theta(\varepsilon_0\lambda)=U\cdot
L(\lambda)$ with $U$ integral and $L$ the exponential of the sum of the one-place series, and let $\psi$ be the inverse
substitution of $Z\mapsto\varepsilon_0\lambda$, so that $\Theta=\psi^*(U\cdot L)$. Split $U=\sum_{k\in S}\mathrm{coeff}_k(U)Z^k+V$
with $S$ finite and $V$ of order larger than the total degree of $G_\nu(I)$; since $\psi$ has no constant terms,
$\psi^*(V\cdot L)$ has the same property and does not contribute to $\Theta(\partial)G_\nu(I)$
(\lref{PI.OrbJet.AO.order_psi_gt}). For each $k$, $\psi^*(Z^kL)$ is the product over the places $z$ of one-place
series $T_z(k_z)$ (\lref{PI.OrbJet.AO.factor}), each of which is an admissible one-place datum
(\lref{PI.OrbJet.AO.admissible_jet_term}); by hypothesis $T_z(k_z)$ has kernel valuation at least $\beta_z$, and
by \cref{pf:top-shape:kvalbasic}(b) the product has kernel valuation at least $\sum_z\beta_z$. As in the proof of
\cref{pf:top-shape:o1}, the ultrametric inequality with $|\mathrm{coeff}_k(U)|\le1$ concludes.
\end{proof}

\begin{proposition}[(O3), (O4) Restriction]\label{pf:top-shape:o34}
Statements (O3) and (O4) of \cref{pf:top-shape:orbit} hold for every orbit-entry record.
\formallabel{PI.OrbRestr}
\formallabel{PI.OrbJet.OrbW.orbitRestriction}
\formallabel{PI.OrbRestrOne}
\formallabel{PI.OrbJet.restricted_single_jet_equation}
\end{proposition}

\begin{proof}
Let $\rho$ be the ring homomorphism of power series which sets to zero the variables of the places outside the orbit of
$r$ (for (O4): outside the place $z$). The jet logarithm of a variable involves only variables of the same orbit, so
$\rho$ commutes with the substitution $Z\mapsto\varepsilon\lambda$; applying $\rho$ to $\Theta(\varepsilon\lambda)=U\cdot
\exp(\sum_{r'}\sigma_{r'}\sum_{a,b}\beta_{r'})(\lambda)$ gives
\[
\rho(\Theta)(\varepsilon\lambda)=\rho(U)\cdot\exp\Bigl(\sigma_r\sum_{a,b}\beta_r\Bigr)(\lambda),
\]
because the summands of the other representatives involve only variables outside the orbit and the $\beta_{r'}$ have no
constant term (\lref{PI.OrbJet.OrbW.restricted_jet_equation}). As $\rho(U)$ is integral, the restriction has an
orbit theta expansion of type $(\varepsilon,\sigma_r,\beta_r)$; this is (O3). The same computation at a prime with $n_p=1$
and one place $z$ gives the one-place equation, and with the bounds on $\varepsilon_z$, $\sigma_z$, $\beta_z$ the restricted
datum is admissible; this is (O4).
\end{proof}

\begin{proposition}[(O5) Orbit valuation bound]\label{pf:top-shape:o5}
If $\sfname{OLB}(C_6,C)$ holds, $E$ has the shape and $\sfname{CLarge}(C_6,C,K,c_e)$ holds, then statement (O5) of
\cref{pf:top-shape:orbit} holds.
\formallabel{PI.OrbBoundF}
\formallabel{PI.OrbBoundFW.prvF_of_orbitLetterBounds}
\formallabel{PI.OrbBoundFW.orbit_place_term_le_sum}
\formallabel{PI.OrbBoundFW.sum_tops_on_orbit}
\formallabel{PI.OrbBoundFW.summable_truncated_geometric}
\formallabel{PI.OrbBoundFW.sum_Wc_le}
\formallabel{PI.OrbBoundFW.loss_le}
\formallabel{PI.OrbBoundFW.setting_sq}
\formallabel{PI.OrbBoundFW.eOrb_bounds}
\formallabel{PI.OrbBoundFW.dl_off}
\formallabel{PI.OrbBoundFW.dl_on}
\end{proposition}

\begin{proof}
Fix $p$, $u$, $i<m$ and a representative $r$ of kind B or C, put $j=i+1$, $n=n_p$, $w_s=w(j,r_s)$ and
\[
E_s=|\nu_{u,i,r_s}|-(k-1)w_s\ \ (r_s\text{ not top}),\qquad E_s=|\nu_{u,i,r_s}|-\nu^{(0)}_{u,i,r_s}-(k-3)w_s\ \ (r_s\text{ top}),
\]
where $\nu^{(0)}$ is the first row. Let $\mathcal L$ be the orbit local datum of $r$ (prime $p$, $d=n$, ranks $N_i<N_{i+1}$,
diagrams $\nu_{u,i,r_s}$, top set $\mathrm{TT}=\{s:r_s\text{ top}\}$).

\emph{Step 1 (Corollary 2 applies).} The datum is admissible (it is built from the admissible $(\varepsilon,\sigma,\beta)$,
\lref{PI.OrbBoundFW.admissible_of_jet_bounds}). The Hida setting holds with $a^2$: $p\le X''(n)\le a\le a^2$;
$|\nu_{u,i,z}|\le a^2$; and $2\log_p(a^2)+2\le2\log_5(a^2)+2\le k-1\le N_i-1$ by the shape (rank condition) and
$N_i\ge k$. Moreover $k\ge m+3\ge3$, $w_s>0$, $0\le E_s\le4w_s$ by the shape, and $\mathrm{Dl}(s)=(k-1)w_s+E_s$ off
$\mathrm{TT}$, $\mathrm{Dl}(s)=(k-1)w_s-2w_s+E_s$ on $\mathrm{TT}$ (as $\mathrm{Dl}(s)$ is the size, minus the first row at top
places). So Corollary~2 of $\sfname{OLB}(C_6,C)$ gives kernel valuation at least
\[
B_0=\sum_{s\in\mathrm{TT}}\mathrm{base}(s)+\mathrm{lp}_F\bigl((k-1)w\bigr)-C_6\frac{\sum_sw_s}p-\gamma_p\sum_{s\in\mathrm{TT}}W_c(s)
-C\,nN_iN_{i+1}\bigl(\log_p(a^2)+1\bigr),
\]
with $\mathrm{base}(s)=\sum_{e\ \rm odd,\ e\ne d_{\rm top}(s)}\max\bigl((\nu^{(0)}_s-\nu^{(1)}_s)/p^e-\mathrm{cap}(s+e),0\bigr)$ and
$W_c(s)=\sup\{\mathrm{cap}(s-e):e\text{ odd},\ p^e\le|\nu_{s-e}|\}$. By \cref{pf:top-shape:kvalbasic}(c) it suffices to show
$g_0(r)+F_{\rm Orb}(j,p,r)-\mathrm{err}_F(j,p,r)\le B_0$, where $F_{\rm Orb}(j,p,r)=\mathrm{lp}_F((k-1)w)$ by definition.

\emph{Step 2 (base gain).} If $r$ is of the second kind, $g_0(r)=0\le\sum_{\rm TT}\mathrm{base}(s)$ as each $\mathrm{base}(s)\ge0$.
Otherwise $r=(0,t_0)$, the prime is of signal type, and
$g_0(r)=\max(0,\sum_{t\in\mathrm{tops}(t_0)}\sum_{e\ \rm odd,\ \neg\mathrm{Excl}}\max(0,\mathrm{al}(t)/p^e-\mathrm{cap}_U(t,e))-\mathrm{dl})$ with
$\mathrm{dl}=\delta_1n\,\mathrm{Nm}^2\log_p(2a)+\delta_2N_v\log_p(2a)^2\ge0$. Drop $\mathrm{dl}$. The tops of $t_0$ correspond
bijectively to the top positions $s$ of the orbit via $s\mapsto(r_s)_2$ (injectivity because the positions are distinct
modulo the period, surjectivity because every top in the base orbit of $t_0$ is reached at an even
position). For a top position $s$ with place $(0,t)$ and odd $e$, the translate $(0,t)_e$ is of the second kind at a signal
prime, hence no odd $e$ is excluded and $d_{\rm top}(s)=\infty$. Termwise,
$\mathrm{al}(t)/p^e-\mathrm{cap}_U(t,e)\le(\nu^{(0)}_s-\nu^{(1)}_s)/p^e-\mathrm{cap}(s+e)$, by the shape conditions
$\mathrm{al}(t)\le\nu^{(0)}-\nu^{(1)}$ and $\mathrm{cap}(z)\le\mathrm{fill}_G(z)+c_G=\mathrm{cap}_U$. The series are dominated by
$|x|p^{-e}$ and converge. Hence $g_0(r)\le\sum_{\rm TT}\mathrm{base}(s)$.

\emph{Step 3 (capacities).} Every capacity is at most $b+\eta+1+3S_c$ by the shape, hence so is every $W_c(s)$ (a supremum
of capacities), and $\sum_{\rm TT}W_c(s)\le\mathrm{tops}(r)(b+\eta+1+3S_c)$; multiplying by $\gamma_p\ge0$ gives the first
term of $\mathrm{err}_F$.

\emph{Step 4 (loss).} Put $P=nN_iN_{i+1}$ and $\ell=\log_pa\ge1$ (as $p\le a$). Then $\log_p(a^2)=2\ell$ and, as $C\ge0$,
$C(2\ell+1)\le3C\ell\le(1+\delta_1)\ell$ by $\sfname{CLarge}$. Since $N_i,N_{i+1}\le\mathrm{Nm}$, $P\le n\,\mathrm{Nm}^2$, and
$\ell\le\log_p(2a)$ and $\ell\le\log_p(4n\,\mathrm{Nm}\,a)$ (as $4n\,\mathrm{Nm}\ge8$):
\[
CP(\log_p(a^2)+1)\le P\ell+\delta_1P\ell\le P\bigl(\log_p(4n\,\mathrm{Nm}\,a)+1\bigr)+\delta_1n\,\mathrm{Nm}^2\log_p(2a)
+\delta_2\,\mathrm{tops}(r)\log_p(2a)^2 .
\]

\emph{Step 5.} $C_6\le c_6$ and $\sum_sw_s\ge0$ give $C_6\sum_sw_s/p\le c_6\sum_sw_s/p$. Adding Steps 2--5,
$g_0+F_{\rm Orb}-\mathrm{err}_F\le B_0$.
\end{proof}

The ramified places are handled by a counting argument for Laplace operators on products of minors. For a place pair
$t$ (the places $(0,t)$, $(1,t)$) put $\Delta_t=\sum_{a,b}\partial_{Z_{(0,t),a,b}}\partial_{Z_{(1,t),a,b}}$ and, for a row $r$ and
$b<N_1$, $d_{r,b}=\sum_aX_{(0,t),r,a}\,\partial_{Z_{(1,t),a,b}}$; for Young diagrams $\alpha$, $\beta$ put
$c(\alpha,\beta)=(|\alpha|-\alpha^{(0)})+\min(\alpha^{(0)},\beta^{(0)})$. A polynomial $g$ has \emph{row-$0$ degree} $\le d$ if
any $d+1$ operators $d_{0,b}$ annihilate it.

\begin{lemma}[Laplace operators on minors]\label{pf:top-shape:lap}
\begin{enumerate}[label=(\alph*)]
\item Every minor $M_{(1,t)}(I)$ has row-$0$ degree $\le1$, minors at other places have row-$0$ degree $0$, and row-$0$
degrees add under products.
\item If $g$ does not involve the variables $Z_{(0,t),\cdot,\cdot}$ and has row-$0$ degree $\le d$, then
$\Delta_t^n\bigl(\prod_cM_{(0,t)}(I_c)\cdot g\bigr)=0$ for $n>(|\alpha|-\alpha^{(0)})+\min(\alpha^{(0)},d)$, the product running
over the columns $c$ of $\alpha$.
\item $\Delta_t^nG_\nu(I)=0$ if $n>c(\nu_{(0,t)},\nu_{(1,t)})$.
\end{enumerate}
\formallabel{PI.EntAtD.W.pMinorAff}
\formallabel{PI.EntAtD.W.pSDegMul}
\formallabel{PI.EntAtD.W.pLapForms}
\formallabel{PI.EntAtD.W.lapZero_of}
\end{lemma}

\begin{proof}
(a) $d_{0,b}$ replaces the entries $\sum_aX_{(1,t),i,a}Z_{(1,t),a,b}$ of the column $b$ of the matrix of $M_{(1,t)}(I)$ by
$\sum_aX_{(1,t),i,a}X_{(0,t),0,a}$, which does not depend on $Z$; the determinant is linear in each column, so applying
$d_{0,b}$ and then $d_{0,b'}$ gives $0$ (for $b=b'$ the second derivative of a linear function; for $b\ne b'$ a
determinant with two equal columns). The operators $d_{0,b}$ only differentiate $Z_{(1,t)}$-variables, and they are
derivations, so row-$0$ degrees add by the Leibniz rule.

(b) Expand each minor of size $k$ by the Leibniz formula: it is a sum of products of $k$ entries
$\ell_{q,b}=\sum_aX_{(0,t),q,a}Z_{(0,t),a,b}$, exactly one from row $q=0$ and $k-1$ from rows $q\ge1$. So the product over the
columns of $\alpha$ is a sum of products containing $\alpha^{(0)}$ entries of row $0$ and $|\alpha|-\alpha^{(0)}$ entries of
rows $\ge1$ (the column lengths add up to $|\alpha|$). Each $\ell_{q,b}$ is linear in the $Z_{(0,t)}$-variables and $g$ does
not involve them, so in $\Delta_t^n$ applied to such a product times $g$ every $\partial_{Z_{(0,t),a,b}}$ must consume one
entry $\ell_{q,b}$, replacing it by $X_{(0,t),q,a}$, while the paired $\partial_{Z_{(1,t),a,b}}$ acts on $g$; summed over
$a$ this applies $d_{q,b}$ to $g$. Each entry can be consumed once, so at most $|\alpha|-\alpha^{(0)}$ steps consume
entries of rows $\ge1$ and at most $\alpha^{(0)}$ consume row-$0$ entries; the latter apply operators $d_{0,b}$ to $g$, which
commute with the $d_{q,b}$ and with each other and do not raise the row-$0$ degree, so at most $d$ of them act
non-trivially. Hence every term vanishes when $n$ exceeds $(|\alpha|-\alpha^{(0)})+\min(\alpha^{(0)},d)$
(\lref{PI.EntAtD.W.lap_tm_zero}, \lref{PI.EntAtD.W.minor_mem}, and \lref{PI.EntAtD.W.ms_prod}).

(c) Write $G_\nu(I)=\prod_cM_{(0,t)}(I_c)\cdot g$, where $g$ collects the minors at the other places; $g$ does not involve
$Z_{(0,t)}$ and by (a) has row-$0$ degree at most the number $\nu^{(0)}_{(1,t)}$ of minors at $(1,t)$. Apply (b) with
$d=\nu^{(0)}_{(1,t)}$.
\end{proof}

\begin{proposition}[(O6) Discriminant places]\label{pf:top-shape:o6}
Let $\Theta=\exp\bigl(-S\sum_tu_t\sum_{a,b}Z_{(0,t),a,b}Z_{(1,t),a,b}\bigr)\cdot\sigma'$ with $|S|\le p^{1/2}$, $|u_t|\le1$ and
$|m!\,\mathrm{coeff}_m(\sigma')|\le1$ for all $m$. Then $\Theta$ has kernel valuation at least
$-\frac12\sum_tc(\nu_{(0,t)},\nu_{(1,t)})$ on $\nu$. In particular statement (O6) of \cref{pf:top-shape:orbit} holds.
\formallabel{jet_bound_ramified_prime}
\formallabel{PI.EntAtD.W.entAtD_of}
\formallabel{PI.EntAtD.W.entAtD_of_core}
\formallabel{PI.EntAtD.W.core_of}
\formallabel{PI.EntAtD.W.pExpCoeff}
\formallabel{PI.EntAtD.W.expand_of}
\formallabel{PI.EntAtD.W.pBridge}
\formallabel{PI.EntAtD.W.pLapComm}
\formallabel{PI.EntAtD.W.vanish_of}
\end{proposition}

\begin{proof}
\emph{Step 1 (the exponential).} With $Y_{t,a,b}=Z_{(0,t),a,b}Z_{(1,t),a,b}$ and $c_t=-Su_t$,
$\exp(\sum c_tY_{t,a,b})=\prod_{t,a,b}\sum_ec_t^eY_{t,a,b}^e/e!$. Hence $\kappa!\,\mathrm{coeff}_\kappa$ of the exponential is
$\prod_{t,a,b}e_{t,a,b}!\,c_t^{e_{t,a,b}}$ if $\kappa$ is balanced ($\kappa((0,t),a,b)=\kappa((1,t),a,b)=e_{t,a,b}$), since then
$\kappa!=\prod(e!)^2$, and $0$ otherwise.

\emph{Step 2 (expansion).} For a polynomial $P$ and an $X$-exponent $m$, writing $\mu_Z=\kappa_1+\kappa_2$ and
$\mu_Z!=\binom{\mu_Z}{\kappa_2}\kappa_1!\kappa_2!$, Step~1 gives
\[
\mathrm{coeff}_m\Theta(\partial)P=\sum_{\kappa_2}\sum_n\bigl(\kappa_2!\,\mathrm{coeff}_{\kappa_2}\sigma'\bigr)\prod_tc_t^{n_t}\,
A(P,m,\kappa_2,n),
\]
with integers $A(P,m,\kappa_2,n)=\sum_\mu\mathrm{coeff}_\mu(P)\binom{\mu_Z}{\kappa_2}\prod_{t,a,b}e_{t,a,b}!$, the sum over the $\mu$ with
$\mu_X=m$, $\kappa_2\le\mu_Z$ and $\mu_Z-\kappa_2$ balanced with $\sum_{a,b}e_{t,a,b}=n_t$.

\emph{Step 3 (bridge).} Expanding $\Delta_t^{n_t}=\sum_{\sum e=n_t}\frac{n_t!}{\prod e!}\prod(\partial_{Z_{(0,t),a,b}}
\partial_{Z_{(1,t),a,b}})^{e_{a,b}}$ and using $\partial^eZ^\mu=\frac{\mu!}{(\mu-e)!}Z^{\mu-e}$, one finds
$\bigl(\prod_tn_t!\bigr)A(P,m,\kappa_2,n)=\mathrm{coeff}_{(m,\kappa_2)}\bigl(\prod_t\Delta_t^{n_t}P\bigr)$; indeed
$\binom{\mu_Z}{\kappa_2}\prod e!=\mu_Z!/(\kappa_2!\prod e!)$ because $(\mu_Z-\kappa_2)!=\prod(e!)^2$.

\emph{Step 4 (vanishing).} The $\Delta_t$ commute. If $n_t>c(\nu_{(0,t)},\nu_{(1,t)})$ for some $t$, apply
$\Delta_t^{n_t}$ first: by \cref{pf:top-shape:lap}(c) it kills $G_\nu(I)$, so $A(G_\nu(I),m,\kappa_2,n)=0$.

\emph{Step 5 (estimate).} In the expansion of $\mathrm{coeff}_m\Theta(\partial)G_\nu(I)$ only terms with
$n_t\le c(\nu_{(0,t)},\nu_{(1,t)})$ survive, and each has absolute value at most
$1\cdot\prod_t(p^{1/2})^{n_t}\cdot1\le p^{\frac12\sum_tc(\nu_{(0,t)},\nu_{(1,t)})}$, as $|A|\le1$. The ultrametric inequality
gives the first claim.

\emph{Step 6 (comparison with $\mathrm{ct}_D$).} For $\alpha=\nu_{(0,t)}$, $\beta=\nu_{(1,t)}$, with rows the weight rows of
$u$: if $t$ is not top, $\mathrm{ct}_D(u,i+1,t)=|\alpha|\ge c(\alpha,\beta)$; if $t$ is top,
$\mathrm{ct}_D=|\alpha|-(\alpha^{(0)}-\alpha^{(1)})+\beta^{(0)}\ge|\alpha|-\alpha^{(0)}+\beta^{(0)}\ge c(\alpha,\beta)$. By
\cref{pf:top-shape:kvalbasic}(c) the kernel valuation is at least $-\frac12\sum_t\mathrm{ct}_D(u,i+1,t)$.
\end{proof}

\begin{proposition}[(O7) Valuation sums]\label{pf:top-shape:o7}
If $E$ has the shape and $\sfname{CLarge}(C_6,C,K,c_e)$ holds, then statement (O7) of \cref{pf:top-shape:orbit} holds.
\formallabel{PI.SumOrbC}
\formallabel{SumOrbCW.sumOrbC}
\formallabel{SumOrbCW.size_ge}
\formallabel{SumOrbCW.NN_le}
\end{proposition}

\begin{proof}
Let $5\le p\le X''(n)$ be of residue class zero with $h_p=0$, $j=i+1$, $S=\sum_\tau w_j(\tau)=\sum_zw(j,z)\ge0$ and
$\ell=\log_pa\ge1$. By definition $F_{\rm inert}(j,p)=(k-1)S/(p^2-1)$ and
$\mathrm{Eat}_C(j,p)=c_{\rm dig}f\,\mathrm{Nm}^2\ell+S/(p^2-1)$, and $\rho_p=1/(p^2-1)$. The shape gives
$(k-1)w(j,z)\le|\nu_{u,i,z}|$ for every $z$ (at a top place: $|\nu|\ge\nu^{(0)}+(k-3)w\ge(k-1)w$, as $\nu^{(0)}\ge2w$), so
$(k-1)S/(p^2-1)\le\rho_p\sum_z|\nu_{u,i,z}|$. Moreover $C(2\ell+1)\le3C\ell\le c_{\rm dig}\ell$ and $N_iN_{i+1}\le\mathrm{Nm}^2$,
so $C\,fN_iN_{i+1}(\log_p(a^2)+1)\le c_{\rm dig}f\,\mathrm{Nm}^2\ell$. Since $S/(p^2-1)\ge0$,
\[
F_{\rm inert}-\mathrm{Eat}_C\le\rho_p\sum_z|\nu_{u,i,z}|-c_{\rm dig}f\,\mathrm{Nm}^2\ell
\le\rho_p\sum_z|\nu_{u,i,z}|-C\,fN_iN_{i+1}\bigl(\log_p(a^2)+1\bigr).
\qedhere
\]
\end{proof}

The scaled-generator bound rests on a linear program along an orbit. For $d\ge2$, $l\ge3$ and demands $D_s\ge0$
($s\in\Z/d$), a \emph{flow} is $x_{s,e}\ge0$ ($e<d$ odd) with $\sum_el^ex_{s,e}+\sum_ex_{s-e,e}\le D_s$ for all $s$, of value
$\sum_{s,e}x_{s,e}$; $\mathrm{lp}_F(D)=\gamma_l\sum_sD_s-\max(\text{flow value})$ with $\gamma_l=l/(l^2-1)$. A \emph{dual} is
$\sigma$ with $0\le\sigma_s\le1$ and $l^e\sigma_s+\sigma_{s+e}\ge1$ for odd $e<d$. By linear programming duality
(\lref{Formalism.IncomingBounds.FlowLoadBounds.strong} and \lref{Formalism.IncomingBounds.FlowLoadBounds.exists_opt_flow};
\cite{Schrijver1986}) the maximum is attained and there is an optimal dual with $\mathrm{lp}_F(D)=\sum_sD_s(\gamma_l-\sigma_s)$;
by complementary slackness, at every $s$ with $\sigma_s\ge2/l$ the demand $D_s$ equals the incoming flow
$\sum_ex_{s-e,e}$ of an optimal flow (\lref{Formalism.IncomingBounds.FlowDualityReindexing.fd_part1}).

\begin{lemma}[Perturbed duals]\label{pf:top-shape:lp}
\begin{enumerate}[label=(\alph*)]
\item For an optimal dual $\sigma$ of demands $D\ge0$, $\sum_{s:\sigma_s\ge2/l}D_s\le\frac1l\sum_sD_s$.
\item Let $l\ge5$, $d\ge2$, $\kappa>0$, $w_s\ge0$ and $\kappa w_s\le D_s\le\kappa w_s+4w_s$. Then there is $\sigma\ge0$ with
$l^e\sigma_s+\sigma_{s+e}\ge1$ for all $s$ and all odd $e$, and
$\mathrm{lp}_F(\kappa w)-\kappa\sum_sw_s/l^3-8\sum_sw_s/l\le\sum_sD_s(\gamma_l-\sigma_s)$.
\end{enumerate}
\formallabel{UpperTheory.sum_demand_filter_two_div_le_of_isOptDual}
\formallabel{UpperTheory.exists_nonneg_dual_fnum_sub_le_of_demand_near_mul}
\formallabel{UpperTheory.exists_maxFlow_attained_and_isOptDual}
\end{lemma}

\begin{proof}
(a) Let $x$ be an optimal flow of value $V$. Summing the capacity constraints over $s$ and using $l^e\ge l$ for odd $e$ and
the reindexing $\sum_s\sum_ex_{s-e,e}=V$ gives $lV+V\le\sum_sD_s$. By complementary slackness the heavy demands are
incoming flows, so $\sum_{\sigma_s\ge2/l}D_s\le V\le\sum_sD_s/(l+1)\le\sum_sD_s/l$.

(b) Take an optimal dual $\sigma$ for $\kappa w$, so $\mathrm{lp}_F(\kappa w)=\sum_s\kappa w_s(\gamma_l-\sigma_s)$, and put
$c=l^{-3}$, $\sigma'_s=\max(\sigma_s,c)\ge0$. For odd $e$: if $e<d$ then $l^e\sigma'_s+\sigma'_{s+e}\ge l^e\sigma_s+\sigma_{s+e}\ge1$; if
$e\ge3$ then $l^e\sigma'_s\ge l^3c=1$; as $d\ge2$ this covers all odd $e$. For each $s$,
\[
D_s(\gamma_l-\sigma'_s)=\kappa w_s(\gamma_l-\sigma_s)-\kappa w_s(\sigma'_s-\sigma_s)+(D_s-\kappa w_s)(\gamma_l-\sigma'_s),
\]
with $0\le\sigma'_s-\sigma_s\le c$ and $0\le D_s-\kappa w_s\le4w_s$. If $\sigma_s\ge2/l$, then $\gamma_l-\sigma'_s\ge-1$ (as
$\sigma'\le1$, $\gamma_l\ge0$), so the last term is $\ge-4w_s$. If $\sigma_s<2/l$, then $\sigma'_s<2/l$ (as $c<2/l$) and
$\gamma_l\ge1/l$, so $\gamma_l-\sigma'_s\ge-1/l$ and the last term is $\ge-4w_s/l$. Summing, and using (a) for $D=\kappa w$
(which gives $\sum_{\sigma_s\ge2/l}w_s\le\sum_sw_s/l$),
\[
\sum_sD_s(\gamma_l-\sigma'_s)\ge\mathrm{lp}_F(\kappa w)-\kappa c\sum_sw_s-4\frac{\sum_sw_s}l-4\frac{\sum_sw_s}l .
\qedhere
\]
\end{proof}

\begin{proposition}[(O8) Scaled generators]\label{pf:top-shape:o8}
If $E$ has the shape and $\sfname{CLarge}(C_6,C,K,c_e)$ holds, then statement (O8) of \cref{pf:top-shape:orbit} holds.
More precisely, for $5\le p\le X''(n)$ of residue class zero with $h_p\ne0$, $i<m$, $j=i+1$, and every kernel $\Theta$ of
theta product form, $\Theta$ has kernel valuation at least
\[
B=\frac1{n_p}\sum_{z\in\mathrm{BC}}F_{\rm Orb}(j,p,z)-\frac{(k-1)\sum_{\rm BC}w(j,z)}{p^3}-\frac{8\sum_{\rm BC}w(j,z)}p
-3f\,\mathrm{Nm}^2\log_pa
\]
on $\nu_{u,i}$, where $\mathrm{BC}$ is the set of places of kind B or C; and at the primes of type $C_s$, if the error
constants are large, $F_{\rm inert}(j,p)-\mathrm{Eat}_I(j,p)\le B$.
\formallabel{UpperTheory.kvalGe_theta_of_shape}
\formallabel{UpperTheory.fInert_sub_eatI_le_orbit_average_sub_errors}
\formallabel{UpperTheory.norm_mfact_mul_coeff_le_of_B0S}
\end{proposition}

\begin{proof}
\emph{Step 1 (radii).} The prime is not of signal type, so $\mathrm{BC}$ consists of the places of the second kind, none of
which is top; by the shape $(k-1)w_s\le|\nu_s|\le(k-1)w_s+4w_s$ there. On an orbit of $\mathrm{BC}$ apply
\cref{pf:top-shape:lp}(b) with $\kappa=k-1>0$, $w_s=w(j,r_s)$, $D_s=|\nu_{u,i,r_s}|$ (and $\mathrm{lp}_F=F_{\rm Orb}$ on the
orbit, \lref{UpperTheory.OrbitRadiiBounds.lpF_eq}); on an orbit of kind A take $\sigma\equiv1/(p+1)$, which is dual feasible
($p^e/(p+1)+1/(p+1)\ge1$) and has $\gamma_p-\sigma=1/(p^2-1)\ge0$. Transport these to the places (a place has a unique
representative and position, and $F_{\rm Orb}$ is constant along orbits). Summing over the orbits, with each $\mathrm{BC}$ orbit
of size $n_p$, we get radii $\sigma_z\ge0$ with $p^e\sigma_z+\sigma_{z_e}\ge1$ for odd $e$ and
\[
\frac1{n_p}\sum_{\rm BC}F_{\rm Orb}-\frac{(k-1)\sum_{\rm BC}w}{p^3}-\frac{8\sum_{\rm BC}w}p\le\sum_z|\nu_{u,i,z}|(\gamma_p-\sigma_z).
\]

\emph{Step 2 (coefficients).} For a theta product form $\Theta$ and radii $\sigma$, every exponent $\kappa$ with all
$\kappa_v\le A$ ($A\ge1$) satisfies
\[
|\kappa!\,\mathrm{coeff}_\kappa\Theta|\le p^{-\left(\sum_v\kappa_v(\gamma_p-\sigma_{z(v)})-N_{\rm var}(\log_pA+1)\right)},
\]
$N_{\rm var}$ the number of variables and $z(v)$ the place of $v$. Indeed, by Legendre's formula
$v_p(N!)=(N-s_p(N))/(p-1)\ge N/(p-1)-\log_pN-1$ \cite{Robert2000}, so $|\kappa!|\le p^{-(\sum_v\kappa_v/(p-1)-N_{\rm var}(\log_pA+1))}$.
Inverting the generalised logarithm in the product form, $\Theta$ (rescaled by the unit $\varepsilon$) is the product of an
integral series substituted into the inverse logarithm and the exponential of the $\beta$-series. The $\beta$-series has no
constant term and $|\mathrm{coeff}_n|\le|p|^{1/(p-1)-\sum_vn_v(\rho_p+\sigma_{z(v)})}$, $\rho_p=1/(p^2-1)$, by the radii
condition (\lref{UpperTheory.beS_constantCoeff_eq_zero_and_wtGe}); Dwork's estimates for exponentials, products,
substitutions and the inverse of the generalised logarithm
(\lref{Formalism.DworkPowerSeriesWeights.exponential_weight_le}, \lref{Formalism.DworkPowerSeriesWeights.multiplicative_weight_le}, \lref{Formalism.DworkPowerSeriesWeights.substitution_weight_bound}, and \lref{Formalism.DworkPowerSeriesWeights.generalized_log_weight};
\cite{Dwork1962,Robert2000}) then give $|\mathrm{coeff}_\kappa\Theta|\le|p|^{-\sum_v\kappa_v(\rho_p+\sigma_{z(v)})}$. Multiply, using
$1/(p-1)-\rho_p=\gamma_p$.

\emph{Step 3.} For $\kappa$ with $\sum_{a,b}\kappa(z,a,b)=|\nu_{u,i,z}|$ we have $\kappa_v\le|\nu_{u,i,z}|\le a^2$, so Step~2
applies with $A=a^2$; $\sum_v\kappa_v(\gamma_p-\sigma_{z(v)})=\sum_z|\nu_{u,i,z}|(\gamma_p-\sigma_z)$ and
$N_{\rm var}=fN_iN_{i+1}$. With $\ell=\log_pa\ge1$, $N_{\rm var}(\log_p(a^2)+1)=fN_iN_{i+1}(2\ell+1)\le3f\,\mathrm{Nm}^2\ell$.
By Step~1, $|\kappa!\,\mathrm{coeff}_\kappa\Theta|\le p^{-B}$, and \cref{pf:top-shape:kvalbasic}(a) gives the kernel bound.

\emph{Step 4 ($F_{\rm inert}-\mathrm{Eat}_I$).} At a prime of type $C_s$ (residue class zero, $h_p$ of odd order and $\ne0$)
the prime is not of signal type, so $F_{\rm inert}=\frac1{n_p}\sum_{\rm A}F_{\rm Orb}+\frac1{n_p}\sum_{\rm BC}F_{\rm Orb}$ and there
are no top places of kind B or C. By definition $\mathrm{Eat}_I$ is
$\max(0,\frac1{n_p}\sum_{\rm A}F_{\rm Orb})+c_6\frac{\sum_{\rm BC}w}p+|\mathrm{BC}|N_iN_{i+1}(\log_p(4n_p\mathrm{Nm}\,a)+1)
+\delta_1|\mathrm{BC}|\mathrm{Nm}^2\log_p(2a)+\frac{(k-1)\sum_{\rm BC}w}{p^3}+c_6f\,\mathrm{Nm}^2\ell$ (the terms with tops of kind
B or C vanish). Large error constants give $c_6\ge12$, so $8\sum w/p\le c_6\sum w/p$ and $3f\,\mathrm{Nm}^2\ell\le
c_6f\,\mathrm{Nm}^2\ell$; the remaining terms are non-negative. Hence $F_{\rm inert}-\mathrm{Eat}_I\le B$.
\end{proof}

\emph{Formal-group systems.} For a prime $p$, a \emph{formal-group system} $X$ at $p$ consists of local group data
over a discrete valuation ring $O$ with fraction field $K_f$, fillings with weights $w$ and type sets, levels and a
chain of lattices $\mathrm{Int}(j)$ with operators $T_m$ and transition maps $K_z$, standard data, a uniformiser
$\varpi$ with $p\sim\varpi^{e_O}$, bounds $d_G(j)$ for the $p$-adic valuations of the stabilisers of lattice tuples,
and bases at the places. \emph{Pending data} $Y$ for $X$ consist of a field $K_c\supseteq K_f$ with a rank-one valuation
ring $O_c$ extending the integers of $K_f$ and containing the Witt vectors, representatives of the lattice tuples
modulo the level groups, frames with theta series $\Theta_A$, $\Theta_F$, coefficients and the transfer maps
$\Delta_F$. For a datum $D$, $Y$ is \emph{admitted} if it satisfies thirteen local conditions (the first conditions
$\mathrm{adm}_1$--$\mathrm{adm}_{3b}$, archimedean signature matching, $\mathrm{adm}_5$, Minkowski counting,
$\mathrm{adm}_6$--$\mathrm{adm}_8$ and $\mathrm{adm}_{11}$).

\begin{theorem}[Formal-group statements]\label{pf:top-shape:packet}
Let $D$ be a datum, $p$ a prime, $X$ a formal-group system at $p$ and $Y$ pending data for $X$ which are admitted for
$D$. Then the eleven local conclusions hold: the conclusions $0$, $1$ and $2$ of the formal-group statement, the
multiplicativity of the expansions, the projection lemma at the base level, the volume ratio and the mass ratio, the
points lemma, the conclusions $4$ and $5$, and the bound for the low terms.
\formallabel{Formalism.MainNotions.FormalGroupSystem.Pf.packet_statements_hold}
\end{theorem}

\begin{proof}
This is the first half of \cref{pf:top-shape:prv} below, applied to the restriction statements of
\cref{pf:top-shape:rests}.
\end{proof}

\subsubsection*{The restriction statements and the chain estimates}

We prove \cref{pf:top-shape:packet} one level further. Let $X$ be a formal-group system at $p$ and $Y$ pending data.
The system carries a chain of lattices $\mathrm{Int}(j)$ ($0\le j\le\mathrm{mm}$) in $K_f$-spaces $F(j)$ with $K_f$-linear
transition maps $K_z\colon F(j)\to F(j+1)$; we write $\mathrm{val}_j(F)\ge g$ if $F\in\varpi^{\lceil ge_O\rceil}\mathrm{Int}(j)$.
At every place $v$ there are Hecke operators $T^{(i)}_v$ on $F(j)$ ($0\le i\le r$, $T^{(0)}_v=\mathrm{id}$), and for data
$i_0$ (with $i_0(v)\le r$, supported on a list $\mathrm{vs}$ of distinct places) $T_{i_0}=\prod_{v\in\mathrm{vs}}T_v^{(i_0(v))}$.
The system provides the gains $g(v,j)$, the base gains $g_0(v,j)$, the even parts $\mathrm{al}_{\rm ev}(v)$, the bounds
$d_G(j)$ for the $p$-adic valuations of the stabilisers, and the losses $\mathrm{Pr}(v,j)$. For $x\in K_c$ we write
$x\succeq g$ if $x^bp^{-a}\in O_c$ whenever $a\in\Z$, $b\ge1$ and $a\le gb$. The \emph{kernel entries} at level $j$ with
Hecke data $i_0$ are the coefficients of the polynomials obtained by applying the local differential operators
$\Pi(j,I,i_0)$ (for column choices $I$) to the theta series $\Theta_A$ of the frames between representative lattice tuples;
$\mathrm{EntB}(j,i_0,g)$ means that they all satisfy $\succeq g$. $\mathrm{LowAt}(j,v,g)$ is the local lower bound at $v$:
for every series $U$ with integral parameters, all coefficients of the local operator at $v$ applied to
$\exp(-\Phi_v)\,U$ satisfy $\succeq g$, where $\exp(-\Phi_v)$ is the Dwork exponential factor of $v$. Finally
\begin{itemize}
\item $\mathrm{PlaceB}(j,v,i,g)$: either $i=0$ and ($\mathrm{LowAt}(j,v,g)$ or $g\le g(v,j+1)-\mathrm{al}_{\rm ev}(v)$), or $1\le i\le r$
and $g\le g(v,j+1)$;
\item $\mathrm{FormB}$: for all $j<\mathrm{mm}$, $i_0$, $\mathrm{vs}$, $g$, if $\mathrm{EntB}(j,i_0,g)$ then
$\mathrm{val}_{j+1}(T_{i_0}K_zG)\ge g-d_G(j)$ for all $G\in\mathrm{Int}(j)$;
\item $\mathrm{EntCond}$: for all $j<\mathrm{mm}$, $i_0$ and rational $g_P(v)$, if $\mathrm{PlaceB}(j,v,i_0(v),g_P(v))$ for all $v$, then
$\mathrm{EntB}(j,i_0,\sum_vg_P(v))$.
\end{itemize}

\begin{lemma}[Additivity over places]\label{pf:top-shape:formplace}
Assume $\mathrm{FormB}$ and $\mathrm{EntCond}$. Let $j<\mathrm{mm}$, $i_0$, $\mathrm{vs}$ as above and real numbers $g_P(v)$ with
$\mathrm{PlaceB}(j,v,i_0(v),g_P(v))$ for all $v$. Then $\mathrm{val}_{j+1}(T_{i_0}K_zG)\ge\sum_vg_P(v)-d_G(j)$ for all $G\in\mathrm{Int}(j)$.
\formallabel{Formalism.MainNotions.FormalGroupSystem.Pf.form_of_place}
\formallabel{Formalism.MainNotions.FormalGroupSystem.Pf.ent_strict}
\formallabel{Formalism.MainNotions.FormalGroupSystem.Pf.formB_lim}
\end{lemma}

\begin{proof}
Put $g=\sum_vg_P(v)$ and let $\#$ be the number of places. \emph{Rational approximation.} Let $g'<g$,
$\epsilon=(g-g')/(\#+1)>0$ and choose rationals $q_v\in(g_P(v)-\epsilon,g_P(v))$. All alternatives of $\mathrm{PlaceB}$ are
monotone in $g$, so $\mathrm{PlaceB}(j,v,i_0(v),q_v)$ holds, and $\mathrm{EntCond}$ gives $\mathrm{EntB}(j,i_0,\sum_vq_v)$. As
$\sum_vq_v>g-\#\epsilon>g'$ and $\mathrm{EntB}$ is monotone, $\mathrm{EntB}(j,i_0,g')$ holds for every $g'<g$.
\emph{Discreteness.} Let $x=(g-d_G(j))e_O$ and $\delta=(x-\lceil x\rceil+1)/(e_O+1)>0$. Then $\delta e_O<x-\lceil x\rceil+1$, so
$\lceil x\rceil-1<x-\delta e_O\le\lceil x\rceil$ and $\lceil(g-\delta-d_G(j))e_O\rceil=\lceil x\rceil$. By $\mathrm{FormB}$ with
$g-\delta<g$, $T_{i_0}K_zG\in\varpi^{\lceil x\rceil}\mathrm{Int}(j+1)$, which is the claim.
\end{proof}

\begin{proposition}[Chain estimates]\label{pf:top-shape:chainest}
Assume $\mathrm{FormB}$, $\mathrm{EntCond}$ and $\mathrm{LowAt}(j,v,g_0(v,j+1))$ for all $v$ and $j<\mathrm{mm}$ (the volume
ratio). Let $j<\mathrm{mm}$ and $G\in\mathrm{Int}(j)$.
\begin{enumerate}[label=(\alph*)]
\item (Conclusion 0.) If all kernel entries with $i_0=0$ satisfy $\succeq0$ (admission condition $\mathrm{adm}_3$), then
$\mathrm{val}_{j+1}(K_zG)\ge-d_G(j)$.
\item (Conclusion 2.) $\mathrm{val}_{j+1}(K_zG)\ge\sum_vg_0(v,j+1)-d_G(j)$.
\item (Projection lemma at the base level.) If the $g_0(v,j+1)$ are rational, the kernel entries with $i_0=0$ satisfy
$\succeq\sum_vg_0(v,j+1)$.
\item (Mass ratio.) $\mathrm{val}_{j+1}(K_zG)\ge g(v,j+1)-\mathrm{al}_{\rm ev}(v)+\sum_{v'\ne v}g_0(v',j+1)-d_G(j)$ for every $v$.
\item (Points lemma.) If the $T_v^{(i)}$ preserve $\mathrm{Int}(j+1)$, then for $1\le i\le r$,
$\mathrm{val}_{j+1}(T^{(i)}_vK_zG)\ge g(v,j+1)+\sum_{v'\ne v}g_0(v',j+1)-d_G(j)$.
\item If moreover $\mathrm{LowAt}(j,v,0)$ (condition $\mathrm{adm}_{11}$), then $\mathrm{LowAt}(j,v,g)$ for every $g\le\max(g_0(v,j+1),0)$.
\end{enumerate}
\formallabel{Formalism.MainNotions.FormalGroupSystem.Pf.form_valuation_of_bound}
\formallabel{Formalism.MainNotions.FormalGroupSystem.Pf.slopeEstimate_of_lemmaB_slopeLine}
\formallabel{Formalism.MainNotions.FormalGroupSystem.Pf.theorem_A_entries_of_corollary}
\formallabel{Formalism.MainNotions.FormalGroupSystem.Pf.hecke_integral_of_forms}
\formallabel{Formalism.MainNotions.FormalGroupSystem.Pf.hecke_valuation_of_entries}
\formallabel{Formalism.MainNotions.FormalGroupSystem.Pf.latticeChain_succ_of_partial}
\formallabel{Formalism.MainNotions.FormalGroupSystem.lowAt_max}
\end{proposition}

\begin{proof}
(a) $\mathrm{adm}_3$ is $\mathrm{EntB}(j,0,0)$; apply $\mathrm{FormB}$ with $i_0=0$, $\mathrm{vs}=\emptyset$, $g=0$.
(b) $\mathrm{PlaceB}(j,v,0,g_0(v,j+1))$ holds by the volume ratio; apply \cref{pf:top-shape:formplace} with $i_0=0$.
(c) With $q_v=g_0(v,j+1)\in\Q$, $\mathrm{PlaceB}(j,v,0,q_v)$ holds; apply $\mathrm{EntCond}$.
(d) Take $g_P(v)=g(v,j+1)-\mathrm{al}_{\rm ev}(v)$ (the second alternative of $\mathrm{PlaceB}$ with $i=0$) and
$g_P(v')=g_0(v',j+1)$ for $v'\ne v$; apply \cref{pf:top-shape:formplace} with $i_0=0$.
(e) Take $i_0=i$ at $v$ and $0$ elsewhere, $\mathrm{vs}=[v]$, $g_P(v)=g(v,j+1)$ (third alternative, as $1\le i\le r$) and
$g_P(v')=g_0(v',j+1)$; apply \cref{pf:top-shape:formplace}. With the integrality of the Hecke operators
(\lref{Formalism.MainNotions.FormalGroupSystem.Pf.restricted_hecke_integral}) this is the points lemma.
(f) If $g_0(v,j+1)\le0$ then $g\le0$ and $\mathrm{LowAt}(j,v,0)$ gives the claim by monotonicity; otherwise $g\le g_0(v,j+1)$ and
the volume ratio gives it.
\end{proof}

\begin{lemma}[Low terms]\label{pf:top-shape:lowterm}
Let $G_n,G_{0n},\mathrm{PR}_n$ be functions of the places and levels, $S$ a function of the places and $\mathrm{VV}$ a set of places
with $G_n(v,j)\le g(v,j)$ and $\mathrm{Pr}(v,j)\le\mathrm{PR}_n(v,j)$ for $v\in\mathrm{VV}$, $1\le j\le\mathrm{mm}$,
$\mathrm{LowAt}(j,v,G_{0n}(v,j+1))$ for all $v$, $j<\mathrm{mm}$, and $\sum_{j=1}^{\mathrm{mm}}\max(G_n-\mathrm{PR}_n-G_{0n},0)(v,j)\le S(v)$
for $v\in\mathrm{VV}$. Assume $\mathrm{FormB}$, $\mathrm{EntCond}$, conclusion $4$ and that the operators $y_v$ preserve $\mathrm{Int}(0)$
(or that $e_O=0$ and conclusion $0$ holds). Then for every $F$ in the product lattice $\mathrm{Iprod}(S,\mathrm{VV})$,
\[
\mathrm{val}_{\mathrm{mm}}\bigl(K_z^{(\mathrm{mm})}F\bigr)\ge\sum_{j=1}^{\mathrm{mm}}\Bigl(\sum_{v\in\mathrm{VV}}\max\bigl(G_n-\mathrm{PR}_n,G_{0n}\bigr)(v,j)
+\sum_{v\notin\mathrm{VV}}G_{0n}(v,j)\Bigr)-\sum_{j<\mathrm{mm}}d_G(j).
\]
\formallabel{Formalism.MainNotions.FormalGroupSystem.Pf.low_term_bound_of_eO_eq_zero_of_aux0}
\formallabel{Formalism.MainNotions.FormalGroupSystem.Pf.main_ind}
\formallabel{Formalism.MainNotions.FormalGroupSystem.Pf.lift_bound_of_integrality}
\end{lemma}

\begin{proof}
\emph{Case $e_O=0$.} Every place has $e_O>0$ (\lref{Formalism.MainNotions.FormalGroupSystem.Pf.eO_pos}), so
there are no places; and $\varpi^{\lceil g\cdot0\rceil}=1$, so $\mathrm{val}_j(F)\ge g$ just means $F\in\mathrm{Int}(j)$. The product
lattice is then $\mathrm{Int}(0)$, and conclusion $0$ gives $K_z\mathrm{Int}(j)\subseteq\mathrm{Int}(j+1)$; by induction
$K_z^{(m')}F\in\mathrm{Int}(m')$ for $m'\le\mathrm{mm}$, which is the claim.

\emph{Case $e_O>0$ (induction over the levels).} Let $\mathrm{Lv}$ be a list of distinct places and, for every level $j$, let
$B_j\subseteq\mathrm{Lv}$ be a set of places to be ``boosted''. Conclusion $4$ provides for every $v$ units $u_l$ and elements
$m_l$ of valuation $v_m(v,l)$ and an operator $Z_v$ with $Z_v(K_z^{(j+1)}F)=p^{\mathrm{Pr}_Z(v,j+1)}K_z^{(j+1)}(y_vF)$, whose
coefficients have valuation $\ge0$ (and $\ge\mathrm{al}_{\rm ev}(v)$ for the last one). By induction on $m'$ one shows
\[
\mathrm{val}_{m'}\bigl(K_z^{(m')}(P^{m'}F_1)\bigr)\ge\sum_{j<m'}\Bigl(\sum_{v\in B_j}\bigl(G_n(v,j+1)-\mathrm{Pr}_Z(v,j+1)\bigr)
+\sum_{v\notin B_j}G_{0n}(v,j+1)-d_G(j)\Bigr)
\]
for $F_1\in\mathrm{Int}(0)$, where $P$ is the product of the $y_v$, $v\in\mathrm{Lv}$: in the step $m'\to m'+1$ one writes
$P^{m'+1}=P_{B}\cdot P^{m'}\cdot P_{C}$ ($B=B_{m'}$, $C=\mathrm{Lv}\setminus B$, the $y_v$ commute), converts $P_B$ into the operators
$Z_v$ at the cost $p^{-\sum_{v\in B}\mathrm{Pr}_Z(v,m'+1)}$, writes $K_z^{(m'+1)}(P^{m'}P_CF_1)=\varpi^{\lceil\cdot\rceil}K_zG'$ with $G'\in\mathrm{Int}(m')$ by
the induction hypothesis, and bounds the $Z$-operators applied to $K_zG'$ by \cref{pf:top-shape:formplace} with
$g_P(v)=G_n(v,m'+1)$ on $B$ and $g_P(v)=G_{0n}(v,m'+1)$ elsewhere
(\lref{Formalism.MainNotions.FormalGroupSystem.Pf.ZL_bound}). The lemma follows by choosing $B_j$ to be the places of
$\mathrm{VV}$ where $G_n-\mathrm{PR}_n\ge G_{0n}$ at level $j$, so that every place contributes $\max(G_n-\mathrm{PR}_n,G_{0n})$;
the integrality of the elements of $\mathrm{Iprod}(S,\mathrm{VV})$ is where $\sum_j\max(G_n-\mathrm{PR}_n-G_{0n},0)\le S$ enters.
\end{proof}

The Hecke condition rests on a count of lattices. For a discrete valuation ring $W$ with uniformiser $p$ and residue field
$k$ with $q=p^{n'}$ elements, an \emph{edge} of type $i$ is a pair of lattices $L_1\subseteq L_2$ with $pL_2\subseteq L_1$ and
$L_2/L_1\cong k^i$. Let $\binom{n}{l}_q$ be the Gaussian binomial coefficient.

\begin{lemma}[Hecke counting]\label{pf:top-shape:countloc}
Let $A\subseteq\mathrm{Frac}(W)^N$ and $B\subseteq\mathrm{Frac}(W)^{N+g}$ be lattices and $x$ an $(N+g)\times N$ matrix. Then
\begin{multline*}
\#\{B_1:(B_1,B)\text{ an edge of type }i,\ xA\subseteq B_1\}\\
=\sum_{l=0}^iq^{(i-l)(g-l)}\binom gl_q\,\#\{A_1:(A,A_1)\text{ an edge of type }i-l,\ xA_1\subseteq B\}.
\end{multline*}
\formallabel{Formalism.MainNotions.FormalGroupSystem.Pf.countLoc}
\formallabel{Formalism.MainNotions.FormalGroupSystem.Pf.Cnt.core}
\end{lemma}

\begin{proof}
If $xA\not\subseteq B$ both sides vanish ($B_1\subseteq B$, resp.\ $A\subseteq A_1$). Let $xA\subseteq B$, let $V=B/pB\cong
k^{N+g}$ and let $U\subseteq V$ be the image of $xA$, of dimension $u$. The lattices $B_1$ of the left side correspond to the
subspaces $B_1/pB$ of codimension $i$ of $V$ containing $U$, so the left side is the number $\binom{N+g-u}{i}_q$ of
codimension-$i$ subspaces of $V/U$. The lattices $A_1$ of the right side correspond to the subspaces $A_1/A$ of
$p^{-1}A/A\cong k^N$ of dimension $i-l$ contained in the kernel of the map $p^{-1}A/A\to p^{-1}B/B$ induced by $x$;
multiplication by $p$ identifies this map with the reduction $A/pA\to B/pB$ of $x$, whose image is $U$, so the kernel has
dimension $N-u$ and the number is $\binom{N-u}{i-l}_q$. The claim is the $q$-Vandermonde identity
$\binom{M+g}{i}_q=\sum_lq^{(i-l)(g-l)}\binom gl_q\binom{M}{i-l}_q$ with $M=N-u$ \cite[Thm.~3.4]{Andrews1976}.
\end{proof}

\begin{theorem}[Restriction statements]\label{pf:top-shape:rests}
If $Y$ is admitted for $D$, then: (R1) conclusion $1$ holds (the local theta series satisfy the $S$-conclusion at every
place); (R2) the expansions multiply; (R3) the volume ratio $\mathrm{LowAt}(j,v,g_0(v,j+1))$ holds for all $v$ and
$j<\mathrm{mm}$; (R4) $\mathrm{EntCond}$ holds.
\formallabel{Formalism.MainNotions.FormalGroupSystem.Pf.rests}
\formallabel{Formalism.MainNotions.FormalGroupSystem.Pf.theta_kernel_expansion}
\formallabel{Formalism.MainNotions.FormalGroupSystem.Pf.expansions_multiply}
\formallabel{Formalism.MainNotions.FormalGroupSystem.Pf.first_theorem_place_bound}
\formallabel{Formalism.MainNotions.FormalGroupSystem.Pf.second_lemma_place_bound}
\formallabel{Formalism.MainNotions.FormalGroupSystem.Pf.hecke_place_bound}
\formallabel{Formalism.MainNotions.FormalGroupSystem.Pf.ea_A1U}
\formallabel{Formalism.MainNotions.FormalGroupSystem.Pf.ea_A2U}
\formallabel{Formalism.MainNotions.FormalGroupSystem.Pf.ea_BU}
\formallabel{Formalism.MainNotions.FormalGroupSystem.Pf.auxU_of_firstTheoremPlaceBound}
\end{theorem}

\begin{proof}
Two algebraic inputs are used: the bi-additivity in the second variable of the Dwork-type series attached to the theta
data (\lref{Formalism.MainNotions.FormalGroupSystem.Pf.restBeta}) and a congruence for two sums of Frobenius twists,
derived from the Frobenius congruences of Witt vectors
(\lref{Formalism.MainNotions.FormalGroupSystem.Pf.restGc} and \lref{Formalism.MainNotions.FormalGroupSystem.Pf.restGcA}).
(R1) Under admission (condition $\mathrm{adm}_2$: the local extension condition at every place) these two inputs give
the $S$-conclusion for the local theta series
(\lref{Formalism.MainNotions.FormalGroupSystem.Pf.baseChange_preserves_unitary_of}). (R2) Likewise they give the
multiplicativity of the expansions
(\lref{Formalism.MainNotions.FormalGroupSystem.Pf.expansion_all_places_of_computation}).

(R3) Fix $v$ and $j<\mathrm{mm}$. By the reduction to a local norm bound
(\lref{Formalism.MainNotions.FormalGroupSystem.Pf.ea_lowN}, using the norm restriction at $v$ and the restriction of
the $S$-expansion, \lref{Formalism.MainNotions.FormalGroupSystem.Pf.nm_restNorm_of_place} and \lref{Formalism.MainNotions.FormalGroupSystem.Pf.ea_restSU}),
it suffices to show: over every ultrametric field $K'$ containing a root $\varpi'$ of $\varpi'^{\,p^2-1}=p$ with $|p|=p^{-1}$, for
units $u_s$ ($|u_s|\le1$) and every series $F$ whose rescaling $F(\varpi'\cdot)$ lies in the local class of $v$, all
coefficients of the local operator applied to $F$ have absolute value $\le p^{-g_0(v,j+1)}$. Each such coefficient is $0$
or a coefficient of the frame jet of the restriction of $F$ to the variables of $v$
(\lref{Formalism.MainNotions.FormalGroupSystem.Pf.eb_brJet}); the restriction lies in the class of the generic kernel
theorem (\lref{Formalism.MainNotions.FormalGroupSystem.Pf.ea_brClass}); and the numbers of the generic theorem are
those of the system: its excess gain is $g_0(v,j+1)$, its gain is $g(v,j+1)$, and the sum of its even
$\mathrm{al}$-terms is $\mathrm{al}_{\rm ev}(v)$ (\lref{Formalism.MainNotions.FormalGroupSystem.Pf.ea_brNum}). The generic
kernel bound A1 (\lref{Formalism.InvariantFactorPaths.G7.kernel_a_one_generic}) gives the bound $p^{-g_0(v,j+1)}$.

(R4) In the same way, the generic bound A2
(\lref{Formalism.InvariantFactorPaths.G7.kernel_a_two_generic}) gives $\mathrm{LowAt}(j,v,g(v,j+1)-\mathrm{al}_{\rm ev}(v))$,
and, for $1\le i\le r$, the generic bound B (\lref{Formalism.InvariantFactorPaths.G7.kernel_b_generic}) together
with the scaling of the local operator by the Hecke data
(\lref{Formalism.MainNotions.FormalGroupSystem.Pf.ea_brScale}: the coefficient is $0$ or $p^{\mathrm{Exp}}$ times the
unscaled one) gives the local bound $g(v,j+1)$ for the operator with Hecke data $i$. The local bounds of the places
combine into bounds for the kernel entries, which is $\mathrm{EntCond}$
(\lref{Formalism.MainNotions.FormalGroupSystem.Pf.restEnt_of}).
\end{proof}

\begin{theorem}\label{pf:top-shape:prv}
If $Y$ is admitted for $D$ and the restriction statements (R1)--(R4) hold, then the eleven conclusions of
\cref{pf:top-shape:packet} hold, and the kernels commute with the transition maps up to units.
\formallabel{Formalism.MainNotions.FormalGroupSystem.Pf.prv_and_kernel_commutation_up_to_units_of_rests}
\formallabel{Formalism.MainNotions.FormalGroupSystem.Pf.ramifiedLattice_of_chain_estimates}
\end{theorem}

\begin{proof}
Under admission, the kernel formula $\mathrm{FormB}$ holds (\lref{Formalism.MainNotions.FormalGroupSystem.Pf.formB}),
and the Hecke condition holds
(\lref{Formalism.MainNotions.FormalGroupSystem.Pf.heckeCondition_of_adm}, from the counting of
\cref{pf:top-shape:countloc} and the Hecke operators on the quotient spaces,
\lref{Formalism.MainNotions.FormalGroupSystem.Pf.restTmFQ}). Conclusions $1$, the multiplicativity of the
expansions and the volume ratio are (R1)--(R3). With $\mathrm{FormB}$, (R4) and the volume ratio,
\cref{pf:top-shape:chainest} gives conclusion $0$ (with $\mathrm{adm}_3$), conclusion $2$, the projection lemma (the
$g_0$ are rational, \lref{Formalism.MainNotions.FormalGroupSystem.Pf.restRat}), the mass ratio and the points
lemma. The Hecke condition, the integrality of the Hecke operators, $\mathrm{adm}_{3b}$, the points lemma and the mass
ratio give the commutation of the kernels up to units
(\lref{Formalism.MainNotions.FormalGroupSystem.Pf.alternate_expansion_of_relations}), of which conclusion $4$ is a
reformulation (\lref{Formalism.MainNotions.FormalGroupSystem.Pf.theorem_of_alternate_form}). The low-term bound is
\cref{pf:top-shape:lowterm} (with $\mathrm{adm}_8$ for the $y_v$). Finally conclusion $5$: its first half
$S_v\le\mathrm{mm}\,\gamma_p\sum_{s\in\mathrm{TT}(v)}\mathrm{al}(v,s)$ is the relative presentation bound
(\lref{Formalism.MainNotions.FormalGroupSystem.Pf.relative_presentation_of_d7}), and its second half is the
low-term bound with $G_n=g$, $G_{0n}=g_0$ (by the volume ratio), $\mathrm{PR}_n=\mathrm{Pr}$ and $S(v)=S_v$.
\end{proof}

\emph{The local service law.} Let $(D,L)$ be a member with its instance data, twisted forms, level data and error data.
For $C\in\R$, $\sfname{Serv}(C)$ is the conjunction of: (S1) the size conditions at the type primes with constant
$C$; (S2) under the ties and the error definitions of the member: for every word $u$, every place $w$ whose residue
characteristic $p$ is different from $T_u$, $g_0$, $g_1$, every used twist $\mu$ with $p$ not in the prime set of
$\mu$, and every $x$ in the lattice of $u$, the image of $x$ under the column chain lies in $a\cdot\Lambda_w$ for an
element $a$ which is $0$ or has $w$-valuation at least $\mathrm{ur}(D,w,s)$, where $\Lambda_w$ is the local lattice and
\[
s=\sum_{j=1}^m\bigl(\mathrm{gain}_S(p,j)+F(j,w)-E_S(u,j,w)\bigr)+v_p(D_\Gamma(u));
\]
(S3) the same statement for the elements $x$ of the ramified part, with $s=\sum_j(F(j,w)-E_S(u,j,w))+v_p(D_\Gamma(u))$.
$\sfname{ServF}$ is the family property ``there is $C\ge0$ such that eventually every member satisfies
$\sfname{Serv}(C)$''.

The \emph{sum-over-primes hypotheses} on a family class $\mathcal C$ are the closedness on $\mathcal C$ of: (twU) the
numerical comparison of the members and the existence of twisted formal-group systems at the relevant places; (pdv)
the existence, eventually at every member, of an orbit-entry kernel record $E$ satisfying the unit, rationality,
filtration and stabiliser conditions of the kernels, of theta constants $(\varepsilon,\beta,\sigma)$ satisfying the
even conditions at the inert primes and at the primes of class $C$, of the orbit bound without operator, of the jet
integrality, of the shape of the kernels at the ramified places and of the shape of $E$.

\begin{theorem}[Sum-over-primes service]\label{pf:top-shape:serv}
Let $\mathcal C$ be a family class on which $\sfname{DGamSize}$ is closed and which satisfies the sum-over-primes
hypotheses. Then $\sfname{ServF}$ is closed on $\mathcal C$.
\formallabel{Formalism.MainNotions.Pk.sum_over_primes_service}
\end{theorem}

\begin{proof}
By the closedness calculus (\cref{pf:top-shape:calculus}(b)) the four properties $\sfname{DGamSize}$, (twU), (pdv)
and the numerical exponent statement (which is closed on every class,
\lref{Formalism.MainNotions.G.num_closed}) are closed together, and it suffices to derive $\sfname{ServF}$ at every
family $\topshapeF$ with these properties whose constants are large and positive
(\lref{Formalism.MainNotions.Pk.closed_large_pos}). Fix such $\topshapeF$.

\emph{Step 1 (constants).} The fixed constants of the family give $\mathrm{Pr}_1\ge1$ and $c_6=12$, $c_{\rm dig}=3$,
$\delta_1=5$, $\delta_2\ge0$, hence $\sfname{CLarge}(11,1,K,c_e)$: indeed $1\ge0$, $11\le12$, $3\le1+5$, $5\ge0$,
$\delta_2\ge0$ and $3\le3$. The orbit letter bounds $\sfname{OLB}(11,1)$ hold
(\lref{Formalism.MainNotions.DIV.Bridge.DworkPowerSeriesWeights.split_bound_eleven_one}).

\emph{Step 2 (type primes).} From $\sfname{DGamSize}$, \lref{Formalism.MainNotions.Pk.admissible_size_family}
gives $C_s\ge0$ and a threshold beyond which the members satisfy (S1) with constant $C_s$. We show that $C_s$ serves
for $\sfname{ServF}$; let $x$ be beyond all thresholds below, $x>x_1$, and $(D,L)$ a member.

\emph{Step 3 (unit bound).} (twU) provides the comparison data and the twisted formal-group systems at the member.
The formal-group statements (\cref{pf:top-shape:packet}, in the form of the statement $\sfname{Stmt}(D)$,
\lref{Formalism.MainNotions.FormalGroupSystem.Pf.stmt_proved}), the numerical exponents and the largeness of the
parameters (\lref{Formalism.MainNotions.G.TopChoiceBounds.hLarge_proved}) give, by
\lref{Formalism.MainNotions.G.Joints.unit_bound_of_sum_over_primes}, the unit bound at the twisted places (a
lower bound for the growth of the local lattices at the twisted places in terms of the gains at the orbit
representatives).

\emph{Step 4 (orbit admissibility).} (pdv) provides $E$ with its kernel conditions and theta constants. The type
primes differ from $d$ (\lref{Formalism.MainNotions.Pk.typeG_ne_d}); the orbits are admissible because the split
of the family is at least $3$ (\lref{Formalism.MainNotions.DIV.Inh.orbit_admissible_of_split_bound}); the theta
isometry facts (\lref{Formalism.MainNotions.DIV.Lemmas.theta_isometry_facts_hold}) and the kernel conditions give
the kernel form of the induced record with parameter $s_p=1$
(\lref{Formalism.MainNotions.OrbitPlaces.kform_of_ker}). With Step~1 and the remaining fields of (pdv), the
orbit data are admitted with constants $(11,1)$.

\emph{Step 5 (orbit-entry bounds).} By \cref{pf:top-shape:orbit} with $(C_6,C)=(11,1)$, the eight orbit statements
hold. Together with the admission they give the lifting conclusion for the orbit data
(\lref{Formalism.MainNotions.OrbitPlaces.lift_bounds_of_essential}), from which
\lref{Formalism.MainNotions.J.four_orbit_conditions_of_proof} extracts four local conditions: the orbits without
operator, the length bound at medium inert primes, the loss at the discriminant, and the integral kernel valuation
(the isotypic parts form an internal direct sum).

\emph{Step 6 (conclusion).} The error-exponent form of the formal-group statements
(\lref{Formalism.MainNotions.FormalGroupSystem.Pf.stmtAt_proved}, again a consequence of
\cref{pf:top-shape:packet}) gives the length bound at small inert primes of small order
(\lref{Formalism.MainNotions.WeilJacquetShalikaLaw.local_theorem_of_statement}). The unit bound of Step~3, this
length bound and the four conditions of Step~5 are the six hypotheses of the local service law; with the comparison
data and (S1), \lref{Formalism.MainNotions.WeilJacquetShalikaLaw.serv_of} (applied to the proof of the
Weil--Jacquet--Shalika statement, \lref{Formalism.MainNotions.WeilJacquetShalikaLaw.Proof.stmt_proved}) gives
$\sfname{Serv}(C_s)$ for the re-indexed data, and hence for the member
(\lref{Formalism.MainNotions.Pk.serv_re}).
\end{proof}

\begin{corollary}\label{pf:top-shape:servT}
Let $\mathcal C$ be a family class with a twist selection $\sigma$ which is served (the selection's tied records provide
the comparison data, the twisted formal-group systems and the kernel entries of the sum-over-primes hypotheses). If
$\sfname{DGamSize}$ is closed on $\mathcal C$, then $\sfname{ServF}$ is closed on $\mathcal C$.
\formallabel{Formalism.L2.G.ServT.closed_servF_of_servT_of_dGamSize}
\end{corollary}

\begin{proof}
A served selection yields the sum-over-primes hypotheses: the field (twU) is read from the selection, and (pdv) from
the kernel entries of the selection's tied records
(\lref{Formalism.L2.G.ServT.sumOverPrimesService_of_servT}). \Cref{pf:top-shape:serv} concludes.
\end{proof}

\begin{proposition}[Admissibility from the service law]\label{pf:top-shape:adminc}
Let $\mathcal C$ be a family class on which the following eight eventual properties are closed: (1) the sixth
conclusion of the zeroth-first parts at every word; (2) integrality of the twisted points; (3) the Weyl class-sum of
the periods; (4) admissibility of the twisted periods; (5) $\sfname{ServF}$; (6) the sharp fifth conclusion and the
twisted column bound at $T_u$ at every word; (7) the level primes are coprime to the stabilisers; (8) there is $C_1$
such that the used tuples are bounded by $C_1$ and the family size bound at the type primes holds with $C_1$. Then the
property ``eventually $\mathrm{adm}_2(D,L)$, and eventually $\mathrm{adm}_{12}(D,L)$'' is closed on $\mathcal C$.
\formallabel{Formalism.L2.G.Adm2InU.admissible_family_of_inclusion}
\end{proposition}

\begin{proof}
The conjunction of two closed properties is closed: if $(\eps_1,K_1)$ and $(\eps_2,K_2)$ witness the closedness, then
$\min(\eps_1,\eps_2)$ and the componentwise maxima of the thresholds $K_1$, $K_2$ witness it for the conjunction,
since a family large for the maximum is large for both. Iterating, the eight properties are closed together. Let
$\topshapeF$ satisfy them, with thresholds $y_1,\dots,y_7$ and the constant $C$ of $\sfname{ServF}$. For
$x\ge\max(y_1,\dots,y_7,3^{5/\eps_z})$ with $x>x_1$ and every member, the eight properties, the error data, the size
data at the level primes and the local compatibility of the member give the second-level bound and the fourth
complement condition (\lref{Formalism.L2.G.Adm2InU.admissible_level_two}), which are $\mathrm{adm}_2$ and
$\mathrm{adm}_{12}$ (\lref{Formalism.L2.G.Adm2InU.famU}).
\end{proof}

% ===========================================================================
\subsection{Mass ratios, slopes and the global link}\label{sec:pf:top-shape:gl}

For a word $u$ of a datum $D$ let $M_u$ be the hermitian lattice at level $0$ attached to $u$ (denoted $\Lambda_u$ in
\cref{prop:families:famgl}), with mean slope $\mu(M_u)$ and maximal slope $\mu_{\max}(M_u)$. A \emph{column valuation}
of $u$ is a record describing $M_u$ in polynomial coordinates: for each class $c$ of the level-zero genus of $u$ and
each embedding type $\tau$, a hermitian lattice $Q_{c,\tau}$ over $\OO_{\mathit{LL}}$ with a semilinear base-change map
from the base space, and an injective linear map from the space of $M_u$ to tuples of polynomials whose image is
described by a span condition, whose local membership at every finite place is described by integrality after the
action of local bases, and whose norm at every infinite place is a sum of Fock-type norms. For a word with column
valuations put $a_\tau$, $a'_\tau$ (column sizes), $s_\tau$ (spread of the first column) and $\lambda_\tau$ (last entry
of the first column), as in \cref{prop:families:famgl}.

\begin{definition}[The eleven properties]\label{pf:top-shape:famglmin}
For a family class $\mathcal C$ and $C_5\in\R$, $\sfname{FamGLmin}(\mathcal C,C_5)$ is the conjunction of the
closedness on $\mathcal C$ of the following eventual statements on the members $D$ of a family (those marked $(\ast)$
are required for every word $u$ of $D$ carrying a Hida setting):
\begin{enumerate}[label=(G\arabic*)]
\item\label{G:mass} the mass ratios $\mathrm{mass}(D,M)/\mathrm{mass}(D,M-1)$, $M\ge3$, equal the right-hand side of
\cref{prop:families:mass};
\item the volume of the base space $Y_0$ of $D$ equals the mass $\mathrm{mass}(D,r)$;
\item\label{G:sigma} for $e^{-r}\le\eps_0<\frac12$ and every row type $\rho$ whose parameter has gap at least an explicit
bound, $(1-\eps_0)\dim V(\rho)\le d_\sigma(\rho)\le4\,d_\tau(\rho)\,(3pq\,|\lambda(\rho)|)^{pq}$;
\item\label{G:low} every word $u$ with column valuations satisfies
\[
\mu(M_u)\ge-\sum_\tau\frac{a_\tau}k\log T_u-\sum_\tau\bigl(a'_\tau\log k+k^2\log(s_\tau+k)\bigr);
\]
\item\label{G:max} every word $u$ with column valuations satisfies
\[
\mu_{\max}(M_u)\le\sum_\tau\bigl(C_5\,a'_\tau\log(kv)+k^2\log(s_\tau+k)+|\lambda_\tau|\log T_u\bigr);
\]
\item\label{G:colval} every word has column valuations;
\item[(G7)--(G11)] $(\ast)$ free $T$-data, the slope line, Hecke integrality at level zero, ramified level zero, and the
second-branch inclusion length.
\end{enumerate}
\formallabel{Formalism.MainNotions.Wire.FamGLmin}
\end{definition}

\begin{proposition}\label{pf:top-shape:famgl}
Let $\mathcal C$ be a family class on which the real-place second-layer conditions $\mathrm{tie}$ are closed, and
assume $h_U$ and $h_G$. Then $\sfname{FamGLmin}(\mathcal C,4)$ holds.
\formallabel{Formalism.MainNotions.Wire.famGLmin_of_realPlaceSecondLayer}
\end{proposition}

\begin{proof}
Each property is obtained from a lemma, combined with $\mathrm{tie}$ by \cref{pf:top-shape:calculus}(a),(c).
\ref{G:mass} holds for every family, by the Siegel mass ratio of \cref{prop:families:mass}
(\lref{Formalism.L2.GlobalTheoryTheory.GenusMassVolumes.siegel_mass_ratio_family}). For (G2): the real-place
ties give eventually the volume tie of every member
(\lref{Formalism.MainNotions.family_conditions_of_layer_two}); then $h_U$ gives the volume ratio identity, which
unfolds to the mass-ratio form of the volume recursion
(\lref{Formalism.MainNotions.StandardLattice.volumeMass_of_unfolding}), $h_G$ gives the volume at level zero in
rank $2$ (\lref{Formalism.L2.GlobalTheoryTheory.GenusMassVolumes.volume_mass_of_genus_rank}), and the two give
the equality in rank $r$ (\lref{Formalism.L2.GlobalTheoryTheory.GenusMassVolumes.volume_eq_mass_rank_two}).
\ref{G:sigma} holds for every family (\lref{Formalism.MainNotions.GappedMultiplicityBounds.explicit_bound_family}).
\ref{G:low} and \ref{G:max} with $C_5=4$ hold for every family: these are the Arakelov slope inequalities
\cref{pf:top-shape:lowslope,pf:top-shape:upslope} below, which rest on the hook-length bound \cref{pf:top-shape:hook}.
\ref{G:colval} follows from the ties by \cref{pf:top-shape:tcv}. (G7), (G8), (G11) follow from the
second-layer conditions and \lref{Formalism.MainNotions.RT.slope_free_branch_family}, with
\lref{Formalism.MainNotions.freeTT_holds}, \lref{Formalism.MainNotions.slopeLineT_holds}, and \lref{Formalism.L2.GlobalTheoryTheory.LemB.second_branch_all_members};
(G9) from \lref{Formalism.MainNotions.RT.hI_tied} and the integrality at the base point lattice
\lref{PI.IntegralAtBasePointLattice.integral_at_t0_lattice}; (G10) from
\lref{Formalism.MainNotions.ramified_level_zero_holds}.
\end{proof}

\subsubsection*{The hook-length bound}

Both slope inequalities use the following bound for the Fock norms of highest weight vectors. Let $N\ge1$ and let
$\Q[x]$ be the polynomial ring in the $N^2$ variables $x_{ab}$ ($0\le a,b<N$), with the Fock (Fischer) inner product
\[
\langle f,g\rangle=\sum_mf_mg_m\,m!,\qquad m!=\prod_{a,b}m_{ab}!,
\]
the sum running over exponent vectors $m$ and $f_m,g_m$ denoting coefficients. Let $\Delta_s$ be the leading principal
$s\times s$ minor of the matrix $X=(x_{ab})$, $\partial_{ab}=\partial/\partial x_{ab}$, and
$E_{ab}=\sum_sx_{as}\partial_{bs}$ the polarization operators. For a non-increasing $\nu\colon\{0,\dots,N-1\}\to\N$ put
$\nu_N=0$,
\begin{gather*}
h_\nu=\prod_{s=0}^{N-1}\Delta_{s+1}^{\nu_s-\nu_{s+1}},\qquad c(\nu)=\frac{\langle h_\nu,h_\nu\rangle}{\prod_l\nu_l!},\\
P(\nu)=\prod_{i<j}(\nu_i-\nu_j+j-i),\qquad Q(\nu)=\prod_i(\nu_i+N-1-i)!.
\end{gather*}

\begin{lemma}[Hook-length bound]\label{pf:top-shape:hook}
For every non-increasing $\nu$, $\langle h_\nu,h_\nu\rangle\,P(\nu)=Q(\nu)$ and $1\le c(\nu)\le(\nu_0+N)^{N(N-1)/2}$.
\formallabel{hook_length_bound}
\formallabel{PI.HookSteps.hookBound_of}
\formallabel{PI.HookSteps.hook_norm}
\formallabel{PI.HookSteps.hook_bound}
\formallabel{PI.HookSteps.determinant_shift_norm}
\formallabel{PI.HookLower.one_le_cNu}
\formallabel{PI.HookSteps.prod_factorial_le}
\end{lemma}

\begin{proof}
(i) \emph{Adjointness.} On monomials, $\langle x_{ab}x^m,x^{m'}\rangle=\langle x^m,\partial_{ab}x^{m'}\rangle$: both sides
vanish unless $m'=m+e_{ab}$, in which case both equal $(m_{ab}+1)\,m!$. Hence multiplication by $x_{ab}$ is adjoint to
$\partial_{ab}$, and $\langle\det(X)f,g\rangle=\langle f,\det(\partial)g\rangle$ with
$\det(\partial)=\sum_\sigma\mathrm{sgn}(\sigma)\prod_j\partial_{\sigma(j)j}$
(\lref{PI.HookSteps.determinant_adjugate}).

(ii) \emph{Capelli identity.} As operators on $\Q[x]$,
\[
\det(X)\circ\det(\partial)=\sum_\sigma\mathrm{sgn}(\sigma)\prod_{j=0}^{N-1}\bigl(E_{\sigma(j)j}+(N-1-j)\,\delta_{\sigma(j)j}\bigr),
\]
the product being ordered by $j$ (\lref{PI.HookSteps.Cap.capelli_determinant_identity}; cf.\
\cite{HoweUmeda1991}).

(iii) \emph{Highest weight vectors.} For non-increasing $\mu$, $E_{ab}h_\mu=0$ for $a<b$ (the operator replaces row $b$ by
row $a$ in each minor, producing two equal rows or not acting), and $E_{aa}h_\mu=\mu_ah_\mu$ (the degree of $h_\mu$ in the
row $a$ is $\sum_{s\ge a}(\mu_s-\mu_{s+1})=\mu_a$). For a vector $v$ with $E_{ab}v=0$ ($a<b$) and $E_{aa}v=\lambda_av$, the
Capelli determinant with shift $u$ acts on $v$ by $\prod_j(u+\lambda_j-j)$: in its expansion only the identity
permutation contributes on such a vector (\lref{Formalism.CohomologicalDescent.HeightDescentMinors.c2_holds_alt}).
Moreover $h_{\mu+1}=\det(X)\,h_\mu$, since $\Delta_N=\det X$ and only the exponent of $\Delta_N$ changes. Applying (ii) and
(iii) to $h_{\mu+1}$ with $u=N-1$, $\lambda_j=\mu_j+1$, we get $\det(X)\det(\partial)(\det(X)h_\mu)=\prod_j(\mu_j+N-j)\det(X)h_\mu$,
and since $\det X\ne0$ in the domain $\Q[x]$, $\det(\partial)(\det(X)h_\mu)=\prod_j(\mu_j+N-j)\,h_\mu$. By (i),
\begin{equation}\label{eq:pf:top-shape:shift}
\langle h_{\mu+1},h_{\mu+1}\rangle=\langle h_\mu,\det(\partial)(\det(X)h_\mu)\rangle=\prod_j(\mu_j+N-j)\,\langle h_\mu,h_\mu\rangle .
\end{equation}

(iv) \emph{The norm formula}, by induction on $N$ (the case $N=0$ being trivial) and, for $N=n+1$, on the last entry
$\nu_n$. If $\nu_n=0$, $h_\nu$ is the image of $h_{\nu'}$, $\nu'=(\nu_0,\dots,\nu_{n-1})$, under the injective renaming of
the variables $x_{ab}\mapsto x_{ab}$ ($a,b<n$), which preserves Fock norms; and $P(\nu)=P(\nu')\prod_{i<n}(\nu_i+n-i)$,
$Q(\nu)=Q(\nu')\prod_{i<n}(\nu_i+n-i)$ (the last factor of $Q(\nu)$ being $0!=1$). If $\nu_n\ge1$, then $\nu=\mu+1$ with $\mu$
non-increasing and $\mu_n=\nu_n-1$; $P(\nu)=P(\mu)$, $Q(\nu)=Q(\mu)\prod_j(\mu_j+N-j)$, and \eqref{eq:pf:top-shape:shift}
gives the claim for $\nu$ from that for $\mu$.

(v) \emph{The bounds.} The coefficient of the diagonal monomial $\prod_ix_{ii}^{\nu_i}$ in $h_\nu$ is $1$, its weight $m!$
is $\prod_i\nu_i!$, and all terms of $\langle h_\nu,h_\nu\rangle$ are non-negative; hence $\langle h_\nu,h_\nu\rangle\ge\prod\nu_i!$,
i.e.\ $c(\nu)\ge1$ (\lref{PI.HookLower.fockQ_hw_ge}). Each factor of $P(\nu)$ is a positive integer, so
$\langle h_\nu,h_\nu\rangle\le Q(\nu)$. Finally $(\nu_i+N-1-i)!\le\nu_i!\,(\nu_0+N)^{N-1-i}$, since each of the factors
$\nu_i+1,\dots,\nu_i+N-1-i$ is at most $\nu_0+N$, and $\sum_i(N-1-i)=N(N-1)/2$.
\end{proof}

\subsubsection*{The Arakelov slope inequalities}

Let $D$ be a datum with host field $E$ (a CM field), $v\ge2$ the half root discriminant of the host
(\lref{Formalism.MainNotions.flipTwoThree.two_le_half_rootDiscr}), $k$ the column count and $f$ the rank $f$ of
$D$. For a word $u$ and an embedding type $\tau$ let $\lambda_{\tau,0}\ge\dots\ge\lambda_{\tau,k-1}$ be the entries of the
first column of the weight of $u$ at level $0$, and put
\begin{gather*}
a_\tau=\sum_{l<k}\lambda_{\tau,l},\qquad a'_\tau=\sum_{l<k}(\lambda_{\tau,l}-\lambda_{\tau,k-1}),\qquad
s_\tau=\lambda_{\tau,0}-\lambda_{\tau,k-1},\\
\lambda_\tau=\lambda_{\tau,k-1},\qquad\nu_{u,\tau}=(\lambda_{\tau,l}-\lambda_{\tau,k-1})_{l<k}.
\end{gather*}
Then $\nu_{u,\tau}$ is a non-increasing sequence of natural numbers with $\sum_l\nu_{u,\tau,l}=a'_\tau\ge0$ and
$\nu_{u,\tau,0}=s_\tau\ge0$ (\lref{Formalism.MainNotions.flipTwoThree.sum_natWeight_eq_colAbsShifted}). For a
non-increasing $\nu\in\N^N$ let $c(\nu)$ be as in \cref{pf:top-shape:hook}, and let $\mu_{\rm std}(\nu)$ be the normalised
slope of the integral Weyl module $W_\Z(\nu)\subset\Z[x]$ of weight $\nu$ of the formal development: for a $\Z$-basis
$e_1,\dots,e_r$ ($r>0$), $\mu_{\rm std}(\nu)=-\frac1r\log\det\bigl(\langle e_i,e_j\rangle/\langle h_\nu,h_\nu\rangle\bigr)$.

\begin{lemma}\label{pf:top-shape:loghook}
For every word $u$ and type $\tau$, $\log c(\nu_{u,\tau})\le k^2\log(s_\tau+k)$.
\formallabel{Formalism.MainNotions.flipTwoThree.log_hookConst_le_sq_mul_log}
\formallabel{Formalism.MainNotions.flipTwoThree.hookBound_holds}
\end{lemma}

\begin{proof}
By \cref{pf:top-shape:hook} with $N=k$, $1\le c(\nu_{u,\tau})\le(s_\tau+k)^{k(k-1)/2}$; as $s_\tau+k\ge1$ and
$\lfloor k(k-1)/2\rfloor\le k(k-1)\le k^2$, taking logarithms gives the claim.
\end{proof}

\begin{theorem}[Lower slope inequality]\label{pf:top-shape:lowslope}
For every datum $D$ and every word $u$ with column valuations,
\[
\mu(M_u)\ge-\sum_\tau\frac{a_\tau}k\log T_u-\sum_\tau\bigl(a'_\tau\log k+k^2\log(s_\tau+k)\bigr).
\]
In particular property \ref{G:low} holds for every family (with threshold $0$).
\formallabel{Formalism.L2.GlobalTheoryTheory.TraceIntegrality.lower_slope_inequality_family}
\formallabel{Formalism.MainNotions.flipTwoThree.familySlopeLowerOpen_of_std}
\formallabel{Formalism.MainNotions.flipTwoThree.slopeLowerOpen_of_std_hook}
\formallabel{Formalism.L2.GlobalTheoryTheory.TraceIntegrality.column_exact_slope_family}
\end{theorem}

\begin{proof}
\emph{Exact slope.} For a word with column valuations,
\[
\mu(M_u)=-\sum_\tau\frac{a_\tau}k\log T_u+\sum_\tau\mu_\tau,\qquad\mu_\tau=\mu_{\rm std}(\nu_{u,\tau})
\]
(\lref{Formalism.L2.GlobalTheoryTheory.TraceIntegrality.column_lattice_exact_slope}).

\emph{Standard slopes.} For every non-increasing $\nu\in\N^N$,
$\mu_{\rm std}(\nu)\ge-\bigl(\sum_l\nu_l\log N+\log c(\nu)\bigr)$ (\lref{standard_slope_lower}). With $N=k$,
$\sum_l\nu_{u,\tau,l}=a'_\tau$ and \cref{pf:top-shape:loghook},
$\mu_\tau\ge-\bigl(a'_\tau\log k+k^2\log(s_\tau+k)\bigr)$. Summing over $\tau$ and inserting into the exact slope gives the
inequality. The family statement holds at every member, hence with threshold $0$.
\end{proof}

For the upper inequality we use short vectors. Let $c$ be a class of the level-zero genus of $u$ and $\Lambda_c$ its
class lattice in the hermitian space $X_0$ (with form $h$). A \emph{short family} for $c$ is a family $y_0,\dots,y_{k-1}$
of $E$-linearly independent vectors of the dual lattice of $\Lambda_c$ with
$\mathrm{tr}_{E/\Q}h(y_a,y_a)/(2f)\le k^2(2v)^2$ for all $a$. Let $N_c\subseteq M_u$ be the part of $M_u$ belonging to the
class $c$; the column valuation decomposes $M_u$ as the orthogonal sum of the $N_c$.

\begin{theorem}[Maximal slope given short vectors]\label{pf:top-shape:slopemaxgiven}
Let $D$ be a datum and $u$ a word with column valuations such that every class $c$ of level zero has a short family.
Then
\[
\mu_{\max}(M_u)\le\sum_\tau\Bigl(a'_\tau\log\bigl(4k^4v^2\bigr)+\log c(\nu_{u,\tau})-\lambda_\tau\log T_u\Bigr).
\]
\formallabel{Formalism.L2.GlobalTheoryTheory.TraceIntegrality.slopeMaxGiven_free}
\formallabel{Formalism.L2.GlobalTheoryTheory.TraceIntegrality.slopeMaxGiven_of_class}
\formallabel{Formalism.L2.GlobalTheoryTheory.TraceIntegrality.slopeMaxClass_of_map}
\formallabel{Formalism.L2.GlobalTheoryTheory.TraceIntegrality.class_map_of_archimedean}
\end{theorem}

\begin{proof}
Let $\mathit{LL}$ be the coefficient field of the lattices, $n_\infty$ the number of its infinite places and
$[\mathit{LL}:\Q]$ its degree, and let $\mathrm{Mon}$ be the set of monomial indices of the weights
$\lambda_u=(\lambda_{u,\tau})_\tau$ at level $0$.

\emph{Step 1 (reduction to one class).} The pieces $N_c$ are non-zero, pairwise orthogonal at every infinite place, of
equal ranks, and span $M_u$; hence there is a class $c$ with $\mu_{\max}(M_u)\le\mu_{\max}(N_c)$
(\lref{Formalism.L2.GlobalTheoryTheory.TraceIntegrality.muMax_le_class}, from the slope inequality for orthogonal
sums, \lref{Formalism.L2.Col.hlat_orth_sum}). Fix this $c$ and a short family $y$ for it.

\emph{Step 2 (the archimedean class map).} By \lref{Formalism.L2.GlobalTheoryTheory.TraceIntegrality.class_map_archimedean}
there are an $\mathit{LL}$-linear map $\varphi$ from the space of $M_u$ to $\mathit{LL}^{\mathrm{Mon}}$, injective on $N_c$, a
non-zero fractional ideal $I$ with $\varphi(N_c)\subseteq I^{\mathrm{Mon}}$, and positive numbers $D_v$ (for the infinite
places $v$) such that, for all $v$ and $x$,
\[
\|\varphi(x)\|_{v,\mathrm{mon}}^2\le\Bigl(\prod_\tau G_v(\tau)^{a'_\tau}\Bigr)D_v\,\|x\|_v^2,\qquad
G_v(\tau)=\sum_{a<k}\mathrm{Re}\,v\bigl(\rho_\tau(h(y_a,y_a))\bigr)>0,
\]
where $\|\cdot\|_{v,\mathrm{mon}}$ is the monomial norm (the coordinate norm weighted by the inverse monomial
normalisations) and $\rho_\tau$ the embeddings attached to the types, and
\[
-\frac{2\log\mathrm N(I)}{[\mathit{LL}:\Q]}+\frac1{n_\infty}\sum_v\log D_v\le-\sum_\tau\lambda_\tau\log T_u .
\]
($G_v(\tau)>0$ because $h$ is positive definite at every place and the $y_a$ are non-zero.)

\emph{Step 3 (shortness).} For every $\tau$, the mean of $\log G_v(\tau)$ over the infinite places $v$ of $\mathit{LL}$
equals the mean of $\log\sum_a\mathrm{Re}\,w(h(y_a,y_a))$ over the infinite places $w$ of $E$
(\lref{Formalism.L2.Col.mean_over_places_comap}); by
the shortness of $y$ this mean is at most $\log(k^3(2v)^2)$
(\lref{Formalism.L2.GlobalTheoryTheory.TraceIntegrality.shortAt_mean_log}). Put
$B_v=\bigl(\prod_\tau G_v(\tau)^{a'_\tau}\bigr)D_v$. Since $a'_\tau\ge0$,
\begin{align*}
\frac1{n_\infty}\sum_v\log B_v&=\sum_\tau a'_\tau\cdot\frac1{n_\infty}\sum_v\log G_v(\tau)+\frac1{n_\infty}\sum_v\log D_v\\
&\le\sum_\tau a'_\tau\log\bigl(k^3(2v)^2\bigr)+\frac1{n_\infty}\sum_v\log D_v,
\end{align*}
so $-\frac{2\log\mathrm N(I)}{[\mathit{LL}:\Q]}+\frac1{n_\infty}\sum_v\log B_v\le\sum_\tau\bigl(a'_\tau\log(k^3(2v)^2)-
\lambda_\tau\log T_u\bigr)$.

\emph{Step 4 (slopes of lattices mapping to monomials).} Let $c_{\rm mon}(\lambda_u)>0$ be the monomial constant, a
lower bound for all monomial weights (\lref{Formalism.L2.GlobalTheoryTheory.TraceIntegrality.monW_ge_cNu}),
with $-\log c_{\rm mon}(\lambda_u)=\sum_\tau\bigl(a'_\tau\log k+\log c(\nu_{u,\tau})\bigr)$
(\lref{Formalism.L2.GlobalTheoryTheory.TraceIntegrality.log_monCnu}). For a finitely generated
$\OO_{\mathit{LL}}$-submodule $N$ mapped injectively by $\varphi$ into $I^{\mathrm{Mon}}$ with the norm comparison of Step~2
(with $B_v$), \lref{Formalism.L2.GlobalTheoryTheory.TraceIntegrality.muMax_le_mon} gives
\[
\mu_{\max}(N)\le-\log c_{\rm mon}(\lambda_u)-\frac{2\log\mathrm N(I)}{[\mathit{LL}:\Q]}+\frac1{n_\infty}\sum_v\log B_v :
\]
(the image of $N$ lies in $I$ times the standard lattice of $\mathit{LL}^{\mathrm{Mon}}$, an orthogonal sum of rank-one
lattices whose monomial weights are at least $c_{\rm mon}(\lambda_u)$).

\emph{Step 5.} Applying Step~4 to $N=N_c$ and Step~3,
\[
\mu_{\max}(N_c)\le\sum_\tau\bigl(a'_\tau\log k+\log c(\nu_{u,\tau})\bigr)+\sum_\tau\bigl(a'_\tau\log(k^3(2v)^2)-\lambda_\tau
\log T_u\bigr),
\]
and $\log k+\log(k^3(2v)^2)=\log(4k^4v^2)$. With Step~1 this is the claim.
\end{proof}

\begin{theorem}[Upper slope inequality]\label{pf:top-shape:upslope}
For every datum $D$ and every word $u$ with column valuations,
\[
\mu_{\max}(M_u)\le\sum_\tau\bigl(4a'_\tau\log(kv)+k^2\log(s_\tau+k)+|\lambda_\tau|\log T_u\bigr).
\]
In particular property \ref{G:max} with $C_5=4$ holds for every family (with threshold $0$).
\formallabel{Formalism.L2.GlobalTheoryTheory.TraceIntegrality.slope_inequality_family_four}
\formallabel{Formalism.L2.GlobalTheoryTheory.TraceIntegrality.slope_inequality_constant_four}
\formallabel{Formalism.MainNotions.flipTwoThree.slope_bound_of_max_slope}
\formallabel{Formalism.MainNotions.flipTwoThree.slope_max_open_four_of_hookBound}
\formallabel{Formalism.MainNotions.flipTwoThree.slopeMaxOpen_of_hook}
\formallabel{Formalism.MainNotions.flipTwoThree.slopeMaxOpen_of_hook_two_le_v}
\formallabel{Formalism.MainNotions.flipTwoThree.slopeMaxC_of}
\end{theorem}

\begin{proof}
The column valuations give, for every class $c$, the lattice condition at $T_u^{2f}$ for the class lattice, and under
this condition a short family exists, by Minkowski-type transference for hermitian lattices
(\lref{Formalism.MainNotions.flipTwoThree.shortDual_holds}). By \cref{pf:top-shape:slopemaxgiven},
$\mu_{\max}(M_u)\le\sum_\tau\bigl(a'_\tau\log(4k^4v^2)+\log c(\nu_{u,\tau})-\lambda_\tau\log T_u\bigr)$. We bound the three
terms. (i) As $v\ge2$, $4\le v^2$ and $4k^4v^2\le(kv)^4$, so $\log(4k^4v^2)\le4\log(kv)$; multiply by $a'_\tau\ge0$.
(ii) $\log c(\nu_{u,\tau})\le k^2\log(s_\tau+k)$ by \cref{pf:top-shape:loghook}. (iii) $T_u\ge1$, so $\log T_u\ge0$ and
$-\lambda_\tau\log T_u\le|\lambda_\tau|\log T_u$. Summing gives the claim; it holds at every member, hence with threshold
$0$.
\end{proof}

\begin{lemma}[Column valuations]\label{pf:top-shape:tcv}
If eventually the members of a family carry symmetric column ties and satisfy the column admissibility
$\mathrm{adm}_3$, then eventually every word of every member has column valuations.
\formallabel{Formalism.MainNotions.tCV_of_ties}
\formallabel{Formalism.MainNotions.RT.tCV_of_tied}
\end{lemma}

\begin{proof}
For a datum with symmetric column ties satisfying $\mathrm{adm}_3$, every word has column valuations: this is the
construction of the column valuations from the twist compatibility of the columns
(\lref{Formalism.L2.GlobalTheoryTheory.ColValS.colValOfTiedS_holds}). Apply it at every member beyond the
threshold of the hypothesis.
\end{proof}

\begin{proposition}[\cref{prop:families:famgl}]\label{pf:top-shape:linkgl}
Let $\mathcal A$ be a family class all of whose families satisfy $\sfname{Duty}$, and assume that $\mathrm{tie}$ is closed
on $\topshapeD(\mathcal A)$, and $h_U$, $h_G$. Then the eleven properties \ref{G:mass}--(G11) are closed on
$\topshapeD(\mathcal A)$ (with $C_5=4$), and the properties
\begin{gather*}
\sfname{ChosenE}(3,C_H)\Rightarrow\sfname{NuFullF}\wedge\sfname{MuRank},\\
\sfname{TuSize}\Rightarrow\sfname{DGamSize}\Rightarrow\sfname{LinkGLOBAL}(2,C_H,3,2,0)
\end{gather*}
are closed on $\topshapeD(\mathcal A)$.
\formallabel{GlobalTheory.Top.link_gl}
\formallabel{Formalism.MainNotions.Wire.FamGLmin.muRank}
\formallabel{Formalism.MainNotions.Wire.FamGLmin.link8}
\formallabel{Formalism.MainNotions.Wire.FamGLmin.link7}
\formallabel{Formalism.MainNotions.Wire.linkGLOBAL_of_7w}
\end{proposition}

\begin{proof}
The first assertion is \cref{pf:top-shape:famgl}. (The development also provides the totally positive
representation property used in the slope estimates, \lref{GlobalTheory.hermRepUnram}: in a CM field unramified
over its maximal real subfield, every totally positive $n$ is of the form $as\bar s+bt\bar t$ for given totally positive
$a,b$.)

\emph{Ranks.} Let $\topshapeF$ satisfy (G7)--(G11), $\sfname{TuSize}$ and $\sfname{ChosenE}(3,C_H)$. $\sfname{TuSize}$
gives the eventual level-prime setting (\lref{Formalism.L2.GlobalTheoryTheory.Places.settingF_of}); (G7)--(G11)
give the family properties ``two'' and ``five''
(\lref{Formalism.MainNotions.Wire.family_theorems_two_and_five}); these give the level statement
$\sfname{LvlF}$ (\lref{Formalism.MainNotions.lvlF_of}), hence the commensurability of the $N_u$ and
$\sfname{NuFullF}$ (\lref{Formalism.MainNotions.nuCommensF_min} and \lref{Formalism.MainNotions.Lnk.nuFullF_of}), and,
with $\sfname{ChosenProviso}(3)$, $\sfname{MuRank}$ (\lref{Formalism.MainNotions.muRank_min}). Since
$\sfname{TuSize}$ holds on $\topshapeD(\mathcal A)$, the first property is closed.

\emph{Global link.} Let $\topshapeF$ satisfy the eleven properties, $C_H>0$ (closed by
\cref{pf:top-shape:calculus}(f)), $\sfname{Duty}$ and $\sfname{TuSize}$. By the global theorem
(\cref{pf:top-shape:globalthm}) with $C_{\rm abs}=2$, $C_0=3$, $\topshapeF$ satisfies $\sfname{LinkGLOBAL7}(2,C_H,3,2,0)$,
i.e.\ the conclusion of $\sfname{LinkGLOBAL}$ under the seven-share condition $\sfname{WeightVectorSmallSharesSeven}(3,C_H)$
in place of $\sfname{ChosenE}(3,C_H)$. The wall condition $\sfname{ChosenWall}$ is part of $\sfname{Duty}$, and
$\sfname{ChosenE}(C_0,C_H)$ together with $\sfname{ChosenWall}$ implies $\sfname{WeightVectorSmallSharesSeven}(C_0,C_H)$
(\lref{Formalism.L2.GlobalTheoryTheory.Loss7.loss_requirements_of_small_shares}: the thick-set lower bound gives
the thick density, the bad-top bound gives its refined form, and the wall gives the chosen weights). Hence
$\sfname{LinkGLOBAL}(2,C_H,3,2,0)$ holds, and the hypothesis $\sfname{DGamSize}$ is not used.
\end{proof}

% ===========================================================================
\subsection{The global theorem}\label{sec:pf:top-shape:global}

We now prove \cref{thm:global-chain}. Its formal statement is $\sfname{LinkGLOBAL7}$ for a single family, under the
eleven properties of \cref{pf:top-shape:famglmin} as eventual statements and $\sfname{TuSize}$. The seven-share
condition $\sfname{WeightVectorSmallSharesSeven}(C_0,C_H)$ asserts that there are $C_5,c_6>0$ with the thick-set
density condition, $\sfname{ChosenTT}$, a refined bad-top condition, the chosen weights $\sfname{ChosenWS}$,
$\sfname{ChosenQD}(C_5)$, $\sfname{ChosenThin}(c_6)$, $\sfname{ChosenProviso}(C_0)$ and $\sfname{ChosenCode}(C_0,C_H)$.

\begin{theorem}[\cref{thm:global-chain}]\label{pf:top-shape:globalthm}
Let $\topshapeF$ be a family with parameter $\eps$, and $C_{\rm abs}\ge0$, $C_H>0$, $C_0\ge3$ and $C_5\in\R$. Assume
that the statements \ref{G:mass}--(G11) of \cref{pf:top-shape:famglmin} hold eventually for $\topshapeF$ (with this
$C_5$ in \ref{G:max}). If $\sfname{TuSize}(\topshapeF)$, then $\sfname{LinkGLOBAL7}(C_{\rm abs},C_H,C_0,C_{\rm abs},0)$:
if $\sfname{WeightVectorSmallSharesSeven}(C_0,C_H)$ and $\sfname{GlobalTheoremInputs}(C_{\rm abs},C_H,C_0)$ hold, then
there are $K,c_{HH},e_K,C_L,x_0$ such that $\topshapeF$ is chain-tied with exponents $(C_{\rm abs},0)$ and these
constants at every window scale $x\ge x_0$.
\formallabel{Formalism.MainNotions.linkGLOBAL7_min}
\formallabel{Formalism.MainNotions.linkGLOBAL7_items}
\formallabel{Formalism.MainNotions.LinkGLOBAL7}
\end{theorem}

The proof has two steps: the inputs are converted into eventual statements about the members of the family, one
for each system of size conditions, and the record is then assembled from the members (\cref{pf:top-shape:core}).

\begin{definition}[Family statements]\label{pf:top-shape:famstmts}
For a family $\topshapeF$ the following statements mean that there are constants (and $C\ge0$ where a constant
appears) such that eventually every member $(D,L)$ satisfies:
\begin{itemize}
\item $\sfname{CountF}$: for every word $u$, $\log T_u\ge\sum_{j}\frac1f\sum_\tau\log(w_j(\tau)v/\pi)+m\,\mathrm{lg}
-C(L_a+m^2L_a/k)$, $\log T_u\le Cm\Lambda$, and $|\log T_u-\log T_{u'}|\le C\log k$ (this is \ref{S:count});
\item $\sfname{SlopesMF}$: $-C|TT|aL_a(1+m/k)\le\mu(M_u)$ and $\mu_{\max}(M_u)\le C|TT|aL_a$ (this is \ref{S:slopes});
\item $\sfname{SlopesNF}$: $\mu_{\max}(N_u)\le\mu_{\max}(M_u)+\mathrm{Ga}+F_o-E(u)+CmB_0$ and
$\mu(M_u)+\mathrm{Ga}+F_o-E(u)-\mathrm{Lo}(u)\le\mu(N_u)$ (this is \ref{S:quot});
\item $\sfname{LatF}$: $\rank N_u>0$ and $\mu(N_u)\le\mu_{\max}(N_u)$ for all $u$, and
$\mu^{\min}_{\rm dual}(\mathit{RR})=-\mu_{\max}(\mathit{RR})$;
\item $\sfname{NodesF}$: the weights $w_j(\rho)$ read off from the member are $\ge1$, lie in $[b,b+\eta_{\rm sz}+1]$ at
thick nodes and are at most $2b'$ elsewhere;
\item $\sfname{FrameF}(C_H)$: for every large $x$ there are $\bar\eta,\bar e\ge0$ with $\bar e\le\frac14$ and
$(\bar\eta+4\sqrt{\bar e})x^{C_H}\le6$ such that every member at $x$ satisfies the frame inequality: there are
$\eta\in[0,\bar\eta]$ and $e\in[0,\bar e]$ such that for all $\mathrm{HH}>0$ with $s_{\rm Frame}-\mu^{\min}_{\rm dual}\le\mathrm{HH}$,
\[
s_{\rm Frame}-\log|W|-\frac{4u_F\log(1+\mathrm{HH}/u_F)+(\eta+4\sqrt e)\mathrm{HH}}{1-\delta_p}\le\mu_{\rm dual}(\mathit{RR}),
\]
where $s_{\rm Frame}$ and $u_F$ are formed from the lattices $N_u$ and heights $h(u)$ (this is \ref{S:slopeineq} with
$\delta=\delta_p$);
\item $\sfname{LossFS}(e_L)$: at every window scale $x$ there are a member $i_0$ and integer sequences $a_v,b_v$ with
$a_v(t)$ in the admissible window and $b'\le b_v(t)\le b_{\rm Top}$ for the tops $t$, some member with top weights
$a_v,b_v$, and such that every member with these top weights satisfies
$\mathrm{IO}\le C|TT|mb\,x^{e_L\eps}$ and $\sum_{t\in TT}\sum_{j<m}\mathrm{ASK}(\cdots)-\mathrm{Ga}\le C|TT|mb\,x^{e_L\eps}$
for all admissible readings of the weights (this is \ref{S:loss});
\item $\sfname{CharF}$ (for all $c_6,C_5$): there is an additive-character functional $\ell$ with
$\ell(\psi)=L'/L(1,\chi)$ for the primitive Dirichlet character $\chi$ attached to every non-trivial $\psi$, with
$L_{pL}=\Re\ell(\chi_C)$, and with thin-scale and value conditions under $\sfname{ChosenThin}(c_6)$ and
$\sfname{ChosenQD}(C_5)$;
\item $\sfname{WinF}$ (for all $c_6>0$, with some $e_K\ge0$): the size windows, e.g.\ $m\le x^{1/2+\eps/2}$, $xm\le k\le
x^{e_K}$, $k\le a$, $0\le L_a\le2x^{\eps}$, $X''\le a$, $a_T(t)\le2a$, $x^{1/2-\eps}\le m$ and $|a_T(t)/b-x|\le1$.
\end{itemize}
\formallabel{Formalism.MainNotions.Lnk.CountF}\formallabel{Formalism.MainNotions.Lnk.SlopesMF}
\formallabel{Formalism.MainNotions.Lnk.SlopesNF}\formallabel{Formalism.MainNotions.Lnk.LatF}
\formallabel{Formalism.MainNotions.Lnk.NodesF}\formallabel{Formalism.MainNotions.Lnk.FrameF}
\formallabel{Formalism.MainNotions.Lnk.LossFS}\formallabel{Formalism.MainNotions.Lnk.CharF}
\formallabel{Formalism.MainNotions.Lnk.WinF}
\end{definition}

\begin{proof}[Proof of \cref{pf:top-shape:globalthm}]
Assume $\sfname{TuSize}$, the seven-share condition (with constants $C_5',c_6$) and
$\sfname{GlobalTheoremInputs}$, i.e.\ $\sfname{LinkG}$, $\sfname{LinkS}(C_0)$, $\sfname{LinkIRR}$,
$\sfname{LinkDIV}(C_{\rm abs},C_H)$ and $\sfname{LinkMu}$. We derive the family statements; each step cites the lemma
that proves it.
\begin{enumerate}[label=(\arabic*)]
\item $\sfname{TuSize}$ gives the eventual level-prime setting
(\lref{Formalism.L2.GlobalTheoryTheory.Places.settingF_of}). (G7), (G8) give the division-side conditions
(\lref{Formalism.MainNotions.product_first_zeroth_parts.division_side_conditions_at_sizeData}) and (G9)--(G11)
the congruence-lattice length bound (\lref{Formalism.L2.GlobalTheoryTheory.LocalLatticeBases.congruence_no_kernel_minimal});
together they give $\sfname{LvlF}$ (\lref{Formalism.MainNotions.lvlF_of}), hence the commensurability of the
$N_u$ (\lref{Formalism.MainNotions.nuCommensF_min}), $\sfname{NuFullF}$
(\lref{Formalism.MainNotions.Lnk.nuFullF_of}) and, with $\sfname{ChosenProviso}(C_0)$ and $C_0>0$,
$\sfname{MuRank}$ (\lref{Formalism.MainNotions.muRank_min}).
\item $\sfname{LinkS}(C_0)$, $\sfname{ChosenProviso}(C_0)$, $\sfname{LinkIRR}$ and $\sfname{LinkDIV}$ give the rank of
$\mathit{RR}$ (\lref{Formalism.MainNotions.Lnk.rrRankF_of}); with $\sfname{NuFullF}$ and $\sfname{MuRank}$,
$\sfname{LinkDIV}$ gives $\sfname{LatF}$ (\lref{Formalism.MainNotions.Lnk.latF}).
\item \ref{G:mass}--\ref{G:sigma} give the explicit Vandermonde count
(\lref{Formalism.MainNotions.StandardLattice.division_level_prime_family}), hence the part-five family bounds
(\lref{Formalism.MainNotions.GlobalTheoryTheory.family_part_five_of_explicit_constants}) and $\sfname{CountF}$
(\lref{Formalism.MainNotions.GlobalTheoryTheory.countF_fam}).
\item $\sfname{ChosenQD}(C_5')$ gives eventually $|L_{pL}|\le x^{\eps}$ (\lref{Formalism.MainNotions.Lnk.dl_of_chosen}).
With $\sfname{CountF}$, \ref{G:low}, \ref{G:max} and \ref{G:colval}, this gives $\sfname{SlopesMF}$
(\lref{Formalism.MainNotions.flipTwoThree.slopesMF_of_count_fam}).
\item $\sfname{TuSize}$ and $\sfname{LvlF}$ give the level statement for $T$
(\lref{Formalism.MainNotions.J.lemma_five_one_of_theorem_two}); with the commensurability, $\sfname{TuSize}$,
$\sfname{MuRank}$, $\sfname{LinkG}$ and $\sfname{LinkMu}$ this gives the lower bound for the places function
(\lref{Formalism.L2.GlobalTheoryTheory.Places.placesF_min'}), proved with a Mertens-type estimate for
$\tau$-sums, and hence $\sfname{SlopesNF}$ (\lref{Formalism.MainNotions.Lnk.slopesNF_of}).
\item $\sfname{LinkDIV}$, $\sfname{LinkIRR}$, $\sfname{LinkS}(C_0)$, $\sfname{ChosenProviso}(C_0)$, $\sfname{NuFullF}$,
$\sfname{MuRank}$, $\sfname{TuSize}$, the ramified ranks (from $\sfname{LvlF}$,
\lref{Formalism.MainNotions.ramRankF_min}) and a matrix bound (from $\sfname{LinkS}$, $\sfname{ChosenProviso}$
and $\sfname{LinkDIV}$, \lref{Formalism.MainNotions.Lnk.matKF_of}) give the real frame statement
(\lref{Formalism.MainNotions.Lnk.frameRealF_of''}); with $\sfname{LatF}$, $\sfname{ChosenProviso}$,
$\sfname{ChosenCode}(C_0,C_H)$, $\sfname{NuFullF}$, $C_H>0$ and $C_0\ge1$ this gives $\sfname{FrameF}(C_H)$
(\lref{Formalism.MainNotions.Lnk.frameF_of}).
\item $\sfname{ChosenTT}$, the refined bad-top condition and $\sfname{LinkG}$ give the losses of small shares
(\lref{Formalism.L2.GlobalTheoryTheory.Loss7.stLoss7}), and with the chosen weights $\sfname{ChosenWS}$ this gives
$\sfname{LossFS}(0)$ (\lref{Formalism.MainNotions.Lnk.lossF_of_7}).
\item $\sfname{CharF}$ and $\sfname{NodesF}$ hold for every family
(\lref{Formalism.MainNotions.Lnk.charF} and \lref{Formalism.MainNotions.Lnk.nodesF}), and the thick-set density gives
$\sfname{WinF}$ (\lref{Formalism.MainNotions.Lnk.size_windows_of_thick_density}).
\end{enumerate}
Now \cref{pf:top-shape:core} below, applied with $e_L=0$, $C_5'$, $c_6$, $\sfname{ChosenTT}$, $\sfname{ChosenQD}(C_5')$,
$\sfname{ChosenThin}(c_6)$, $\sfname{LinkIRR}$ and $\sfname{LinkDIV}(C_{\rm abs},C_H)$, gives the conclusion.
\end{proof}

\begin{definition}[Models and their records]\label{pf:top-shape:model}
Let $x$ be a window scale of a family $\topshapeF$. A \emph{model} at $x$ consists of a member $i_0\in I_x$, integer
sequences $a_v,b_v$ and a map $\mathrm{pt}$ from the frame group $A$ of $D_{x,i_0}$ to $I_x$ such that for every $u_0\in A$
the host base point of $D_{x,\mathrm{pt}(u_0)}$ is $u_0$ and its top weights are $a_v,b_v$. Given a model, an integer
$N$, real numbers $C_H,\bar\eta,\bar e$ and functions $\eta,e$ on $A$, the \emph{record of the model} is the record
$\mathbb D$ of \cref{def:record} with $U_0=A$, words $W(u_0)$ the words of $D_{x,\mathrm{pt}(u_0)}$,
$\mathrm{Top}$ the set of tops of the family at $x$, and entries:
\begin{itemize}
\item $x,\eps,a,b,b',k,L,L_a,B_0,N_c$ the scale quantities of $(D_{x,i_0},L_{x,i_0})$, $\gamma$ Euler's constant,
$L_{pL}$ the log-derivative entry of $L_{x,i_0}$, $v$ the host constant, $\mathrm{lg}=\mathrm{lg}(D_{x,i_0})$ and
$\log\Gamma$ the logarithm of the $\Gamma$-entry;
\item $a_T(t)=a_v(t)$, and $w(u_0,j,\rho)$ the weights read off from the nodes;
\item for a word $u$ of $D'=D_{x,\mathrm{pt}(u_0)}$: $d(u_0,u)=\rank N_u$, $T(u_0,u)=T_u$, $\mu_N,\mu_N^{\max}$ the mean and
maximal slopes of $N_u$, $h(u_0,u)=h(u)$, $\mu_M,\mu_M^{\max}$ the mean and maximal slopes of $M_u$, and
$E_j(u_0,u)=\sum_w\mathrm{sz}(w,E(u,j+1,w))$;
\item $\mu_{CC}(u_0)=\mu_{\rm dual}(\mathit{RR})$, $\mu^{\min}_{CC}(u_0)=\mu^{\min}_{\rm dual}(\mathit{RR})$, $F_o$, $\mathrm{Ga}$
and $\mathrm{IO}$ the entries of $L_{x,\mathrm{pt}(u_0)}$, $\mathrm{Ta}=\sum_w\mathrm{sz}(w,\tau(w))$,
$\mathrm{ASK}(u_0)=\sum_{t\in TT}\sum_{j<m}\mathrm{ASK}(\mathrm{arith}(H),N,a(t),t,w(u_0,j,\cdot))$,
$\mathrm{HH}=x^{C_H}b$, $\delta(u_0)=\delta_p(D')$, $\eta,e$ as given;
\item $\bar\delta=\exp(-x^{\eps_z/7})$, and $\bar\eta,\bar e$ as given.
\end{itemize}
\formallabel{Formalism.MainNotions.Lnk.Mod.data}
\end{definition}

Thus the quantities of a record produced by the global theorem are slopes and heights of the lattices of the members:
$\mu_N$ and $\mu_M$ are the slopes of $N_u$ and of the level-zero lattice $M_u$, $\mu_{CC}$ is the dual slope of the glued
lattice $\mathit{RR}$, and the parameter $\delta$ of the slope inequality is the non-generic proportion $\delta_p$
controlled by $\sfname{LinkIRR}$.

\begin{theorem}[Assembly of the record]\label{pf:top-shape:core}
Let $\topshapeF$ be a family with parameter $\eps$, and $C_{\rm abs},C_H,e_L\ge0$, $C_5\in\R$, $c_6>0$. Assume
$\sfname{CountF}$, $\sfname{SlopesMF}$, $\sfname{SlopesNF}$, $\sfname{LatF}$, $\sfname{NodesF}$, $\sfname{FrameF}(C_H)$,
$\sfname{LossFS}(e_L)$, $\sfname{CharF}$, $\sfname{WinF}$, $\sfname{ChosenTT}$, $\sfname{ChosenQD}(C_5)$,
$\sfname{ChosenThin}(c_6)$, $\sfname{LinkIRR}$ and $\sfname{LinkDIV}(C_{\rm abs},C_H)$. Then there are
$K,c_{HH},e_K,C_L,x_0$ such that $\topshapeF$ is chain-tied with exponents $(C_{\rm abs},e_L)$ and these constants at
every window scale $x\ge x_0$.
\formallabel{Formalism.MainNotions.Lnk.linkCore}
\end{theorem}

\begin{proof}
Let $C_d,C_c,C_s,C_1,C_l$ be the constants of $\sfname{LinkDIV}$, $\sfname{CountF}$, $\sfname{SlopesMF}$,
$\sfname{SlopesNF}$, $\sfname{LossFS}$, and let $e_K$ be given by $\sfname{WinF}$ for $c_6$. Put
\[
K=\max(C_d,0)+C_c+C_s+C_1+C_l+6,\qquad c_{HH}=C_H,\qquad C_L=C_5,
\]
and let $x_0$ exceed $x_1+1$ and all thresholds of the hypotheses. Let $x\ge x_0$ be a window scale. The constants of
the record are $C_{\rm Div}=\max(C_d,0)$, $e_{\rm Div}=C_{\rm abs}$, $e_{\rm Loss}=e_L$, $c_{HH}=C_H$, and $C_c,C_s,C_1,C_l$
for the count, slope, quotient and loss systems; so $C.e_{\rm Div}\le C_{\rm abs}$, $C.e_{\rm Loss}\le e_L$ and
$C.c_{HH}\le c_{HH}$.

\emph{The model.} $\sfname{LossFS}(e_L)$ at $x$ gives $i_0$, $a_v$, $b_v$ and a member with top weights $a_v,b_v$; the
covering property of the family then gives, for every $u_0\in A$, a member $\mathrm{pt}(u_0)$ with base point $u_0$ and
top weights $a_v,b_v$. This is a model $M$, and every member $\mathrm{pt}(u_0)$ satisfies the loss inequalities.
$\sfname{FrameF}(C_H)$ at $x$ gives $\bar\eta,\bar e$, and for each $u_0$ numbers $\eta(u_0)\in[0,\bar\eta]$,
$e(u_0)\in[0,\bar e]$ with the frame inequality at $\mathrm{pt}(u_0)$. $\sfname{CharF}$ gives the functional $\ell$. Let
$\mathbb D$ be the record of $M$ with $N=\lceil2a\rceil$ (\cref{pf:top-shape:model}).

\emph{The eleven systems.} We verify $\sfname{SizeData}_2(\mathbb D,C,X'')$, i.e.\ \ref{S:setting}--\ref{S:fparams}:
\begin{itemize}
\item \ref{S:setting}: positivity of $x,b,a,b',k,f$, of the top weights $a_T$ (as $0<a\le a_T(t)$) and the remaining
positivity conditions follow from
$\sfname{WinF}$, $\sfname{NodesF}$ ($w\ge1$), $\sfname{LatF}$ ($d>0$, $\mu_N\le\mu_N^{\max}$) and the primality of the
level primes; $\bar\delta=\exp(-x^{\eps_z/7})<1$, and $\bar e\le\frac14$ by $\sfname{FrameF}$.
\item \ref{S:div}, \ref{S:ecomp}, \ref{S:hcomp}: clauses (7), (8), (9) of $\sfname{LinkDIV}$ at the members
$\mathrm{pt}(u_0)$, with the constant $C_d$ replaced by $\max(C_d,0)$ (the error terms being non-negative, as $L_a\ge0$,
$B_0\ge0$, $N_c\ge0$); in \ref{S:div} the main term $\mathrm{MAIN}(u)$ of $L$ is the right-hand side written there, and
$\mu^{\rm dual}$ is $\mu_{CC}$.
\item \ref{S:count}: $\sfname{CountF}$; \ref{S:slopes}: $\sfname{SlopesMF}$; \ref{S:quot}: $\sfname{SlopesNF}$.
\item \ref{S:loss}: the loss inequalities of $\sfname{LossFS}$ at $\mathrm{pt}(u_0)$, applied to the reading of the weights
of the model; the cut condition $a(t)/w(u_0,j,\rho)\le N$ holds because $w\ge1$ and $a_T(t)\le2a\le\lceil2a\rceil$.
\item \ref{S:slopeineq}: the frame inequality at $\mathrm{pt}(u_0)$, with $\delta=\delta_p$.
\item \ref{S:hh}: $\mathrm{HH}=x^{C_H}b$; clause (2) of $\sfname{LinkDIV}$ gives $s_{\rm Frame}+\mu_{\max}(\mathit{RR})\le
x^{C_H}b$, and $\mu^{\min}_{\rm dual}=-\mu_{\max}(\mathit{RR})$ by $\sfname{LatF}$, so $s-\mu^{\min}_{CC}\le\mathrm{HH}$.
\item \ref{S:fparams}: $0\le\delta_p\le\exp(-x^{\eps_z/7})=\bar\delta$ by $\sfname{LinkIRR}$ at $\mathrm{pt}(u_0)$, and
$0\le\eta\le\bar\eta$, $0\le e\le\bar e$ by construction.
\end{itemize}

\emph{The remaining conditions of \cref{def:chaintied}.} The four ties hold by construction of the record: the members
$\mathrm{pt}(u_0)$ share the frame, $m$, $f$, $q$, $d$ and the top weights (tie of points), their scale quantities coincide
with those of $i_0$ (global tie), the lattice entries are those of the members (local tie), and the arithmetic data
are those of the host with the functional $\ell$ (arithmetic tie); the identification holds because $\ell$ represents
the logarithmic derivatives at $1$ and $L_{pL}=\Re\ell(\chi_C)$. The error windows follow from $\sfname{WinF}$; in
particular $(\bar\eta+4\sqrt{\bar e})x^{C_H}b\le6b\le Kx^{\eps}|TT|a$, because $b\le|TT|\,xb=|TT|a$ (as $|TT|\ge1$ by
$\sfname{ChosenTT}$) and $6\le Kx^{\eps}$. The bounds $B_0=2|A|mb\,x^{\eps}$ and $N_c=|T_2|kb\,x^{\eps}/(|A|^2x^{-\eps})$
are the definitions of these quantities ($|A|=n$). The remaining hypotheses of the growth inequality follow from
$\sfname{WinF}$, $\sfname{NodesF}$, the cut condition above, and the conditions of $\sfname{CharF}$ under
$\sfname{ChosenThin}(c_6)$ and $\sfname{ChosenQD}(C_5)$, with $\eta_{\rm sz}$, $X'$ and $\theta=\exp(-c_6x^{\eps/8})$;
the cost windows and the size conditions of \cref{thm:cheb-windows} follow from $\sfname{WinF}$ and
$\sfname{ChosenTT}$.
\end{proof}

% ===========================================================================
\subsection{Designed families}\label{sec:pf:top-shape:design}

Let $\mathrm M(N)=\prod_{p\le N+1}p^{e_p(N)}$, $e_p(N)=\sum_{i=0}^{N}\lfloor N/(p^i(p-1))\rfloor$, be Minkowski's
number; every prime divisor of $\mathrm M(N)$ is at most $N+1$.

\begin{lemma}[Minkowski's bound over $\Q$]\label{pf:top-shape:mink}
For every $N$ and every finite subgroup $G$ of $\GL_N(\Q)$, the order of $G$ divides $\mathrm M(N)$. In particular
$\sfname{MinkDiv}$ holds.
\formallabel{minkowski_bound_rat}
\formallabel{PI.MkQ.minkowskiBoundQ}
\formallabel{PI.MkQ.phi_injective}
\end{lemma}

\begin{proof}
Let $\Lambda\subseteq\Q^N$ be the $\Z$-span of the vectors $ge_i$, $g\in G$, $1\le i\le N$. It is finitely generated
(as $G$ is finite), contains $\Z^N$ and is stable under $G$; hence it is a free $\Z$-module of rank $N$. Choosing a
basis of $\Lambda$, the action of $G$ on $\Lambda$ defines a homomorphism $\phi\colon G\to\GL_N(\Z)$. It is injective:
if $\phi(g)=\phi(g')$, then $g$ and $g'$ agree on $\Lambda\supseteq\{e_1,\dots,e_N\}$, so $g=g'$. Thus $G$ is
isomorphic to a finite subgroup of $\GL_N(\Z)$, whose order divides $\mathrm M(N)$ by Minkowski's theorem
\cite{Minkowski1887} (\lref{card_dvd_minkowskiNumber}).
\end{proof}

\begin{lemma}\label{pf:top-shape:upN}
Let $\mathcal C$ be a family class, $C_0\in\R$ and $c_T\colon\R^2\to\R$, and put
$\mathcal C^\sharp=\mathrm{WithTw}(\mathrm{WithSm}(\mathrm{WithChN}(\mathcal C)))$.
\begin{enumerate}[label=(\alph*)]
\item If $\sfname{DesignChoices}(\mathrm{WithTw}(\mathcal C),C_0,c_T)$, then
$\sfname{DesignChoices}(\mathcal C^\sharp,C_0,c_T)$.
\item $\sfname{DGamDvdF}$ implies $\sfname{DGamSmooth}$ and $\sfname{DGamSize}$ (with constant $12$).
\item Suppose $\sfname{DesignChoices}(\mathcal C,C_0,c_T)$, and that the following properties are closed on
$\mathcal C$: $c_T(\eps,C_{YL})\le C_{Tu}\Rightarrow\sfname{TuSize}$; $\sfname{ChosenProviso}(C_0)$; and
``$\sfname{BoxesF}\wedge\sfname{QDRowsF}\wedge\exists C_5>0:\sfname{ChosenQD}(C_5)$ implies
$\exists c_6>0:\sfname{ChosenThin}(c_6)$''. Then $\sfname{DesignExists}(\mathcal C,C_0)$.
\end{enumerate}
\formallabel{Formalism.L2.Top3.designChoices_upN}
\formallabel{Formalism.L2.Top3.withSm_of_withChN}
\formallabel{Formalism.MainNotions.J.dGamSmooth_of_dvd}
\formallabel{Formalism.MainNotions.dGamSize_of_dvd}
\formallabel{Formalism.MainNotions.designExists_of_choices}
\end{lemma}

\begin{proof}
(b) If $D_\Gamma(u)$ divides $\mathrm M(2fN(m))$, every prime divisor $p$ of $D_\Gamma(u)$ divides $\mathrm M(2fN(m))$, so
$p\le2fN(m)+1$. Moreover $0<D_\Gamma(u)\le\mathrm M(2fN(m))$ and $\log\mathrm M(2fN(m))\le12fN(m)\log(fN(m))$
(\lref{Formalism.MainNotions.log_minkM_fN}, using $fN(m)\ge2$).

(a) The families produced by $\sfname{DesignChoices}(\mathrm{WithTw}(\mathcal C),C_0,c_T)$ satisfy $\sfname{DGamDvdF}$,
$\sfname{TameDistinctF}$, $\sfname{BoxesF}$, $\sfname{QDRowsF}$ and $C_{YL}\le9$ (they are among the raw choices), hence
lie in $\mathrm{WithChN}(\mathcal C)$, and by (b) in $\mathrm{WithSm}(\mathrm{WithChN}(\mathcal C))$; the other components
are unchanged.

(c) Let $\eps_0$ be the constant of $\sfname{DesignChoices}$ and $\eps_1$ that of the conjunction of the three closed
properties (\cref{pf:top-shape:calculus}(b)). Let $0<\eps\le\min(\eps_0,\eps_1)$ and let $K$ be a staged threshold;
let $K'$ be the threshold of the closed conjunction at $\eps$. $\sfname{DesignChoices}$ applied to $\eps$ and $K\sqcup
K'$ gives $\topshapeF\in\mathcal C$, $C_H\ge0$ and $C_5,C_b>0$ with the raw choices. By \cref{pf:top-shape:large},
$\topshapeF$ is large for $K$ and for $K'$ at level $C_H$. Hence $\sfname{TuSize}$ holds (as $c_T(\eps,C_{YL})\le C_{Tu}$
is a raw choice), $\sfname{ChosenProviso}(C_0)$ holds, and there is $c_6>0$ with $\sfname{ChosenThin}(c_6)$ (by
$\sfname{BoxesF}$, $\sfname{QDRowsF}$ and $\sfname{ChosenQD}(C_5)$). The seven choice conditions now give
$\sfname{ChosenE}(C_0,C_H)$ with $C_5,c_6,C_b$, and (b) gives $\sfname{DGamSize}$. So $\topshapeF$ witnesses
$\sfname{DesignExists}(\mathcal C,C_0)$ for $\eps$ and $K$.
\end{proof}

The hypothesis $\sfname{LitThin}$ of the package is the conjunction of three statements about Dirichlet
$L$-functions, each formally quantified over all data (on which it does not depend): (i) for some $C_1,c>0$, the total
multiplicity of the zeros $\rho$ of $L(s,\chi)$, $\chi$ primitive of conductor at most $Q$, with $\Re\rho\ge1-1/400$ and
$|\Im\rho|\le Q$ is at most $C_1Q^{1/4}$ ($Q\ge3$); a zero with $\Re\rho\ge1-c/\log\max(N,N|\Im\rho|)$ of $L(s,\chi)$,
$\chi$ primitive of conductor $N\ge3$, is real and simple, $\chi$ is real, and there is at most one such zero; and
$\zeta$ has no zero with $1-c/\log(|\Im\rho|+3)<\Re\rho<1$; (ii) a truncated explicit formula for the
Chebyshev function $\psi_0(u,\chi)$, with error $C(u\log^2(NuT)/T+\log(Nu)+\sum_{|\Im\rho|<1}1/|\rho|)$; (iii) the
number of zeros of $L(s,\chi)$ with $|\Im\rho-t|\le1$ is at most $C\log(N(|t|+2))$. These are classical
\cite{Davenport2000,Montgomery1971}; in the formal development they are proved. The condition
$\sfname{ChosenThin}(c_6)$ says that eventually, for the Dirichlet characters $\chi$ attached to the host and all
$y\ge X'/3$, the character sums satisfy $|g_\chi(y)|\le y\exp(-c_6x^{\eps/8})$.

\begin{proposition}[Thin-scale prime sums]\label{pf:top-shape:thin}
Let $\topshapeF$ be a family whose members satisfy (with constants independent of the member) the three statements
(i), (ii), (iii) above. Then there are $c>0$ and $x_0>1$ such that for $x\ge x_0$ and every member $D$: if $D$ satisfies
the box hypothesis at $(x,\eps)$, then for every primitive Dirichlet character $\chi$ attached to the host of $D$ (that
is, $\chi(a)=\psi(\mathrm{art}(a))$ for a character $\psi$ of the group $G_D$ of $D$ and the Artin map $\mathrm{art}$, for
$a$ prime to the conductor of the host) and every $y\ge e^{x^{\eps/3}}/3$, $|g_\chi(y)|\le y\exp(-cx^{\eps/8})$. In particular $\sfname{LitThin}$
implies this conclusion for every family.
\formallabel{Formalism.MainNotions.ShortIntervalPrimes.proposition_six_norm_family_of_admitted}
\formallabel{Formalism.MainNotions.prop_two_six_congruence_family_of_literature}
\end{proposition}

\begin{proof}
\emph{Real characters.} For every datum $D$, every real primitive character $\chi$ attached to the host of $D$ has a
modulus $N$ which divides $8\,\mathrm{odd}(Q_S)$ or is divisible by $q$ or by $d$, where $q,d$ are the auxiliary primes
of the host and $\mathrm{odd}(Q_S)$ is the odd part of its discriminant parameter
(\lref{Formalism.MainNotions.ShortIntervalPrimes.realClassE_holds}).

\emph{The thin-scale bound.} Given (i)--(iii) and this description of the real characters,
\lref{Formalism.MainNotions.ShortIntervalPrimes.proposition_six_family_of_real_classes} proves the bound. In that
proof the box hypothesis excludes real zeros in $(1-1/400,1)$ for the real characters attached to the host whose
modulus is divisible by $q$ or by $d$ (\lref{Formalism.MainNotions.ShortIntervalPrimes.good_real_ne}), so that by
the description above the remaining real characters have modulus dividing $8\,\mathrm{odd}(Q_S)$; the character sums
at the thin scales are then bounded by means of the explicit formula (ii), the zero-density and zero-free-region
statements (i) and the local zero count (iii).

\emph{From $\sfname{LitThin}$.} $\sfname{LitThin}$ provides $C_1,c>0$ for (i), $C_7>0$ for (ii) and $C_{11}>0$ for
(iii), for all data; these constants serve for every member of every family.
\end{proof}

\begin{proposition}[\cref{prop:families:design}]\label{pf:top-shape:designthm}
Let $\mathcal A$ be a family class with a function $c_T$, a twist selection $\sigma$, and $\sfname{Duty}$ for all its
families, and assume $h'_{\rm Ex}\colon\sfname{MinkDiv}\Rightarrow\sfname{DesignChoices}(\mathrm{WithTw}(\mathcal
A),3,c_T)$, the closedness of $h_{Tu}$ and $h_{n4}$ on $\mathrm{WithTw}(\mathrm{WithSm}(\mathrm{WithChN}(\mathcal A)))$, and
$\sfname{LitThin}$. Then $\sfname{DesignExists}(\mathrm{WithTw}(\mathrm{WithSm}(\mathrm{WithChN}(\mathcal A))),3)$.
\formallabel{GlobalTheory.Top.designExists_of_tuSize_aux4n_litThin}
\formallabel{Formalism.MainNotions.chosenThin_closed}
\formallabel{Formalism.MainNotions.chosenThin_of_boxesN}
\end{proposition}

\begin{proof}
Let $\mathcal C=\mathrm{WithTw}(\mathrm{WithSm}(\mathrm{WithChN}(\mathcal A)))\subseteq\mathcal A$. By
\cref{pf:top-shape:mink} and $h'_{\rm Ex}$, the design choices $\sfname{DesignChoices}(\mathrm{WithTw}(\mathcal A),3,c_T)$
exist, and by \cref{pf:top-shape:upN}(a) so do $\sfname{DesignChoices}(\mathcal C,3,c_T)$. We verify the three closed
hypotheses of \cref{pf:top-shape:upN}(c).
\begin{itemize}
\item The first is $h_{Tu}$.
\item $\sfname{ChosenProviso}(3)$: for $\topshapeF\in\mathcal A$, the selection and $\sfname{Duty}$ place $\topshapeF$ in
the second-layer class (as in the proof of \cref{pf:top-shape:linkG}), whose rank ties are closed
(\lref{Formalism.MainNotions.Ext.Lw.second_chern_linked_closed_rank_ties}); restricted to $\mathcal C$ and combined
with $h_{n4}$, the Harish-Chandra reduction \lref{Formalism.MainNotions.StandardLattice.HarishChandraReduction.chosenProviso_min}
(for $C_0\ge e$: the class-number formula for the levels and the rank ties imply $\sfname{ChosenProviso}(C_0)$) gives the
claim, as $e\le3$.
\item The thin choice: for every family with $\sfname{BoxesF}$ and $\sfname{QDRowsF}$ there is $c_6>0$ with
$\sfname{ChosenThin}(c_6)$, so the property is closed on every class (\cref{pf:top-shape:calculus}(a)). Indeed,
$\sfname{LitThin}$ gives $c>0$ and a threshold beyond which a thin-scale form of the short-interval prime theorem holds
for the members (\cref{pf:top-shape:thin}),
and $\sfname{QDRowsF}$, $\sfname{BoxesF}$ give eventually the pair and box hypotheses of that statement
(\lref{Formalism.MainNotions.pair_and_box_hypothesis_of_congruence_data_and_boxes}); together they give
$\sfname{ChosenThin}(c)$.
\end{itemize}
\Cref{pf:top-shape:upN}(c) concludes.
\end{proof}

In the proof of \cref{pf:top-shape:main} this proposition is applied to the class $\mathcal A\wedge\mathcal
G.\mathrm{cls}(\lambda)$, with $h'_{\rm Ex}=h_{\rm Ex}(\cdot)(\lambda)(\mathcal G)(h_H)$; the selection, $\sfname{Duty}$,
$h_{Tu}$ and $h_{n4}$ restrict to this class.

% ===========================================================================
\subsection{Remarks on the formal statements}\label{sec:pf:top-shape:remarks}

\begin{enumerate}[label=(\arabic*)]
\item \emph{Conditional form.} \Cref{pf:top-shape:main} is an implication from $\sfname{Green}$ and the data (a)--(h);
all closedness hypotheses are vacuous for classes without large members, and the existence of large members comes from
$h_{\rm Ex}$ (together with $\sfname{MinkDiv}$, proved in \cref{pf:top-shape:mink}). The closedness of $\sfname{LinkN}$, the
designs $h_{\rm Ex}$, $\sfname{WuDecayOfGap}$, $\sfname{RowsFinDim}$, the row estimates, $h_U$, $h_G$ and
$\sfname{LitThin}$ are hypotheses here and are established in the proof of \cref{thm:hyps-nonempty}.
\item \emph{Unused hypotheses.} The parameters $V$ and $N_R$ of the shape are the true proposition, so $h_{VC}$ and $n_R$
play no role. In $\sfname{LinkDIVt}=(\sfname{LinkEta}\wedge\sfname{DIVinN}\Rightarrow\sfname{LinkDIV})$ the proof of
\cref{pf:top-shape:divmain} does not use $\sfname{LinkEta}$; in \cref{pf:top-shape:linkgl} the hypothesis
$\sfname{DGamSize}$ of the global link is not used; and \cref{pf:top-shape:eta} does not use $\sfname{Duty}$.
\item \emph{Largeness.} The staged thresholds have a component $K_t$ which does not enter the definition of largeness (the
corresponding conjunct is the true proposition).
\item \emph{Dependencies not visible by name.} The global theorem enters the proof of the shape theorem through
\cref{pf:top-shape:linkgl}, where the record of eleven properties is applied to the global theorem by projection
(\lref{Formalism.MainNotions.Wire.FamGLmin.link8}, \lref{Formalism.MainNotions.Wire.FamGLmin.link7},
\lref{Formalism.MainNotions.linkGLOBAL7_min}, \lref{Formalism.MainNotions.Lnk.linkCore}); the same holds for the
word facts of the division packages (\lref{Formalism.L2.DIV.DivisionPackageColumns.wordFacts}) and for the
assembly of the bridge bounds (\lref{Formalism.L2.Top4.RemainingDivisionBounds.toRows_alt}).
\item \emph{Comparison with the outline of Part~I.} In the records produced by the global theorem the parameter $\delta$
of the slope inequality \ref{S:slopeineq} is the non-generic proportion $\delta_p$, bounded by
$\bar\delta=\exp(-x^{\eps_z/7})$ through $\sfname{LinkIRR}$; the parameters $\eta,e$ come from the frame statement
$\sfname{FrameF}$, which is derived (among others) from $\sfname{LinkS}$. The systems \ref{S:div}, \ref{S:ecomp} and
\ref{S:hcomp} are clauses (7)--(9) of $\sfname{LinkDIV}$, and the height bound \ref{S:hh} comes from clause (2) of
$\sfname{LinkDIV}$ and $\sfname{LatF}$. The volume property (G2) is the equality of the base volume with the mass in rank
$r$, deduced from the rank-two genus mass $h_G$ and the volume ratio $h_U$.
\item \emph{Sign conditions.} In $\sfname{LinkS}$ only $e_0$ is required to be non-negative; $\eta_0,\eta_1$ are only
bounded above.
\item \emph{Constants.} The theorem gives $e_D=2$ and $e_L=0$; the constant $C_{\rm abs}=2$ of the division link is the
exponent $e_D$ of the chain-tied records.
\end{enumerate}
